% This LaTeX document needs to be compiled with XeLaTeX.
\documentclass[10pt]{article}
\usepackage[utf8]{inputenc}
\usepackage{amsmath}
\usepackage{amsfonts}
\usepackage{amssymb}
\usepackage[version=4]{mhchem}
\usepackage{stmaryrd}
\usepackage{hyperref}
\hypersetup{colorlinks=true, linkcolor=blue, filecolor=magenta, urlcolor=cyan,}
\urlstyle{same}
\usepackage{graphicx}
\usepackage[export]{adjustbox}
\graphicspath{ {./images/} }
\usepackage{caption}
\usepackage{bbold}
\usepackage{multirow}
\usepackage{fvextra, csquotes}
% \usepackage[fallback]{xeCJK}
\usepackage{polyglossia}
\usepackage{fontspec}
\usepackage{newunicodechar}
\IfFontExistsTF{Noto Serif CJK TC}
{\setCJKmainfont{Noto Serif CJK TC}}
{\IfFontExistsTF{STSong}
  {\setCJKmainfont{STSong}}
  {\IfFontExistsTF{Droid Sans Fallback}
    {\setCJKmainfont{Droid Sans Fallback}}
    {\setCJKmainfont{SimSun}}
}}

\setmainlanguage{english}
\IfFontExistsTF{CMU Serif}
{\setmainfont{CMU Serif}}
{\IfFontExistsTF{DejaVu Sans}
  {\setmainfont{DejaVu Sans}}
  {\setmainfont{Georgia}}
}

 Finite-Sample Properties of OLS }
%New command to display footnote whose markers will always be hidden
\let\svthefootnote\thefootnote
\newcommand\blfootnotetext[1]{%
  \let\thefootnote\relax\footnote{#1}%
  \addtocounter{footnote}{-1}%
  \let\thefootnote\svthefootnote%
}
%Overriding the \footnotetext command to hide the marker if its value is `0`
\let\svfootnotetext\footnotetext
\renewcommand\footnotetext[2][?]{%
  \if\relax#1\relax%
    \ifnum\value{footnote}=0\blfootnotetext{#2}\else\svfootnotetext{#2}\fi%
  \else%
    \if?#1\ifnum\value{footnote}=0\blfootnotetext{#2}\else\svfootnotetext{#2}\fi%
    \else\svfootnotetext[#1]{#2}\fi%
  \fi
}
\newunicodechar{×}{\ifmmode\times\else{$\times$}\fi}
\newunicodechar{⋮}{\ifmmode\vdots\else{$\vdots$}\fi}
\newunicodechar{⇔}{\ifmmode\Leftrightarrow\else{$\Leftrightarrow$}\fi}
\newunicodechar{¥}{\ifmmode\text{\yen}\else\yen\fi}
\title{Chapter 6: Serial Correlation}

\author{}
\date{}

\begin{document}
\maketitle

\section*{CHAPTER 6 Serial Correlation}
\begin{abstract}
This chapter extends the GMM model of Chapter 3 to incorporate serial correlation in the product of the vector of instruments and the error term. To do so, we need to generalize the Central Limit Theorems of Chapter 2 to serially correlated processes. The generalization is possible under certain conditions restricting the degree of serial correlation. The condition is transparent for the stochastic processes called linear processes. Linear processes are important in their own right, and this chapter devotes four sections to covering them. In the application section, the extended GMM model is used to test the efficient market hypothesis, this time for the foreign exchange market.
\end{abstract}

A Note on Notation: Because the issues treated here are specific to time-series data, we use the subscript " $t$ " instead of " $i$ " in this chapter.

\subsection*{6.1 Modeling Serial Correlation: Linear Processes}
Section 2.2 introduced covariance-stationary processes and their autocovariances. In particular, a (scalar or univariate) white noise process $\left\{\varepsilon_{t}\right\}$ is a zero-mean covariance-stationary process with no serial correlation:

$$
\mathrm{E}\left(\varepsilon_{t}\right)=0, \mathrm{E}\left(\varepsilon_{t}^{2}\right)=\sigma^{2}>0, \mathrm{E}\left(\varepsilon_{t} \varepsilon_{t-j}\right)=0 \text { for } j \neq 0 .
$$

A very important class of covariance-stationary processes, called linear processes, can be created by taking a moving average of a white noise process. This section studies general properties of linear processes using an apparatus called the filter.

Because the current value of a linear process can depend on possibly infinite past values of a white noise process, it is convenient to have the white noise process and the linear process it generates defined for $t=0,-1,-2, \ldots$ as well\\
as for $t=1,2, \ldots$. Intuitively, a covariance-stationary process defined for all integers ( $t=0, \pm 1, \pm 2, \ldots$ ) is a process that started so long ago that its mean and autocovariances have stabilized to time-invariant constants.

\section*{MA( $q$ )}
The simplest example of linear processes that exhibit serial correlation is a finiteorder moving average process. A process $\left\{y_{t}\right\}$ is called the $\boldsymbol{q}$-th order moving average process (MA(q)) if it can be written as a weighted average of the current and most recent $q$ values of a white noise process:


\begin{equation*}
y_{t}=\mu+\theta_{0} \varepsilon_{t}+\theta_{1} \varepsilon_{t-1}+\cdots+\theta_{q} \varepsilon_{t-q} \quad \text { with } \theta_{0}=1 . \tag{6.1.1}
\end{equation*}


Evidently, a moving average process is covariance-stationary with mean $\mu$. It is easy to verify that the $j$-th order autocovariance, $\gamma_{j}\left(\equiv \mathrm{E}\left[\left(y_{t}-\mu\right)\left(y_{t-j}-\mu\right)\right]\right)$, is


\begin{align*}
& \gamma_{j}=\left(\theta_{j} \theta_{0}+\theta_{j+1} \theta_{1}+\cdots+\theta_{q} \theta_{q-j}\right) \sigma^{2}=\sigma^{2} \sum_{k=0}^{q-j} \theta_{j+k} \theta_{k} \quad \text { for } j=0,1, \ldots, q,  \tag{6.1.2a}\\
& \gamma_{j}=0 \quad \text { for } j>q, \tag{6.1.2b}
\end{align*}


where $\sigma^{2} \equiv \mathrm{E}\left(\varepsilon_{t}^{2}\right)$. (This formula also covers $\gamma_{-j}$ because $\gamma_{j}=\gamma_{-j}$.) For example, for $q=1$, we have

$$
\begin{aligned}
& \gamma_{0}=\left(\theta_{0}^{2}+\theta_{1}^{2}\right) \sigma^{2}=\left(1+\theta_{1}^{2}\right) \sigma^{2} \\
& \quad\left(\text { by setting } j=0 \text { and } q=1 \text { in (6.1.2a) and noting } \theta_{0}=1\right), \\
& \gamma_{1}=\gamma_{-1}=\left(\theta_{1} \theta_{0}\right) \sigma^{2}=\theta_{1} \sigma^{2} \\
& \quad\left(\text { by setting } j=1 \text { and } q=1 \text { in (6.1.2a) and noting } \theta_{0}=1\right), \\
& \gamma_{j}=0 \quad \text { for } j= \pm 2, \pm 3, \ldots
\end{aligned}
$$

The whole profile of autocovariances, $\left\{\gamma_{j}\right\}$, is described by just $q+1$ parameters $\left(\theta_{1}, \theta_{2}, \ldots, \theta_{q}\right.$, and $\sigma^{2}$ ), and the correlogram $\left\{\rho_{j}\right\}$ (where $\rho_{j} \equiv \gamma_{j} / \gamma_{0}$ ) by $q$ parameters $\left(\theta_{1}, \theta_{2}, \ldots, \theta_{q}\right)$.

\section*{MA( $\infty$ ) as a Mean Square Limit}
As illustrated in the above example, serial correlation in MA( $q$ ) processes dies out completely after $q$ lags. Although some time series have this property (we will see an example in the application below), it is certainly desirable to be able to model serial correlation that does not have this property. An obvious idea is to have $y_{t}$\\
depend on the infinite past by replacing the sum of finitely many terms in (6.1.1),

$$
\theta_{0} \varepsilon_{t}+\theta_{1} \varepsilon_{t-1}+\cdots+\theta_{q} \varepsilon_{t-q}
$$

by an infinite sum


\begin{equation*}
\psi_{0} \varepsilon_{t}+\psi_{1} \varepsilon_{t-1}+\cdots=\sum_{j=0}^{\infty} \psi_{j} \varepsilon_{t-j} \tag{6.1.3}
\end{equation*}


where $\left\{\psi_{j}\right\}$ is a sequence of real numbers.\\
For this idea to work, we must make sure that this sum of infinitely many random variables is well defined, namely, that the partial sum


\begin{equation*}
\sum_{j=0}^{n} \psi_{j} \varepsilon_{t-j} \tag{6.1.4}
\end{equation*}


converges (in some appropriate sense) to a random variable as $n \rightarrow \infty$. One popular condition under which this happens is that the sequence of real numbers $\left\{\psi_{j}\right\}$ be absolutely summable:


\begin{equation*}
\sum_{j=0}^{\infty}\left|\psi_{j}\right|<\infty \tag{6.1.5}
\end{equation*}


(This is a standard short-hand expression in calculus. Equation (6.1.5) reads: the partial sum $\sum_{j=0}^{n}\left|\psi_{j}\right|$ converges to a real number [i.e., a finite limit] as $n \rightarrow \infty$ and the infinite sum in (6.1.5) is defined to be this real number.) For the sequence $\left\{\psi_{j}\right\}$ to be absolutely summable, it is necessary (but certainly not sufficient) that $\psi_{j} \rightarrow 0$ as $j \rightarrow \infty$. Thus, absolute summability requires that the effect of the past shocks represented by $\psi_{j}$ eventually die away.

As claimed in the next proposition, the partial sum (6.1.4) converges in mean square to some random variable for any given $t$. The mean square limit is unique in the following sense. If there exists another random variable to which the partial sum converges in mean square, then (as you will show in Analytical Exercise 1) the two mean square limits are equal with probability one. We define the infinite sum (6.1.3) to be the unique mean square limit of the partial sum (6.1.4) and say "the infinite sum converges in mean square."

Proposition 6.1 (MA( $\infty$ ) with absolutely summable coefficients): Let $\left\{\varepsilon_{t}\right\}$ be white noise and $\left\{\psi_{j}\right\}$ be a sequence of real numbers that is absolutely summable. Then\\
(a) For eacht,


\begin{equation*}
y_{t}=\mu+\sum_{j=0}^{\infty} \psi_{j} \varepsilon_{t-j} \tag{6.1.6}
\end{equation*}


converges in mean square. $\left\{y_{t}\right\}$ is covariance-stationary. (The process $\left\{y_{t}\right\}$ is called the infinite-order moving-average process (MA( $\infty$ )).)\\
(b) The mean of $y_{t}$ is $\mu$. The autocovariances $\left\{\gamma_{j}\right\}$ are given by


\begin{align*}
\gamma_{j} & =\left(\psi_{j} \psi_{0}+\psi_{j+1} \psi_{1}+\psi_{j+2} \psi_{2}+\cdots\right) \sigma^{2} \\
& =\sigma^{2} \sum_{k=0}^{\infty} \psi_{j+k} \psi_{k} \quad(j=0,1,2, \ldots) \tag{6.1.7}
\end{align*}


(c) The autocovariances are absolutely summable:


\begin{equation*}
\sum_{j=0}^{\infty}\left|\gamma_{j}\right|<\infty \tag{6.1.8}
\end{equation*}


(d) If, in addition, $\left\{\varepsilon_{t}\right\}$ is i.i.d, then the process $\left\{y_{t}\right\}$ is (strictly) stationary and ergodic.

This result includes MA( $q$ ) processes as a special case with $\psi_{j}=\theta_{j}$ for $j= 0,1, \ldots, q$ and $\psi_{j}=0$ for $j>q$. Evidently, this sequence $\left\{\psi_{j}\right\}$ is absolutely summable. Proving parts (a)-(c) is Analytical Exercise 2. Proving (a) would use the fact from calculus that Cauchy sequences converge (see Review Question 1 about Cauchy sequences). Part (b) is an obvious extrapolation to infinity of the corresponding results on the mean and autocovariances of finite-order MA processes, but the extrapolation involves interchanging the order of expectations and summations. For example, to prove that the mean of the MA( $\infty$ ) process is $\mu$, you would have to show that the expected value of the mean square limit of (6.1.3) is the limit of the expected value of (6.1.4). The fact that this operation is legitimate under mean square convergence would be used in the proof of part (b). Part (c) says, intuitively, that if the effect of past shocks represented by $\psi_{j}$ dies out fast enough to make $\left\{\psi_{j}\right\}$ absolutely summable, then serial correlation dies out just as fast. Absolute summability of autocovariances will play a key role in the theory of covariance-stationary processes to be developed in this chapter. For a proof of (d), see, e.g., Hannan (1970, p. 204, Theorem 3). ${ }^{1}$

\footnotetext{${ }^{1}$ Absolute summability of $\left\{\psi_{j}\right\}$ is more than enough for part (a) to hold. As shown in Analytical Exercise 2, square summability, which is weaker than absolute summability, is enough to ensure that the infinite sum is well defined as a mean square limit. However, for the other parts of this proposition to hold, absolute summability is needed.
}The MA( $\infty$ ) process in (6.1.6) is said to be one-sided because the moving average does not include future values of $\varepsilon_{t}$. It is a special case of a linear process, which is defined to be two-sided moving averages of the form:

$$
y_{t}=\mu+\sum_{j=-\infty}^{\infty} \psi_{j} \varepsilon_{t-j} \text { with } \sum_{j=-\infty}^{\infty}\left|\psi_{j}\right|<\infty
$$

The proposition (and all the results in this section) can easily be extended to general linear processes. We nevertheless focus on one-sided moving averages, only because two-sided moving averages are very rarely encountered in economics.

Proposition 6.1 can be generalized to the case where the process $\left\{y_{t}\right\}$ is a moving average of a general covariance-stationary process rather than of a white noise process. In particular, absolute summability of autocovariances survives the operation of taking an absolutely summable weighted average.

Proposition 6.2 (Filtering covariance-stationary processes): Let $\left\{x_{t}\right\}$ be a covariance-stationary process and $\left\{h_{j}\right\}$ be a sequence of real numbers that is absolutely summable. Then\\
(a) For each $t$, the infinite sum


\begin{equation*}
y_{t}=\sum_{j=0}^{\infty} h_{j} x_{t-j} \tag{6.1.9}
\end{equation*}


converges in mean square. The process $\left\{y_{t}\right\}$ is covariance-stationary.\\
(b) If, furthermore, the autocovariances of $\left\{x_{t}\right\}$ are absolutely summable, then so are the autocovariances of $\left\{y_{t}\right\}$.

We defer to, e.g., Fuller (1996, Theorem 4.3.1) for a proof, but the basic technique is the same as in the proof of Proposition 6.1. Deriving the autocovariances of $\left\{y_{t}\right\}$, which would generalize (6.1.7), is Analytical Exercise 3.

\section*{Filters}
The operation of taking a weighted average of (possibly infinitely many) successive values of a process, as in (6.1.6) and (6.1.9), is called filtering. It can be expressed compactly if we use the device called the lag operator $L$, defined by the relation $L^{j} x_{t}=x_{t-j}$. We now introduce some concepts associated with the lag operator that will be useful in time-series analysis. ${ }^{2}$

\footnotetext{${ }^{2}$ For a fuller treatment of filters, see, e.g., Gourieroux and Monfort (1997, Sections 5.2 and 7.4).
}\section*{Filters Defined}
For a given arbitrary sequence of real numbers, ( $\alpha_{0}, \alpha_{1}, \alpha_{2}, \ldots$ ), define a filter by


\begin{equation*}
\alpha(L) \equiv \alpha_{0}+\alpha_{1} L+\alpha_{2} L^{2}+\cdots . \tag{6.1.10}
\end{equation*}


If we just mechanically and mindlessly apply the definition of the lag operator to an input process $\left\{x_{t}\right\}$, we get


\begin{align*}
\alpha(L) x_{t} & =\alpha_{0} x_{t}+\alpha_{1} L x_{t}+\alpha_{2} L^{2} x_{t}+\cdots \\
& =\alpha_{0} x_{t}+\alpha_{1} x_{t-1}+\alpha_{2} x_{t-2}+\cdots \\
& =\sum_{j=0}^{\infty} \alpha_{j} x_{t-j} \tag{6.1.11}
\end{align*}


This is indeed the definition of $\alpha(L) x_{t}$. If $x_{t}=c$ (a constant), $\alpha(L) c=\alpha(1) c= c \sum_{j=0}^{\infty} \alpha_{j}$. If $\alpha_{j} \neq 0$ for $j=p$ and $\alpha_{j}=0$ for $j>p$, the filter reduces to a $p$-th degree lag polynomial:

$$
\alpha(L)=\alpha_{0}+\alpha_{1} L+\cdots+\alpha_{p} L^{p}
$$

When applied to an input process, it creates a weighted average of the current and $p$ most recent values of the process. Proposition 6.2 assures us that the object $\alpha(L) x_{t}$ defined by (6.1.11) is a well-defined random variable forming a covariancestationary process if the sequence $\left\{\alpha_{j}\right\}$ is absolutely summable and if the input process $\left\{x_{t}\right\}$ is covariance-stationary. A filter with absolutely summable coefficients will be referred to (in this book) as an absolutely summable filter. To use some fancy set theory language, an absolutely summable filter is a mapping from the set of covariance-stationary processes to itself.

\section*{Products of Filters}
Let $\left\{\alpha_{j}\right\}$ and $\left\{\beta_{j}\right\}$ be two arbitrary sequences of real numbers and define the sequence $\left\{\delta_{j}\right\}$ by the relation


\begin{gather*}
\alpha_{0} \beta_{0}=\delta_{0}, \\
\alpha_{0} \beta_{1}+\alpha_{1} \beta_{0}=\delta_{1}, \\
\alpha_{0} \beta_{2}+\alpha_{1} \beta_{1}+\alpha_{2} \beta_{0}=\delta_{2}, \\
\cdots, \\
\alpha_{0} \beta_{j}+\alpha_{1} \beta_{j-1}+\alpha_{2} \beta_{j-2}+\cdots+\alpha_{j-1} \beta_{1}+\alpha_{j} \beta_{0}=\delta_{j}, \\
\cdots, \tag{6.1.12}
\end{gather*}


The sequence $\left\{\delta_{j}\right\}$ created from this convoluted formula is called the convolution of $\left\{\alpha_{j}\right\}$ and $\left\{\beta_{j}\right\}$. Let $\alpha(L) \equiv \alpha_{0}+\alpha_{1} L+\alpha_{2} L^{2}+\cdots, \beta(L) \equiv \beta_{0}+\beta_{1} L+\beta_{2} L^{2}+\cdots$, and $\delta(L)=\delta_{0}+\delta_{1} L+\delta_{2} L^{2}+\cdots$ be the associated filters. We define the product of two filters $\alpha(L)$ and $\beta(L)$ to be this filter $\delta(L)$ and write it as

$$
\delta(L)=\alpha(L) \beta(L)
$$

For example, for $\alpha(L)=1+\alpha_{1} L$ and $\beta(L)=1+\beta_{1} L$, the convolution formula yields

$$
\left(1+\alpha_{1} L\right)\left(1+\beta_{1} L\right)=1+\left(\alpha_{1}+\beta_{1}\right) L+\alpha_{1} \beta_{1} L^{2}
$$

(If you are familiar with the product of two polynomials of the usual sort from calculus, you can immediately see that this definition for filters is strictly analogous.) As is clear from the definition, filters are commutative:

$$
\alpha(L) \beta(L)=\beta(L) \alpha(L)
$$

Commutativity, however, will not carry over to matrix filters (see Section 6.3 below).

If the two sequences $\left\{\alpha_{j}\right\}$ and $\left\{\beta_{j}\right\}$ are absolutely summable and if $\left\{x_{t}\right\}$ is covariance-stationary, then Proposition 6.2 assures us that $\left\{\beta(L) x_{t}\right\}$ is covariancestationary and so $\alpha(L) \beta(L) x_{t}$ is a well-defined random variable, equal to $\delta(L) x_{t}$. Thus, if $\alpha(L)$ and $\beta(L)$ are absolutely summable and $\left\{x_{t}\right\}$ is covariance-stationary, then


\begin{equation*}
" \alpha(L) \beta(L)=\delta(L) " \Rightarrow " \alpha(L) \beta(L) x_{t}=\delta(L) x_{t} . " \tag{6.1.13}
\end{equation*}


(Incidentally, $\left\{\delta_{j}\right\}$ is absolutely summable [see, e.g., Fuller, 1996, p. 30].)

\section*{Inverses}
Of particular interest is the case when $\delta(L)=1$ (a very special filter) so that the two filters $\alpha(L)$ and $\beta(L)$ satisfy $\alpha(L) \beta(L)=1$. We say that $\beta(L)$ is the inverse of $\alpha(L)$ and denote it as $\alpha(L)^{-1}$ or $1 / \alpha(L)$. That is,


\begin{equation*}
\alpha(L)^{-1} \text { is a filter satisfying } \alpha(L) \alpha(L)^{-1}=1 . \tag{6.1.14}
\end{equation*}


As long as $\alpha_{0} \neq 0$ in $\alpha(L)=\alpha_{0}+\alpha_{1} L+\cdots$, the inverse of $\alpha(L)$ can be defined for any arbitrary sequence $\left\{\alpha_{j}\right\}$ because the set of equations (6.1.12) with $\delta_{0}=1$,\\
$\delta_{j}=0(j=1,2, \ldots)$ can be solved successively for $\left\{\beta_{j}\right\}:$

$$
\beta_{0}=\frac{1}{\alpha_{0}}, \quad \beta_{1}=-\frac{\alpha_{1} \beta_{0}}{\alpha_{0}}=-\frac{\alpha_{1}}{\alpha_{0}^{2}}, \text { etc. }
$$

For example, consider the filter $1-L$. Its inverse is given by

$$
(1-L)^{-1}=1+L+L^{2}+L^{3}+\cdots
$$

The coefficient sequence is obviously not absolutely summable. This example illustrates the point that an inverse filter may not be absolutely summable.

It is left as Review Question 3 to show that

\[
\begin{array}{r}
\alpha(L) \alpha(L)^{-1}=\alpha(L)^{-1} \alpha(L) \quad \text { (so inverses are commutative), } \\
\text { " } \alpha(L) \beta(L)=\delta(L) \text { " } \Leftrightarrow \text { " } \beta(L)=\alpha(L)^{-1} \delta(L) \text { " } \Leftrightarrow \text { " } \alpha(L)=\delta(L) \beta(L)^{-1} \text {," } \tag{6.1.15b}
\end{array}
\]

provided $\alpha_{0} \neq 0$ and $\beta_{0} \neq 0$. Therefore, for example, to solve $\alpha(L) \beta(L)=\delta(L)$ for $\beta(L)$, you "multiply from the left" by the inverse of $\alpha(L)$. We have noted above that commutativity does not necessarily hold for matrix filters. As will be shown in Section 6.3, for inverses, commutativity does generalize to matrix filters.

\section*{Inverting Lag Polynomials}
In the next section on ARMA processes, it will become necessary to calculate the inverse of a $p$-th degree lag polynomial $\phi(L)$

$$
\phi(L)=1-\phi_{1} L-\phi_{2} L^{2}-\cdots-\phi_{p} L^{p}
$$

Since $\phi_{0}=1 \neq 0$, the inverse can be defined. Let $\psi(L) \equiv \phi(L)^{-1}$. By definition, it satisfies

$$
\phi(L) \psi(L)=1
$$

The convolution formula (6.1.12) provides an algorithm for solving this for $\psi(L)$. Setting $\alpha(L)=\phi(L), \beta(L)=\psi(L)$, and $\delta(L)=1$ in the formula yields

\[
\begin{array}{rcr}
\text { constant: } & \psi_{0} & =1 \\
L: & \psi_{1}-\phi_{1} \psi_{0} & =0 \\
L^{2}: & \psi_{2}-\phi_{1} \psi_{1}-\phi_{2} \psi_{0} & =0 \\
\cdots & \cdots & \\
L^{p-1}: & \psi_{p-1}-\phi_{1} \psi_{p-2}-\phi_{2} \psi_{p-3}-\cdots-\phi_{p-1} \psi_{0} & =0 \\
L^{p}: & \psi_{p}-\phi_{1} \psi_{p-1}-\phi_{2} \psi_{p-2}-\cdots-\phi_{p-1} \psi_{1}-\phi_{p} \psi_{0}=0 \\
L^{p+1}: & \psi_{p+1}-\phi_{1} \psi_{p}-\phi_{2} \psi_{p-1}-\cdots-\phi_{p-1} \psi_{2}-\phi_{p} \psi_{1}=0 \\
L^{p+2}: & \psi_{p+2}-\phi_{1} \psi_{p+1}-\phi_{2} \psi_{p}-\cdots-\phi_{p-1} \psi_{3}-\phi_{p} \psi_{2}=0 \tag{6.1.16}
\end{array}
\]

These equations can easily be solved successively for ( $\psi_{0}, \psi_{1}, \psi_{2}, \ldots$ ) as

$$
\psi_{0}=1, \quad \psi_{1}=\phi_{1}, \quad \psi_{2}=\phi_{2}+\phi_{1}^{2}, \quad \text { etc. }
$$

Also, notice that for sufficiently large $j$ (actually for $j \geq p$ ), $\left\{\psi_{j}\right\}$ follows the $p$-th order homogeneous difference equation


\begin{equation*}
\psi_{j}-\phi_{1} \psi_{j-1}-\phi_{2} \psi_{j-2}-\cdots-\phi_{p-1} \psi_{j-p+1}-\phi_{p} \psi_{j-p}=0 \tag{6.1.17}
\end{equation*}


So, once the first $p$ coefficients ( $\psi_{0}, \psi_{1}, \ldots, \psi_{p-1}$ ) are calculated, we can use this homogeneous $p$-th order difference equation to generate the rest of the coefficients $\left(\psi_{j}, j \geq p\right)$ with those first $p$ coefficients as the initial condition.

As we know from the theory of homogeneous difference equations (summarized in Analytical Exercise 4), the solution sequence $\left\{\psi_{j}\right\}$ to (6.1.17) eventually starts declining at a geometric rate if what is known as the stability condition holds. The condition states:

All the roots of the $p$-th degree polynomial equation in $z$


\begin{equation*}
\phi(z)=0 \text { where } \phi(z) \equiv 1-\phi_{1} z-\phi_{2} z^{2}-\cdots-\phi_{p} z^{p} \tag{6.1.18}
\end{equation*}


are greater than 1 in absolute value.\\
The roots of the polynomial equation $\phi(z)=0$ can be complex numbers. If your command of complex numbers is rusty, a complex number $z$ can be written as

$$
z=a+b i
$$

where $a$ and $b$ are two real numbers and $i=\sqrt{-1}$. So a complex number can be represented as a point in a two-dimensional diagram with the horizontal axis measuring $a$ (the real part) and the vertical axis measuring $b$ (the imaginary part).

Its absolute value is $|z|=\sqrt{a^{2}+b^{2}}$, which is the distance between the origin and the point representing $z$. We say that $z$ lies outside the unit circle when $|z|>$ 1. So the stability condition above can be restated as requiring that the roots lie outside the unit circle. Because $z=\lambda$ solves $\phi(z)=0$ whenever $z=1 / \lambda$ solves $z^{p}-\phi_{1} z^{p-1}-\phi_{2} z^{p-2}-\cdots-\phi_{p-1} z-\phi_{p}=0$, the stability condition can be equivalently stated as

All the roots of the $p$-th order polynomial equation


\begin{equation*}
z^{p}-\phi_{1} z^{p-1}-\phi_{2} z^{p-2}-\cdots-\phi_{p-1} z-\phi_{p}=0 \tag{$\prime$}
\end{equation*}


are less than 1 in absolute value (i.e., lie inside the unit circle).\\
Under this condition, as you will be asked to prove (Analytical Exercise 4), there exist two real numbers, $A>0$ and $0<b<1$, such that


\begin{equation*}
\left|\psi_{j}\right|<A b^{j} \quad \text { for all } j . \tag{6.1.19}
\end{equation*}


Given (6.1.19), the absolute summability of $\left\{\psi_{j}\right\}$ follows because

$$
\sum_{j=0}^{\infty}\left|\psi_{j}\right|<\sum_{j=0}^{\infty} A b^{j}=\frac{A}{1-b}<\infty
$$

Thus, we have proved

Proposition 6.3 (Absolutely summable inverses of lag polynomials): Consider a $p$-th degree lag polynomial $\phi(L)=1-\phi_{1} L-\phi_{2} L^{2}-\cdots-\phi_{p} L^{p}$, and let $\psi(L) \equiv \phi(L)^{-1}$. If the associated $p$-th degree polynomial equation $\phi(z)=0$ satisfies the stability condition (6.1.18) or (6.1.18'), then the coefficient sequence $\left\{\psi_{j}\right\}$ of $\psi(L)$ is bounded in absolute value by a geometrically declining sequence (as in (6.1.19)) and hence is absolutely summable.

To illustrate, consider a lag polynomial of degree one, $1-\phi L$. The root of the associated polynomial equation $1-\phi z=0$ is $1 / \phi$. The stability condition is that $|1 / \phi|>1$ or $|\phi|<1$. The filter


\begin{equation*}
1+\phi L+\phi^{2} L^{2}+\cdots \tag{6.1.20}
\end{equation*}


is evidently its inverse. The coefficient sequence $\left\{\phi^{j}\right\}$ is indeed bounded in absolute value by a geometrically declining sequence (for example, set $A=1.1$ and $b=\phi$ ).

\section*{QUESTIONS FOR REVIEW}
\begin{enumerate}
  \item (Cauchy sequences) A sequence of real numbers, $\left\{\alpha_{n}\right\}$, is said to be a Cauchy sequence if $\left|\alpha_{m}-\alpha_{n}\right| \rightarrow 0$ as $m, n \rightarrow \infty$. An important result from calculus is that $\alpha_{n}$ converges to a real number if and only if the sequence is Cauchy. Taking this as given, show that $\left\{\gamma_{j}\right\}$ is absolutely summable if and only if $\sum_{j=n+1}^{m}\left|\gamma_{j}\right| \rightarrow 0$ as $m, n \rightarrow \infty$. (You can assume, without loss of generality, that $m>n$.)
  \item (How filtering alters autocovariances) Let $h(L)=h_{0}+h_{1} L$ and $y_{t}=h(L) x_{t}$ where $\left\{x_{t}\right\}$ is covariance-stationary. Verify that (with $\left\{\gamma_{j}^{x}\right\}$ being the autocovariance of $\left\{x_{t}\right\}$ ) the autocovariances of $\left\{y_{t}\right\}$ are given by
\end{enumerate}

$$
\gamma_{j}^{y}=\sum_{k=0}^{1} \sum_{\ell=0}^{1} h_{k} h_{\ell} \gamma_{j-k+\ell}^{x}
$$

\begin{enumerate}
  \setcounter{enumi}{2}
  \item Prove (6.1.15b). Hint: Here is a proof that " $\alpha(L) \beta(L)=\delta(L) " \Rightarrow " \beta(L)= \alpha(L)^{-1} \delta(L) ":$
\end{enumerate}

$$
\alpha(L)^{-1} \delta(L)=\alpha(L)^{-1} \alpha(L) \beta(L)=\alpha(L) \alpha(L)^{-1} \beta(L)=\beta(L)
$$

\begin{enumerate}
  \setcounter{enumi}{3}
  \item Does $\phi(z)=1-\frac{3}{4} z+\frac{9}{16} z^{2}$ satisfy the stability condition? Hint: The roots of $\phi(z)=0$ are $2(1 \pm \sqrt{3} i) / 3$.
  \item (Coefficients of polynomial equations) Consider a polynomial equation $\phi(z)=0$ where $\phi(z)=1-\phi_{1} z-\phi_{2} z^{2}$ and let ( $\lambda_{1}, \lambda_{2}$ ) be the two roots. Verify that $\phi_{1}=1 / \lambda_{1}+1 / \lambda_{2}, \phi_{2}=-1 /\left(\lambda_{1} \cdot \lambda_{2}\right)$. Hint: If $\lambda_{1}$ and $\lambda_{2}$ are the two roots of $\phi(z)=0, \phi(z)$ can be written as $\phi(z)=\left(1-\frac{1}{\lambda_{1}} z\right)\left(1-\frac{1}{\lambda_{2}} z\right)$.
  \item (Inverting $1-\phi L$ ) Verify that (6.1.20) is the inverse of $1-\phi L$ by checking the convolution formula (6.1.12) with $\delta_{0}=1, \delta_{j}=0$ for $j \geq 1$.
\end{enumerate}

\subsection*{6.2 ARMA Processes}
Having introduced the concept of filters, we can study with ease and grace a class of linear processes called ARMA processes, which are a parameterization of the coefficients of MA( $\infty$ ) processes.

\section*{AR(1) and Its MA( $\infty$ ) Representation}
A first-order autoregressive process (AR(1)) satisfies the following stochastic difference equation:


\begin{equation*}
y_{t}=c+\phi y_{t-1}+\varepsilon_{t} \quad \text { or } \quad y_{t}-\phi y_{t-1}=c+\varepsilon_{t} \quad \text { or } \quad(1-\phi L) y_{t}=c+\varepsilon_{t}, \tag{6.2.1}
\end{equation*}


where $\left\{\varepsilon_{t}\right\}$ is white noise. If $\phi \neq 1$, let $\mu \equiv c /(1-\phi)$ and rewrite this equation as


\begin{equation*}
\left(y_{t}-\mu\right)-\phi \cdot\left(y_{t-1}-\mu\right)=\varepsilon_{t} \quad \text { or } \quad(1-\phi L)\left(y_{t}-\mu\right)=\varepsilon_{t} . \tag{$\prime$}
\end{equation*}


As will be seen in a moment, $\mu$, not $c$, is the mean of $y_{t}$ if $y_{t}$ is covariancestationary. For this reason, we will call (6.2.1') a deviation-from-the-mean form. The moving average is on the successive values of $\left\{y_{t}\right\}$, not on $\left\{\varepsilon_{t}\right\}$. The difference equation is called stochastic because of the presence of the random variable $\varepsilon_{t}$.

We seek a covariance-stationary solution $\left\{y_{t}\right\}$ to this stochastic difference equation. The solution depends on whether $|\phi|$ is less than, equal to, or greater than 1 .

Case 1: $|\phi|<1$\\
The solution can be obtained easily by the use of the inverse $(1-\phi L)^{-1}$ given by (6.1.20). Since this filter is absolutely summable when $|\phi|<1$, we can apply it to both sides of the $\operatorname{AR}(1)$ equation (6.2.1') to obtain

$$
(1-\phi L)^{-1}(1-\phi L)\left(y_{t}-\mu\right)=(1-\phi L)^{-1} \varepsilon_{t} .
$$

This operation is legitimate because, thanks to Proposition 6.2, both sides of this equation are well defined as mean square limits. By the definition (6.1.14) of inverses, $(1-\phi L)(1-\phi L)^{-1}=1$, and by commutativity of inverses $(1-\phi L)^{-1}(1- \phi L)=1$. Therefore, if $\left\{y_{t}\right\}$ is covariance-stationary, the left-hand side equals $y_{t}-\mu$ by (6.1.13). So


\begin{align*}
& y_{t}-\mu=(1-\phi L)^{-1} \varepsilon_{t}=\left(1+\phi L+\phi^{2} L^{2}+\cdots\right) \varepsilon_{t}=\sum_{j=0}^{\infty} \phi^{j} \varepsilon_{t-j} \\
& \text { or } y_{t}=\mu+\sum_{j=0}^{\infty} \phi^{j} \varepsilon_{t-j} . \tag{6.2.2}
\end{align*}


What we have shown is that, if $\left\{y_{t}\right\}$ is a covariance-stationary solution to the stochastic difference equation (6.2.1) or (6.2.1'), then $y_{t}$ has the moving-average representation as in (6.2.2). Conversely, if $y_{t}$ has the representation (6.2.2), then\\
it satisfies the difference equation. This can be derived easily by substitution of (6.2.2) into (6.2.1'):

$$
\begin{aligned}
& (1-\phi L)\left(y_{t}-\mu\right)=(1-\phi L)(1-\phi L)^{-1} \varepsilon_{t} \\
& =\varepsilon_{t} \quad\left(\text { since }(1-\phi L)(1-\phi L)^{-1}=1 \text { by definition of inverses }\right) .
\end{aligned}
$$

Thus, the process $\left\{y_{t}\right\}$ given by (6.2.2) is the unique covariance-stationary solution to the first-order stochastic difference equation if $|\phi|<1$. The mean of this process is $\mu$ by Proposition 6.1(b).

The condition $|\phi|<1$, which is the stability condition (6.1.18) associated with the first-degree polynomial equation $1-\phi z=0$, is called the stationarity condition in the present context of autoregressive processes. Intuitively, this result says that the process, if it started a long time ago, "settles down," provided that $|\phi|<1$ so that the effect of the past dies out geometrically as time progresses.

Case 2: $|\phi|>1$\\
By shifting time forward by one period (i.e., by replacing " $t$ " by " $t+1$ "), multiplying both sides by $\phi^{-1}$, and rearranging, the stochastic difference equation (6.2.1') can be written as


\begin{equation*}
y_{t}-\mu=\phi^{-1}\left(y_{t+1}-\mu\right)-\phi^{-1} \varepsilon_{t+1} . \tag{6.2.3}
\end{equation*}


Applying the above argument but this time moving in the opposite direction in time shows that the unique covariance-stationary solution is


\begin{equation*}
y_{t}=\mu-\sum_{j=1}^{\infty} \phi^{-j} \varepsilon_{t+j} . \tag{6.2.4}
\end{equation*}


That is, the current value of $y$ is a moving average of future values of $\varepsilon$. The infinite sum is well defined because the sequence $\left\{\phi^{-j}\right\}$ is absolutely summable if $|\phi|>1$.

Case 3: $|\phi|=1$\\
The stochastic difference equation has no covariance-stationary solution. For example, if $\phi=1$, the stochastic difference equation becomes

$$
\begin{aligned}
y_{t} & =c+y_{t-1}+\varepsilon_{t} \\
& =c+\left(c+y_{t-2}+\varepsilon_{t-1}\right)+\varepsilon_{t} \quad\left(\text { since } y_{t-1}=c+y_{t-2}+\varepsilon_{t-1}\right) \\
& =c+\left(c+\left(c+y_{t-3}+\varepsilon_{t-2}\right)+\varepsilon_{t-1}\right)+\varepsilon_{t}, \text { etc. }
\end{aligned}
$$

Repeating this type of successive substitution $j$ times, we obtain

$$
y_{t}-y_{t-j}=c \cdot j+\left(\varepsilon_{t}+\varepsilon_{t-1}+\cdots+\varepsilon_{t-j+1}\right) .
$$

(If $\left\{\varepsilon_{t}\right\}$ is independent white noise, this process is a random walk with drift $\boldsymbol{c}$.) Suppose, contrary to our claim, that the process is covariance-stationary. Then the variance of $y_{t}-y_{t-j}$ is $2\left(\gamma_{0}-\gamma_{j}\right)$. The variance of the right-hand side is $\sigma^{2} \cdot j$. So

$$
\begin{gathered}
2 \cdot\left(\gamma_{0}-\gamma_{j}\right)=\sigma^{2} \cdot j \text { or } \\
\rho_{j}=1-\frac{\sigma^{2}}{2 \gamma_{0}} \cdot j<-1 \quad \text { for } j \text { large enough }\left(\text { recall: } \rho_{j}=\gamma_{j} / \gamma_{0}\right) .
\end{gathered}
$$

This is a contradiction since autocorrelation coefficients cannot be greater than one in absolute value. So $\left\{y_{t}\right\}$ cannot be covariance-stationary. Processes with $\phi=1$ are examples of "unit-root processes" and will be studied in Chapter 9.

The solution (6.2.4) for the case of $|\phi|>1$, with the current value of $y$ linked to the future values of the forcing process $\left\{\varepsilon_{t}\right\}$, is not useful in economics, which does not confer perfect foresight on economic agents. In what follows, unless otherwise stated, we will use the term "an $\operatorname{AR}(1)$ process" as the unique covariancestationary solution to an $\operatorname{AR}(1)$ equation (6.2.1) or (6.2.1') when the stationarity condition holds.

\section*{Autocovariances of AR(1)}
The autocovariances $\left\{\gamma_{j}\right\}$ of $\operatorname{AR}(1)$ (for the case of $|\phi|<1$ ) can be calculated in two ways. The first method is to represent $\operatorname{AR}(1)$ as MA( $\infty$ ) as in (6.2.2) and use (6.1.7). Since $\psi_{j}=\phi^{j}$ for $\operatorname{AR}(1)$, we have


\begin{align*}
& \gamma_{j}=\left(\phi^{j}+\phi^{j+2}+\phi^{j+4}+\cdots\right) \sigma^{2}=\frac{\phi^{j} \sigma^{2}}{1-\phi^{2}}  \tag{6.2.5}\\
& \text { and } \rho_{j}\left(\equiv \gamma_{j} / \gamma_{0}\right)=\phi^{j} \quad(j=0,1, \ldots)
\end{align*}


So the correlogram of an $\operatorname{AR}(1)$ process is very simple: it declines with $j$ at a constant geometric rate.

The second method bases the calculation on what is called the Yule-Walker equations; see Analytical Exercise 5 for the use of the Yule-Walker equations in the calculation of autocovariances.

\section*{AR(p) and its MA( $\infty$ ) Representation}
All the results just derived for $\operatorname{AR}(1)$ can be generalized to $\operatorname{AR}(p)$, the $\boldsymbol{p}$-th order autoregressive process, which satisfies the $p$-th order stochastic difference\\
equation:


\begin{gather*}
y_{t}=c+\phi_{1} y_{t-1}+\cdots+\phi_{p} y_{t-p}+\varepsilon_{t} \quad \text { or } \\
y_{t}-\phi_{1} y_{t-1}-\cdots-\phi_{p} y_{t-p}=c+\varepsilon_{t} \quad \text { or } \\
\phi(L) y_{t}=c+\varepsilon_{t} \quad \text { with } \phi(L)=1-\phi_{1} L-\phi_{2} L^{2}-\cdots-\phi_{p} L^{p}, \tag{6.2.6}
\end{gather*}


where $\left\{\varepsilon_{t}\right\}$ is white noise and $\phi_{p} \neq 0$. Evidently, $\phi(1)=1-\phi_{1}-\cdots-\phi_{p}$. If $\phi(1) \neq 0$, let $\mu \equiv c /\left(1-\phi_{1}-\cdots-\phi_{p}\right)=c / \phi(1)$. As shown below, $\mu$ is the mean of $y_{t}$ if $y_{t}$ is covariance-stationary. Substituting for $c$, the $\operatorname{AR}(p)$ equation (6.2.6) can be written equivalently in deviation-from-the-mean form:


\begin{gather*}
\left(y_{t}-\mu\right)-\phi_{1} \cdot\left(y_{t-1}-\mu\right)-\cdots-\phi_{p} \cdot\left(y_{t-p}-\mu\right)=\varepsilon_{t} \\
\text { or } \phi(L)\left(y_{t}-\mu\right)=\varepsilon_{t} . \tag{$\prime$}
\end{gather*}


The generalization to $\operatorname{AR}(p)$ of what we have derived for $\operatorname{AR}(1)$ is

\section*{Proposition 6.4 (AR(p) as MA( $\infty$ ) with absolutely summable coefficients):}
Suppose the $p$-th degree polynomial $\phi(z)$ satisfies the stationarity (stability) condition (6.1.18) or (6.1.18'). Then\\
(a) The unique covariance-stationary solution to the $p$-th order stochastic difference equation (6.2.6) or (6.2.6') has the MA( $\infty$ ) representation


\begin{equation*}
y_{t}=\mu+\psi(L) \varepsilon_{t}, \quad \psi(L)=\psi_{0}+\psi_{1} L+\psi_{2} L^{2}+\psi_{3} L^{3}+\cdots, \tag{6.2.7}
\end{equation*}


where $\psi(L)=\phi(L)^{-1}$. The coefficient sequence $\left\{\psi_{j}\right\}$ is bounded in absolute value by a geometrically declining sequence and hence is absolutely summable.\\
(b) The mean $\mu$ of the process is given by


\begin{equation*}
\mu=\phi(1)^{-1} c \text { where } c \text { is the constant in (6.2.6). } \tag{6.2.8}
\end{equation*}


(c) $\left\{\gamma_{j}\right\}$ is bounded in absolute value by a sequence that declines geometrically with $j$. Hence, the autocovariances are absolutely summable.

Unless otherwise stated, we will use the term "an $\operatorname{AR}(p)$ process" as the unique covariance-stationary solution (6.2.7) to an $\operatorname{AR}(p)$ equation (6.2.6) that satisfies the stationarity condition. The absolute summability of the inverse $\phi(L)^{-1}$ in part (a) of this proposition is immediate from Proposition 6.3. Given absolute summability, the first part of (a) can be proved by the same argument we used for AR(1); just\\
multiply both sides of (6.2.6') by $\phi(L)^{-1}$ and observe that

$$
\phi(L)^{-1} \phi(L)=\phi(L) \phi(L)^{-1}=1
$$

Part (b) is immediate from Proposition 6.1(b). Part (c) can easily be proved by combining (a) and Proposition 6.1(b). Analytical Exercise 7 will ask you to give an alternative proof of the proposition without using filters but using the apparatus called the companion form.

As in AR(1), autocovariances can be obtained in two ways. The first method utilizes the MA( $\infty$ ) representation. The coefficients $\left\{\psi_{j}\right\}$ in the MA representation are the coefficients in the inverse of the lag polynomial $\phi(L)$, and the algorithm for calculating the inverse coefficients has already been presented (see (6.1.16)). Given $\left\{\psi_{j}\right\}$, use (6.1.7) to calculate the autocovariances $\left\{\gamma_{j}\right\}$. The second method uses the Yule-Walker equations.

$$
\begin{aligned}
& \operatorname{ARMA}(p, q) \\
& \operatorname{An} \operatorname{ARMA}(p, q) \text { process combines } \operatorname{AR}(p) \text { and MA }(q) \text { : } \\
& \qquad \begin{aligned}
y_{t} & =c+\phi_{1} y_{t-1}+\phi_{2} y_{t-2}+\cdots+\phi_{p} y_{t-p} \\
& +\theta_{0} \varepsilon_{t}+\theta_{1} \varepsilon_{t-1}+\theta_{2} \varepsilon_{t-2}+\cdots+\theta_{q} \varepsilon_{t-q}
\end{aligned}
\end{aligned}
$$

or

$$
\begin{aligned}
& y_{t}-\phi_{1} y_{t-1}-\phi_{2} y_{t-2}-\cdots-\phi_{p} y_{t-p} \\
& \quad=c+\theta_{0} \varepsilon_{t}+\theta_{1} \varepsilon_{t-1}+\theta_{2} \varepsilon_{t-2}+\cdots+\theta_{q} \varepsilon_{t-q}
\end{aligned}
$$

or


\begin{gather*}
\phi(L) y_{t}=c+\theta(L) \varepsilon_{t}, \quad \phi(L)=1-\phi_{1} L-\cdots-\phi_{p} L^{p} \\
\theta(L)=\theta_{0}+\theta_{1} L+\cdots+\theta_{q} L^{q} \tag{6.2.9}
\end{gather*}


where $\left\{\varepsilon_{t}\right\}$ is white noise. If $\phi(1) \neq 1$, set $\mu=c / \phi(1)$. The deviation-from-themean form is

$$
\begin{gathered}
\left(y_{t}-\mu\right)-\phi_{1} \cdot\left(y_{t-1}-\mu\right)-\cdots-\phi_{p} \cdot\left(y_{t-p}-\mu\right) \\
=\theta_{0} \varepsilon_{t}+\theta_{1} \varepsilon_{t-1}+\cdots+\theta_{q} \varepsilon_{t-q}
\end{gathered}
$$

or


\begin{equation*}
\phi(L)\left(y_{t}-\mu\right)=\theta(L) \varepsilon_{t} \tag{$\prime$}
\end{equation*}


This is still a stochastic $p$-th order difference equation but with a serially correlated forcing process $\left\{\theta(L) \varepsilon_{t}\right\}$ instead of a white noise process $\left\{\varepsilon_{t}\right\}$. Proposition 6.4 generalizes easily to

Proposition 6.5 (ARMA( $p, q$ ) as MA( $\infty$ ) with absolutely summable coefficients): Suppose the $p$-th degree polynomial $\phi(z)$ satisfies stationarity (stability) condition (6.1.18) or (6.1.18'). Then\\
(a) The unique covariance-stationary solution to the $p$-th order stochastic difference equation (6.2.9) or (6.2.9') has the MA( $\infty$ ) representation


\begin{equation*}
y_{t}=\mu+\psi(L) \varepsilon_{t}, \quad \psi(L)=\psi_{0}+\psi_{1} L+\psi_{2} L^{2}+\psi_{3} L^{3}+\cdots \tag{6.2.10}
\end{equation*}


where $\psi(L) \equiv \phi(L)^{-1} \theta(L)$. The coefficient sequence $\left\{\psi_{j}\right\}$ is bounded in absolute value by a geometrically declining sequence and hence is absolutely summable.\\
(b) The mean $\mu$ of the process is given by


\begin{equation*}
\mu=\phi(1)^{-1} c \text { where } c \text { is the constant in (6.2.9). } \tag{6.2.11}
\end{equation*}


(c) $\left\{\gamma_{j}\right\}$ is bounded in absolute value by a sequence that declines geometrically with $j$. Hence, the autocovariances are absolutely summable.

Unless otherwise stated, we will use the term "an ARMA process" as the unique solution (6.2.10) to an ARMA equation (6.2.9) or (6.2.9') that satisfies the stationarity condition. The MA( $\infty$ ) representation can be derived easily by multiplying both sides of (6.2.9') by the inverse $\phi(L)^{-1}$ and observing that $\phi(L)^{-1} \phi(L)= \phi(L) \phi(L)^{-1}=1$. For the rest of part (a) and part (b), the only nontrivial part of the proof is to show that $\left\{\psi_{j}\right\}$ is bounded in absolute value by a geometrically declining sequence. For $\operatorname{AR}(p)$ we used the algorithm based on the set of equations (6.1.16) to solve $\phi(L) \psi(L)=1$ for $\psi(L)$. This algorithm can easily be generalized. This time $\psi(L) \equiv \phi(L)^{-1} \theta(L)$ or

$$
\phi(L) \psi(L)=\theta(L)
$$

To solve this for $\psi(L)$, apply the convolution formula (6.1.12) with $\alpha(L)=\phi(L)$, $\beta(L)=\psi(L)$, and $\delta(L)=\theta(L)$ (instead of 1). The resulting set of equations is the same as (6.1.16), except that the right-hand side of the equation for the constant is $\theta_{0}$ rather than 1 , that of the equation for $L$ is $\theta_{1}$ rather than 0 , that of $L^{2}$ is $\theta_{2}$, and\\
so forth, until the equation for $L^{q}$, whose right-hand side is $\theta_{q}$. The right-hand side of the rest of the equations is 0 . As in $\operatorname{AR}(p)$, these equations can easily be solved successively for ( $\psi_{0}, \psi_{1}, \psi_{2}, \ldots$ ) as

$$
\psi_{0}=\theta_{0}=1, \quad \psi_{1}=\theta_{1}+\phi_{1}, \quad \psi_{2}=\theta_{2}+\phi_{2}+\theta_{1} \phi_{1}+\phi_{1}^{2}, \text { etc. }
$$

Again, as in the $\mathrm{AR}(p)$ case, for sufficiently large $j$ (actually, for $j \geq \max (p$, $q+1$ ), as you will verify in Review Question 5), $\left\{\psi_{j}\right\}$ follows the same $p$-th order homogeneous difference equation (6.1.17). So $\left\{\psi_{j}\right\}$ is bounded in absolute value by a geometrically declining sequence and hence is absolutely summable.

Again there are two ways to derive the expression for autocovariances in terms of ( $\phi_{1}, \ldots, \phi_{p}, \theta_{1}, \ldots, \theta_{q}, \sigma^{2}$ ). You can calculate the coefficients $\left\{\psi_{j}\right\}$ in the MA representation as just described and then use (6.1.7). The other method is based on the Yule-Walker equations.

\section*{ARMA $(p, q)$ with Common Roots}
In the ARMA( $p, q$ ) equation (6.2.9'), suppose $\phi(L)$ satisfies the stationarity condition as in Proposition 6.5, but suppose that one of the roots of $\phi(z)=0$ is also a root of $\theta(z)=0$. If that root is $\lambda$, the polynomials $\phi(z)$ and $\theta(z)$ have a factorization

$$
\phi(z)=\alpha(z) \phi^{*}(z), \quad \theta(z)=\alpha(z) \theta^{*}(z) \text { with } \alpha(z)=1-\lambda^{-1} z
$$

The inverse of $\phi(L)$ is $\phi^{*}(L)^{-1} \alpha(L)^{-1}$. So the $\psi(L)$ in Proposition 6.5 can be written as


\begin{equation*}
\psi(L)=\phi(L)^{-1} \theta(L)=\phi^{*}(L)^{-1} \alpha(L)^{-1} \alpha(L) \theta^{*}(L)=\phi^{*}(L)^{-1} \theta^{*}(L) \tag{6.2.12}
\end{equation*}


which shows that the ARMA( $p, q$ ) equation (6.2.9') and the simpler ARMA( $p- 1, q-1$ ) equation

$$
\phi^{*}(L)\left(y_{t}-\mu\right)=\theta^{*}(L) \varepsilon_{t}
$$

share the same process as the unique covariance-stationary solution. ${ }^{3}$ For reasons of parsimony, ARMA equations with common roots are rarely used to parameterize covariance-stationary processes.

\footnotetext{${ }^{3}$ Since we defined filters for sequences of real numbers, the discussion in the text presumes that the root $\lambda$ is real. If $\lambda$ is complex, then there must be another common root which is the complex conjugate of $\lambda$. Let $\lambda_{1}=a+b i$ and $\lambda_{2}=a-b i$ be the pair of complex roots. Then if we define

$$
\alpha(z)=1-\left(\lambda_{1}^{-1}+\lambda_{2}^{-1}\right) z+\frac{1}{\left(\lambda_{1} \cdot \lambda_{2}\right)} z^{2}=1-\frac{2 a}{\left(a^{2}+b^{2}\right)} z+\frac{1}{\left(a^{2}+b^{2}\right)} z^{2},
$$
}
the discussion in the text carries over.

\section*{Invertibility}
For the $\operatorname{ARMA}(p, q)$ equation (6.2.9), if $\phi(z)=0$ and $\theta(z)=0$ have no common roots and if $\theta(z)$ satisfies the stability condition (6.1.18) or (6.1.18'), then the ARMA process is said to be invertible and the stability condition on $\theta(z)$ is called the invertibility condition. ${ }^{4}$ Since $\theta(L)^{-1}$ is absolutely summable if $\theta(z)$ satisfies the invertibility (stability) condition, we can multiply both sides of the ARMA equation (6.2.9) by $\theta(L)^{-1}$ and use the relation $\theta(L)^{-1} c=c \theta(1)^{-1}=c / \theta(1)$ to obtain


\begin{equation*}
\theta(L)^{-1} \phi(L) y_{t}=\frac{c}{\theta(1)}+\varepsilon_{t} \tag{6.2.13}
\end{equation*}


So, the process has the infinite-order autoregressive ( $\operatorname{AR}(\infty)$ ) representation. The derivation does not require the stationarity condition (that $\phi(z)$ satisfies the stability condition). If both $\phi(z)$ and $\theta(z)$ satisfy the stability condition, then the ARMA $(p, q)$ process has both the MA( $\infty$ ) and AR $(\infty)$ representations.

\section*{Autocovariance-Generating Function and the Spectrum}
A particularly useful way to summarize the whole profile of autocovariances, $\left\{\gamma_{j}\right\}$, of a covariance-stationary process $\left\{y_{t}\right\}$ is the autocovariance-generating function:


\begin{align*}
g_{Y}(z) & \equiv \sum_{j=-\infty}^{\infty} \gamma_{j} z^{j} \\
& =\gamma_{0}+\sum_{j=1}^{\infty} \gamma_{j} \cdot\left(z^{j}+z^{-j}\right) \quad\left(\text { since } \gamma_{-j}=\gamma_{j}\right) . \tag{6.2.14}
\end{align*}


The argument of this function ( $z$ ) is a complex scalar. Because the summation involves infinitely many terms, there arises a question as to whether the function is well defined. A familiar result from calculus says that the infinite sum is well defined at least for $|z|=1$ (the unit circle) if $\left\{\gamma_{j}\right\}$ is absolutely summable. For example, for $z=1$,


\begin{equation*}
g_{Y}(1)=\sum_{j=-\infty}^{\infty} \gamma_{j}=\gamma_{0}+2 \sum_{j=1}^{\infty} \gamma_{j} . \tag{6.2.15}
\end{equation*}


\footnotetext{${ }^{4}$ It is easy to confuse invertibility and the existence of an inverse. Since $\theta_{0}=1 \neq 0, \theta(L)$ does have an inverse without the invertibility condition. The inverse, however, may not be absolutely summable unless the invertibility condition is satisfied.
}This is indeed well defined because the infinite sum converges. ${ }^{5}$\\
Any complex number on the unit circle can be represented as


\begin{equation*}
z=\cos (\omega)-i \sin (\omega)=e^{-i \omega} \tag{6.2.16}
\end{equation*}


where $i=\sqrt{-1}$ and $\omega$ is the negative of the radian angle that $z$ makes with the real axis. If the autocovariance-generating function is evaluated at this $z$ and divided by $2 \pi$, the resulting function of $\omega$,


\begin{equation*}
s_{Y}(\omega) \equiv \frac{1}{2 \pi} g_{Y}(\cos (\omega)-i \sin (\omega))=\frac{1}{2 \pi} g_{Y}\left(e^{-i \omega}\right) \tag{6.2.17}
\end{equation*}


is called the (power) spectrum or the spectral density (function) of $\left\{y_{t}\right\}$, and $\omega$ is referred to as the frequency.

It can be shown that for absolutely summable autocovariances, all of the autocovariances can be calculated from the spectrum (hence the term autocovariancegenerating function). So there is a one-to-one mapping between the $\left\{\gamma_{j}\right\}$ sequence and $g_{Y}(z)$ or $s_{Y}(\omega)$. The time domain approach, which concentrates on $\left\{\gamma_{j}\right\}$ and which is what we have been doing, and the frequency domain approach, which is based on the interpretation of the spectrum, are therefore equivalent, although some results can be more easily stated in one approach than in the other. This book will not cover the frequency domain approach any further than this.

If $\left\{y_{t}\right\}$ is white noise, then evidently $g_{Y}(z)$ is constant and equal to the variance. For MA(1), $g_{Y}(z)$ is obtained by simple substitution of (6.1.2) for $q=1$ into the definition (6.2.14):

$$
\begin{aligned}
g_{Y}(z) & =\gamma_{-1} z^{-1}+\gamma_{0}+\gamma_{1} z \\
& =\left(\theta_{1} \sigma^{2}\right) z^{-1}+\left(1+\theta_{1}^{2}\right) \sigma^{2}+\left(\theta_{1} \sigma^{2}\right) z \\
& =\sigma^{2} \cdot\left[\theta_{1} z^{-1}+\left(1+\theta_{1}^{2}\right)+\theta_{1} z\right] \\
& =\sigma^{2} \cdot\left(1+\theta_{1} z\right)\left(1+\theta_{1} z^{-1}\right)
\end{aligned}
$$

So


\begin{equation*}
g_{Y}(z)=\sigma^{2} \theta(z) \theta\left(z^{-1}\right) \tag{6.2.18}
\end{equation*}


where $\theta(z)=1+\theta_{1} z$. This last expression generalizes to the MA( $q$ ) case where $\theta(z)=1+\theta_{1} z+\theta_{2} z^{2}+\cdots+\theta_{q} z^{q}$. (You should verify this by carrying out the multiplication in this expression for $g_{Y}(z)$ for MA $(q)$, collecting terms by powers of $z$, and looking at the coefficient of $z^{j}$ or $z^{-j}$.) The expression is also good for

\footnotetext{${ }^{5}$ Fact: If $\left\{\gamma_{j}\right\}$ is absolutely summable, then it is summable (i.e., the partial sum of $\left\{\gamma_{j}\right\}$ converges to a finite limit).
}MA( $\infty$ ) processes if the coefficients are absolutely summable, as the following result assures us.

\section*{Proposition 6.6 (Autocovariance-generating function for filtered processes):}
Let $\left\{\varepsilon_{t}\right\}$ be white noise and let $\psi(L)=\psi_{0}+\psi_{1} L+\psi_{2} L^{2}+\cdots$ be an absolutely summable filter (i.e., with $\left\{\psi_{j}\right\}$ absolutely summable). Then the autocovariancegenerating function of the MA $(\infty)$ process $\left\{y_{t}\right\}$ where $y_{t}=\mu+\psi(L) \varepsilon_{t}$ is


\begin{equation*}
g_{Y}(z)=\sigma^{2} \psi(z) \psi\left(z^{-1}\right) \tag{6.2.19}
\end{equation*}


More generally, let $\left\{x_{t}\right\}$ be covariance-stationary with absolutely summable autocovariances and $g_{X}(z)$ be the autocovariance-generating function of $\left\{x_{t}\right\}$. The autocovariance-generating function of the filtered series $\left\{y_{t}\right\}$ where $y_{t}=h(L) x_{t}$ (with absolutely summable $\left\{h_{j}\right\}$ ) is


\begin{equation*}
g_{Y}(z)=h(z) g_{X}(z) h\left(z^{-1}\right) \tag{6.2.20}
\end{equation*}


We will not prove this result; see, e.g., Theorem 4.3.1 of Fuller (1996).\\
Because $\operatorname{AR}(p)$ and $\operatorname{ARMA}(p, q)$ processes have the MA( $\infty$ ) representations, their autocovariance-generating functions can be derived easily from (6.2.19). Since $\psi(L)=1 / \phi(L)$ for $\operatorname{AR}(p)$, and $\psi(L)=\theta(L) / \phi(L)$ for ARMA $(p, q)$, we have


\begin{align*}
\operatorname{AR}(p): g_{Y}(z) & =\sigma^{2} \frac{1}{\phi(z) \phi\left(z^{-1}\right)}  \tag{6.2.21}\\
\operatorname{ARMA}(p, q): g_{Y}(z) & =\sigma^{2} \frac{\theta(z) \theta\left(z^{-1}\right)}{\phi(z) \phi\left(z^{-1}\right)} \tag{6.2.22}
\end{align*} .


\section*{QUESTIONS FOR REVIEW}
\begin{enumerate}
  \item If $\left\{y_{t}\right\}$ is the covariance-stationary solution to the $\mathrm{AR}(1)$ equation (6.2.1) with $|\phi|<1$, what is
\end{enumerate}

$$
\widehat{\mathrm{E}}^{*}\left(y_{t} \mid 1, y_{t-1}\right)
$$

(the least squares projection)? Hint: $\mathrm{E}\left(y_{t-1} \varepsilon_{t}\right)=0$. Do the projection coefficients depend on $t$ ? Is $\widehat{\mathrm{E}}^{*}\left(y_{t} \mid 1, y_{t-1}\right)=\mathrm{E}\left(y_{t} \mid y_{t-1}\right)$ ? Hint: $\mathrm{E}\left(y_{t-1} \varepsilon_{t}\right)=0$ but does it mean $\mathrm{E}\left(\varepsilon_{t} \mid y_{t-1}\right)=0$ ? What is $\widehat{\mathrm{E}}^{*}\left(y_{t} \mid 1, y_{t-1}, y_{t-2}\right)$ ? Now suppose that $\left\{y_{t}\right\}$ is the covariance-stationary solution to the $\operatorname{AR}(1)$ equation when $|\phi|>1$. Is it true that $\mathrm{E}\left(y_{t-1} \varepsilon_{t}\right)=0$ ? [Answer: No.] Does $\widehat{\mathrm{E}}^{*}\left(y_{t} \mid\right. \left.1, y_{t-1}\right)=c+\phi y_{t-1}$ as was the case when $|\phi|<1$ ?\\
2. (AR(1) with $\phi=-1$ ) For simplicity, set $c=0$ in the $\mathrm{AR}(1)$ equation (6.2.1). Show that the $\operatorname{AR}(1)$ equation with $\phi=-1$ has no covariance-stationary solution. Hint: The same sort of successive substitution we did for the case $\phi=1$ yields

$$
y_{t}-(-1)^{j} y_{t-j}=\varepsilon_{t}-\varepsilon_{t-1}+\cdots+(-1)^{j-1} \varepsilon_{t-j+1} .
$$

From this, derive

$$
(-1)^{j} \rho_{j}=1-\frac{\sigma^{2}}{2 \gamma_{0}} \cdot j
$$

\begin{enumerate}
  \setcounter{enumi}{2}
  \item (About Proposition 6.4) How do we know that $\phi(1) \neq 0$ in (6.2.8)? Hint: Use the stationarity condition. Prove (b) of Proposition 6.4. Hint: Take the expectation of both sides of (6.2.6) and exploit the covariance-stationarity of $\left\{y_{t}\right\}$.
  \item (About Proposition 6.5) Verify that $y_{t}-\mu=\phi(L)^{-1} \theta(L) \varepsilon_{t}$ is a solution to the ARMA $(p, q)$ equation.
  \item ( $\left\{\psi_{j}\right\}$ for $\left.\operatorname{ARMA}(p, q)\right)$ For $\operatorname{AR}(p)$, we used the equations (6.1.12) to calculate the MA coefficients $\left\{\psi_{j}\right\}$. For ARMA(3, 1), write down the corresponding equations. From which lag $j$ does $\left\{\psi_{j}\right\}$ start following
\end{enumerate}

$$
\psi_{j}-\phi_{1} \psi_{j-1}-\phi_{2} \psi_{j-2}-\phi_{3} \psi_{j-3}=0
$$

which is (6.1.17) for $p=3$ ? Do the same for ARMA(1,2). From which $j$ does $\left\{\psi_{j}\right\}$ start following (6.1.17) for $p=1$ ? [Answer: $j=3$.]\\
6. (How filtering alters the spectrum) Let $h(L)$ be absolutely summable. By Proposition 6.2, if $\left\{x_{t}\right\}$ is covariance-stationary with absolutely summable autocovariances, then so is $\left\{y_{t}\right\}$ where $y_{t}=h(L) x_{t}$. Using (6.2.20), verify that

$$
s_{Y}(\omega)=h\left(e^{-i \omega}\right) s_{X}(\omega) h\left(e^{i \omega}\right) .
$$

\begin{enumerate}
  \setcounter{enumi}{6}
  \item (The spectrum is real valued) Show that the spectrum is real valued. Hint: Substitute (6.2.16) into (6.2.14). Use the following facts from complex analysis: $[\cos (\omega)-i \sin (\omega)]^{j}=\cos (j \omega)-i \sin (j \omega) ; \sin (-\omega)=-\sin (\omega)$.
\end{enumerate}

\subsection*{6.3 Vector Processes}
It is straightforward to extend all the concepts introduced so far to vector processes.\\
A vector white noise process $\left\{\varepsilon_{t}\right\}$ is a jointly covariance-stationary process satisfying


\begin{gather*}
\mathrm{E}\left(\boldsymbol{\varepsilon}_{t}\right)=\mathbf{0}, \quad \mathrm{E}\left(\boldsymbol{\varepsilon}_{t} \boldsymbol{\varepsilon}_{t}^{\prime}\right)=\boldsymbol{\Omega}(\text { positive definite }), \quad \text { and } \\
\mathrm{E}\left(\boldsymbol{\varepsilon}_{t} \boldsymbol{\varepsilon}_{t-j}^{\prime}\right)=\mathbf{0} \text { for } j \neq 0 . \tag{6.3.1}
\end{gather*}


Since $\Omega$ is not restricted to being diagonal, there can be contemporaneous correlation between the elements of $\varepsilon_{t}$. Perfect correlation among the elements of $\varepsilon_{t}$ is ruled out because $\boldsymbol{\Omega}$ is required to be positive definite.

A vector MA( $\boldsymbol{\infty}$ ) process is the obvious vector version of (6.1.6):


\begin{equation*}
\mathbf{y}_{t}=\boldsymbol{\mu}+\sum_{j=0}^{\infty} \boldsymbol{\Psi}_{j} \boldsymbol{\varepsilon}_{t-j} \quad \text { with } \boldsymbol{\Psi}_{0}=\mathbf{I} \tag{6.3.2}
\end{equation*}


where $\left\{\boldsymbol{\Psi}_{j}\right\}$ are square matrices. The sequence of coefficient matrices is said to be absolutely summable if each element is absolutely summable. That is, if $\psi_{k \ell j}$ is the $(k, \ell)$ element of $\boldsymbol{\Psi}_{j}$,


\begin{equation*}
"\left\{\boldsymbol{\Psi}_{j}\right\} \text { is absolutely summable" } \Leftrightarrow " \sum_{j=0}^{\infty}\left|\psi_{k \ell j}\right|<\infty \text { for all }(k, \ell) . " \tag{6.3.3}
\end{equation*}


With this definition, Proposition 6.1 generalizes to the multivariate case in an obvious way. In particular, if $\boldsymbol{\Gamma}_{j}\left(\equiv \mathrm{E}\left[\left(\mathbf{y}_{t}-\boldsymbol{\mu}\right)\left(\mathbf{y}_{t-j}-\boldsymbol{\mu}\right)^{\prime}\right]\right)$ is the $j$-th order autocovariance matrix, then the expression for autocovariances in part (b) of Proposition 6.1 becomes


\begin{equation*}
\boldsymbol{\Gamma}_{j}=\sum_{k=0}^{\infty} \boldsymbol{\Psi}_{j+k} \boldsymbol{\Omega} \boldsymbol{\Psi}_{k}^{\prime} \quad(j=0,1,2, \ldots) \tag{6.3.4}
\end{equation*}


(This formula also covers $j=-1,-2, \ldots$ because $\boldsymbol{\Gamma}_{-j}=\boldsymbol{\Gamma}_{j}^{\prime}$.)\\
A multivariate filter can be written as

$$
\mathbf{H}(L)=\mathbf{H}_{0}+\mathbf{H}_{1} L+\mathbf{H}_{2} L^{2}+\cdots
$$

where $\left\{\mathbf{H}_{j}\right\}$ is a sequence of (not necessarily square) matrices. So if $H_{k \ell j}$ is the\\
$(k, \ell)$ element of $\mathbf{H}_{j}$, the $(k, \ell)$ element of $\mathbf{H}(L)$ is

$$
H_{k \ell 0}+H_{k \ell 1} L+H_{k \ell 2} L^{2}+\cdots
$$

The multivariate version of Proposition 6.2 is obvious, with $\mathbf{y}_{t}=\mathbf{H}(L) \mathbf{x}_{t}$.

Product of Filters\\
Let $\mathbf{A}(L)$ and $\mathbf{B}(L)$ be two filters where $\left\{\mathbf{A}_{j}\right\}$ is $m \times r$ and $\left\{\mathbf{B}_{j}\right\}$ is $r \times s$ so that the matrix product $\mathbf{A}_{j} \mathbf{B}_{k}$ can be defined. The product of two filters, $\mathbf{D}(L)= \mathbf{A}(L) \mathbf{B}(L)$ is an $m \times s$ filter whose coefficient matrix sequence $\left\{\mathbf{D}_{j}\right\}$ is given by the multivariate version of the convolution formula (6.1.12):


\begin{gather*}
\mathbf{A}_{0} \mathbf{B}_{0}=\mathbf{D}_{0} \\
\mathbf{A}_{0} \mathbf{B}_{1}+\mathbf{A}_{1} \mathbf{B}_{0}=\mathbf{D}_{1}, \\
\mathbf{A}_{0} \mathbf{B}_{2}+\mathbf{A}_{1} \mathbf{B}_{1}+\mathbf{A}_{2} \mathbf{B}_{0}=\mathbf{D}_{2}, \\
\cdots, \\
\mathbf{A}_{0} \mathbf{B}_{j}+\mathbf{A}_{1} \mathbf{B}_{j-1}+\mathbf{A}_{2} \mathbf{B}_{j-2}+\cdots+\mathbf{A}_{j-1} \mathbf{B}_{1}+\mathbf{A}_{j} \mathbf{B}_{0}=\mathbf{D}_{j}, \\
\cdots \tag{6.3.5}
\end{gather*}


Inverses\\
Let $\mathbf{A}(L)$ and $\mathbf{B}(L)$ be two filters whose coefficient matrices are square. $\mathbf{B}(L)$ is said to be the inverse of $\mathbf{A}(L)$ and is denoted $\mathbf{A}(L)^{-1}$ if


\begin{equation*}
\mathbf{A}(L) \mathbf{B}(L)=\mathbf{I} . \tag{6.3.6}
\end{equation*}


For any arbitrary sequence of square matrices $\left\{\mathbf{A}_{j}\right\}$, the inverse of $\mathbf{A}(L)$ exists if $\mathbf{A}_{0}$ is nonsingular. It is easy to show that inverses are commutative (see Review Question 2), so

$$
\mathbf{A}(L) \mathbf{A}(L)^{-1}=\mathbf{A}(L)^{-1} \mathbf{A}(L) .
$$

Absolutely Summable Inverses of Lag Polynomials\\
A $\boldsymbol{p}$-th degree lag matrix polynomial is


\begin{equation*}
\boldsymbol{\Phi}(L)=\mathbf{I}_{r}-\boldsymbol{\Phi}_{1} L-\boldsymbol{\Phi}_{2} L^{2}-\cdots-\boldsymbol{\Phi}_{p} L^{p}, \tag{6.3.7}
\end{equation*}


where $\left\{\boldsymbol{\Phi}_{j}\right\}$ is a sequence of $r \times r$ matrices with $\boldsymbol{\Phi}_{p} \neq \mathbf{0}$. From the theory of homogeneous difference equations, the stability condition for the multivariate case is as follows:

All the roots of the determinantal equation


\begin{equation*}
\left|\mathbf{I}-\boldsymbol{\Phi}_{1} z-\cdots-\boldsymbol{\Phi}_{p} z^{p}\right|=0 \tag{6.3.8}
\end{equation*}


are greater than 1 in absolute value (i.e., lie outside the unit circle).\\
Or, equivalently,\\
All the roots of the determinantal equation


\begin{equation*}
\left|\mathbf{I} z^{p}-\boldsymbol{\Phi}_{1} z^{p-1}-\cdots-\boldsymbol{\Phi}_{p}\right|=0 \tag{$\prime$}
\end{equation*}


are less than 1 in absolute value (i.e., lie inside the unit circle).\\
As an example, consider a first-degree lag polynomial $\boldsymbol{\Phi}(L)=\mathbf{I}-\boldsymbol{\Phi}_{1} L$ where $\boldsymbol{\Phi}_{\mathbf{I}}=\left(\phi_{k \ell}\right)$ is $2 \times 2$. Equation (6.3.8) can be written as

$$
\left|\begin{array}{cc}
1-\phi_{11} z & -\phi_{12} z \\
-\phi_{21} z & 1-\phi_{22} z
\end{array}\right|=1-\left(\phi_{11}+\phi_{22}\right) z+\left(\phi_{11} \phi_{22}-\phi_{21} \phi_{12}\right) z^{2}=0 .
$$

With the stability condition thus generalized, Proposition 6.3 generalizes in an obvious way. Let $\boldsymbol{\Psi}(L)=\boldsymbol{\Phi}(L)^{-1}$. Each component of the coefficient matrix sequence $\left\{\boldsymbol{\Psi}_{j}\right\}$ will be bounded in absolute value by a geometrically declining sequence.

The multivariate analogue of an $\operatorname{AR}(p)$ is a vector autoregressive process of $\boldsymbol{p}$-th order ( $\operatorname{VAR}(\boldsymbol{p})$ ). It is the unique covariance-stationary solution under stationarity condition (6.3.8) to the following vector stochastic difference equation:

$$
\mathbf{y}_{t}-\boldsymbol{\Phi}_{1} \mathbf{y}_{t-1}-\cdots-\boldsymbol{\Phi}_{p} \mathbf{y}_{t-p}=\mathbf{c}+\boldsymbol{\varepsilon}_{t} \quad \text { or } \quad \boldsymbol{\Phi}(L)\left(\mathbf{y}_{t}-\boldsymbol{\mu}\right)=\boldsymbol{\varepsilon}_{t}
$$

where $\boldsymbol{\Phi}(L)=\mathbf{I}_{r}-\boldsymbol{\Phi}_{1} L-\boldsymbol{\Phi}_{2} L^{2}-\cdots-\boldsymbol{\Phi}_{p} L^{p} \quad$ and $\quad \boldsymbol{\mu}=\boldsymbol{\Phi}(1)^{-1} \mathbf{c}$,\\
where $\boldsymbol{\Phi}_{p} \neq \mathbf{0}$. For example, a bivariate VAR(1) can be written out in full as a two-equation system with common regressors:

$$
\begin{aligned}
& y_{1 t}=c_{1}+\phi_{11} y_{1, t-1}+\phi_{12} y_{2, t-1}+\varepsilon_{1 t}, \\
& y_{2 t}=c_{2}+\phi_{21} y_{1, t-1}+\phi_{22} y_{2, t-1}+\varepsilon_{2 t}, \quad \text { where } \boldsymbol{\Phi}_{1}=\left[\begin{array}{ll}
\phi_{11} & \phi_{12} \\
\phi_{21} & \phi_{22}
\end{array}\right] .
\end{aligned}
$$

More generally, an $M$-variate $\operatorname{VAR}(p)$ is a set of $M$ equations with $M p+1$ common regressors. Proposition 6.4 generalizes straightforwardly to the multivariate case.

In particular, each element of $\Gamma_{j}$ is bounded in absolute value by a geometrically declining sequence.

Vector autoregressions are a very popular tool for analyzing the dynamic interrelationship between key macro variables. A thorough treatment can be found in Hamilton (1994, Chapter 11).

The multivariate analogue of an $\operatorname{ARMA}(p, q)$ is a vector $\operatorname{ARMA}(p, q)$ (VARMA( $\boldsymbol{p}, \boldsymbol{q}$ )). It is the unique covariance-stationary solution under the stationarity condition to the following vector stochastic difference equation:

$$
\begin{gathered}
\mathbf{y}_{t}=\mathbf{c}+\boldsymbol{\Phi}_{1} \mathbf{y}_{t-1}+\boldsymbol{\Phi}_{2} \mathbf{y}_{t-2}+\cdots+\boldsymbol{\Phi}_{p} \mathbf{y}_{t-p}+\boldsymbol{\varepsilon}_{t} \\
+\boldsymbol{\Theta}_{1} \boldsymbol{\varepsilon}_{t-1}+\boldsymbol{\Theta}_{2} \boldsymbol{\varepsilon}_{t-2}+\cdots+\boldsymbol{\Theta}_{q} \boldsymbol{\varepsilon}_{t-q} \text { or } \\
\boldsymbol{\Phi}(L)\left(\mathbf{y}_{t}-\boldsymbol{\mu}\right)=\boldsymbol{\Theta}(L) \boldsymbol{\varepsilon}_{t},
\end{gathered}
$$

where


\begin{gather*}
\boldsymbol{\Phi}(L)=\mathbf{I}-\boldsymbol{\Phi}_{1} L-\cdots-\boldsymbol{\Phi}_{p} L^{p}, \quad \boldsymbol{\Theta}(L)=\mathbf{I}+\boldsymbol{\Theta}_{1} L+\cdots+\boldsymbol{\Theta}_{q} L^{q}, \\
\boldsymbol{\mu}=\boldsymbol{\Phi}(1)^{-1} \mathbf{c}, \tag{6.3.10}
\end{gather*}


where $\boldsymbol{\Phi}_{p} \neq \mathbf{0}$ and $\boldsymbol{\Theta}_{q} \neq \mathbf{0}$. Proposition 6.5 generalizes in an obvious way.

\section*{Autocovariance Generating Function}
The multivariate version of the autocovariance-generating function for a vector covariance-stationary process $\left\{\mathbf{y}_{t}\right\}$ is


\begin{equation*}
\mathbf{G}_{\mathbf{Y}}(z)=\sum_{j=-\infty}^{\infty} \boldsymbol{\Gamma}_{j} z^{j}=\boldsymbol{\Gamma}_{0}+\sum_{j=1}^{\infty}\left(\boldsymbol{\Gamma}_{j} z^{j}+\boldsymbol{\Gamma}_{j}^{\prime} z^{-j}\right) \quad\left(\text { since } \boldsymbol{\Gamma}_{-j}=\boldsymbol{\Gamma}_{j}^{\prime}\right) \tag{6.3.11}
\end{equation*}


The spectrum of a vector process $\left\{\mathbf{y}_{t}\right\}$ is defined as


\begin{equation*}
\mathbf{s}_{\mathbf{Y}}(\omega)=\frac{1}{2 \pi} \mathbf{G}_{\mathbf{Y}}\left(e^{-i \omega}\right) \tag{6.3.12}
\end{equation*}


Proposition 6.6 generalizes easily. In particular, if $\mathbf{H}(L)$ is an $r \times s$ absolutely summable filter and $\mathbf{G}_{\mathbf{X}}(z)$ is the autocovariance-generating function of an $s$ dimensional covariance-stationary process $\left\{\mathbf{x}_{t}\right\}$ with absolutely summable autocovariances, then the autocovariance-generating function of $\mathbf{y}_{t}=\mathbf{H}(L) \mathbf{x}_{t}$ is given by


\begin{equation*}
\underset{(r \times r)}{\mathbf{G}_{\mathbf{Y}}(z)}=\underset{(r \times s)}{\mathbf{H}(z)} \underset{(s \times s)}{\mathbf{G}_{\mathbf{x}}(z)} \underset{(s \times r)}{\mathbf{H}\left(z^{-1}\right)^{\prime}} . \tag{6.3.13}
\end{equation*}


As special cases of this, we have


\begin{equation*}
\text { vector white noise: } \mathbf{G}_{\mathbf{Y}}(z)=\boldsymbol{\Omega} \text {, } \tag{6.3.14}
\end{equation*}



\begin{gather*}
\operatorname{VMA}(\infty): \mathbf{G}_{\mathbf{Y}}(z)=\boldsymbol{\Psi}(z) \boldsymbol{\Omega} \boldsymbol{\Psi}\left(z^{-1}\right)^{\prime}  \tag{6.3.15}\\
\operatorname{VAR}(p): \mathbf{G}_{\mathbf{Y}}(z)=\left[\boldsymbol{\Phi}(z)^{-1}\right] \boldsymbol{\Omega}\left[\boldsymbol{\Phi}\left(z^{-1}\right)^{-1}\right]^{\prime} \tag{6.3.16}
\end{gather*}



\begin{equation*}
\operatorname{VARMA}(p, q): \mathbf{G}_{\mathbf{Y}}(z)=\left[\boldsymbol{\Phi}(z)^{-1}\right][\boldsymbol{\Theta}(z)] \boldsymbol{\Omega}\left[\boldsymbol{\Theta}\left(z^{-1}\right)\right]^{\prime}\left[\boldsymbol{\Phi}\left(z^{-1}\right)^{-1}\right]^{\prime} . \tag{6.3.17}
\end{equation*}


For example, if $\mathbf{y}_{t}=\boldsymbol{\varepsilon}_{t}+\boldsymbol{\Theta}_{1} \boldsymbol{\varepsilon}_{t-1}$,

$$
\mathbf{G}_{\mathbf{Y}}(z)=\left(\mathbf{I}+\boldsymbol{\Theta}_{1} z\right) \boldsymbol{\Omega}\left(\mathbf{I}+\boldsymbol{\Theta}_{1}^{\prime} z^{-1}\right) .
$$

\section*{QUESTIONS FOR REVIEW}
\begin{enumerate}
  \item Verify (6.3.4) for vector MA(1).
  \item (Commutativity of inverses) Verify the commutativity of inverses, that is, "A $(L) \mathbf{B}(L)=\mathbf{I}$ " ⇔ "B $(L) \mathbf{A}(L)=\mathbf{I}$ " for two square filters of dimension $r$ with $\left|\mathbf{A}_{0}\right| \neq 0$ and $\left|\mathbf{B}_{0}\right| \neq 0$. Hint: Set $\mathbf{D}_{0}=\mathbf{I}, \mathbf{D}_{j}=\mathbf{0}(j \geq 1)$ in (6.3.5) and solve for B's. Show that this $\mathbf{B}(L)$ satisfies $\mathbf{B}(L) \mathbf{A}(L)=\mathbf{I}$.
  \item (Lack of commutativity for multivariate filters) Provide an example where $\mathbf{A}(L)$ and $\mathbf{B}(L)$ are $2 \times 2$ but $\mathbf{A}(L) \mathbf{B}(L) \neq \mathbf{B}(L) \mathbf{A}(L)$. Hint: How about $\mathbf{A}(L)=\mathbf{I}+\mathbf{A}_{1} L$ and $\mathbf{B}(L)=\mathbf{I}+\mathbf{B}_{1} L$ ? Find two square matrices $\mathbf{A}_{1}$ and $\mathbf{B}_{1}$ such that $\mathbf{A}_{1} \mathbf{B}_{1} \neq \mathbf{B}_{1} \mathbf{A}_{1}$.
  \item Verify the multivariate version of (6.1.15b), namely,
\end{enumerate}

$$
\begin{gathered}
" \mathbf{A}(L) \mathbf{B}(L)=\mathbf{D}(L) " \Leftrightarrow " \mathbf{B}(L)=\mathbf{A}(L)^{-1} \mathbf{D}(L) " \\
\Leftrightarrow " \mathbf{A}(L)=\mathbf{D}(L) \mathbf{B}(L)^{-1}, "
\end{gathered}
$$

provided $\mathbf{A}_{0}$ and $\mathbf{B}_{0}$ are nonsingular.\\
So, to solve $\mathbf{A}(L) \mathbf{B}(L)=\mathbf{D}(L)$ for $\mathbf{B}(L)$, "multiply both sides from the left" by $\mathbf{A}(L)^{-1}$.\\
5. Show that $[\mathbf{A}(L) \mathbf{B}(L)]^{-1}=\mathbf{B}(L)^{-1} \mathbf{A}(L)^{-1}$, provided that $\mathbf{A}_{0}$ and $\mathbf{B}_{0}$ are nonsingular. Hint: By definition, $\mathbf{A}(L) \mathbf{B}(L)[\mathbf{A}(L) \mathbf{B}(L)]^{-1}=\mathbf{I}$.

\subsection*{6.4 Estimating Autoregressions}
Autoregressive processes (autoregressions) are popular in econometrics, not only because they have a natural interpretation but also because they are easy to estimate. This section considers the estimation of autoregressions. Estimation of ARMA processes will be touched upon at the end.

\section*{Estimation of AR(1)}
Recall that an $\operatorname{AR}(1)$ is an MA( $\infty$ ) with absolutely summable coefficients. So if $\left\{\varepsilon_{t}\right\}$ is independent white noise (an i.i.d. sequence with zero mean and finite variance), it is strictly stationary and ergodic (Proposition 6.1(d)). Letting $\mathbf{x}_{t}= \left(1, y_{t-1}\right)^{\prime}$ and $\boldsymbol{\beta}=(c, \boldsymbol{\phi})^{\prime}$, the $\mathrm{AR}(1)$ equation (6.2.1) can be written as a regression equation


\begin{equation*}
y_{t}=\mathbf{x}_{t}^{\prime} \boldsymbol{\beta}+\varepsilon_{t} . \tag{6.4.1}
\end{equation*}


We assume the sample is $\left(y_{0}, y_{1}, y_{2}, \ldots, y_{n}\right), y_{0}$ inclusive, so that (6.4.1) for $t=1$ (which has $y_{0}$ on the right-hand side) can be included in the estimation. So the sample period is from $t=1$ to $n$.

We now show that all the conditions of Proposition 2.5 about the asymptotic properties of the OLS estimator of $\boldsymbol{\beta}$ with conditional homoskedasticity are satisfied here. First, obviously, linearity (Assumption 2.1) is satisfied. Second, as just noted, $\left\{y_{t}, \mathbf{x}_{t}\right\}$ is jointly stationary and ergodic (Assumption 2.2). Third, since $y_{t-1}$ (the nonconstant regressor in $\mathbf{x}_{t}$ ) is a function of $\left\{\varepsilon_{t-1}, \varepsilon_{t-2}, \ldots\right\}$, it is independent of the error term $\varepsilon_{t}$ by the i.i.d. assumption for $\left\{\varepsilon_{t}\right\}$. Thus


\begin{equation*}
\mathrm{E}\left(\varepsilon_{t}^{2} \mid \mathbf{x}_{t}\right)=\sigma^{2} \tag{6.4.2}
\end{equation*}


that is, the error is conditionally homoskedastic (Assumption 2.7). Fourth, if $\mathbf{g}_{t} \equiv \mathbf{x}_{t} \cdot \varepsilon_{t}=\left(\varepsilon_{t}, y_{t-1} \varepsilon_{t}\right)^{\prime},\left\{\mathbf{g}_{t}\right\}$ is a martingale difference sequence, as required in Assumption 2.5, because

\begin{verbatim}
first element of $\mathrm{E}\left(\mathbf{g}_{t} \mid \mathbf{g}_{t-1}, \mathbf{g}_{t-2}, \ldots\right)$
$=\mathrm{E}\left(\varepsilon_{t} \mid \mathbf{g}_{t-1}, \mathbf{g}_{t-2}, \ldots\right)$
$=\mathrm{E}\left(\varepsilon_{t} \mid \varepsilon_{t-1}, \varepsilon_{t-2}, \ldots, y_{t-2} \varepsilon_{t-1}, y_{t-3} \varepsilon_{t-2}, \ldots\right)$
$=0$ (since $\left\{\varepsilon_{t}\right\}$ is i.i.d. and $y_{t-j}$ is a function of $\left(\varepsilon_{t-j}, \varepsilon_{t-j-1}, \ldots\right)$ )
\end{verbatim}

and

$$
\begin{aligned}
& \text { second element of } \mathrm{E}\left(\mathbf{g}_{t} \mid \mathbf{g}_{t-1}, \mathbf{g}_{t-2}, \ldots\right) \\
& =\mathrm{E}\left(\varepsilon_{t} y_{t-1} \mid \mathbf{g}_{t-1}, \mathbf{g}_{t-2}, \ldots\right) \\
& =\mathrm{E}\left[\mathrm{E}\left(\varepsilon_{t} y_{t-1} \mid y_{t-1}, \mathbf{g}_{t-1}, \mathbf{g}_{t-2}, \ldots\right) \mid \mathbf{g}_{t-1}, \mathbf{g}_{t-2}, \ldots\right]
\end{aligned}
$$

(by the Law of Iterated Expectations)


\begin{align*}
& =\mathrm{E}\left[y_{t-1} \mathrm{E}\left(\varepsilon_{t} \mid y_{t-1}, \mathbf{g}_{t-1}, \mathbf{g}_{t-2}, \ldots\right) \mid \mathbf{g}_{t-1}, \mathbf{g}_{t-2}, \ldots\right] \\
& =0 \quad\left(\text { since } \mathrm{E}\left(\varepsilon_{t} \mid y_{t-1}, \varepsilon_{t-1}, \varepsilon_{t-2}, \ldots, y_{t-2} \varepsilon_{t-1}, y_{t-3} \varepsilon_{t-2}, \ldots\right)=0\right) \tag{6.4.3}
\end{align*}


In particular, $\mathrm{E}\left(\varepsilon_{t} y_{t-1}\right)=0$, so $\mathbf{x}_{t}$ is orthogonal to the error term (Assumption 2.3). Finally, the rank condition about the moment matrix $\mathrm{E}\left(\mathbf{x}_{t} \mathbf{x}_{t}^{\prime}\right)$ (Assumption 2.4) is satisfied because the determinant of

\[
\mathrm{E}\left(\mathbf{x}_{t} \mathbf{x}_{t}^{\prime}\right)=\left[\begin{array}{cc}
1 & \mu  \tag{6.4.4}\\
\mu & \gamma_{0}+\mu^{2}
\end{array}\right]
\]

is $\gamma_{0}>0$. This also means that $\mathrm{E}\left(\mathbf{g}_{t} \mathbf{g}_{t}^{\prime}\right)$ is nonsingular because under conditional homoskedasticity $\mathrm{E}\left(\mathbf{g}_{t} \mathbf{g}_{t}^{\prime}\right)=\sigma^{2} \mathrm{E}\left(\mathbf{x}_{t} \mathbf{x}_{t}^{\prime}\right)$. So the nonsingularity condition in Assumption 2.5 is met.

Thus the AR(1) with an independent white noise process satisfies all the assumptions of Proposition 2.5, so parts (a)-(c) of the propositions hold true here, with $(c, \phi)^{\prime}$ as $\boldsymbol{\beta},\left(1, y_{t-1}\right)^{\prime}$ as $\mathbf{x}_{t}$, and $\sigma^{2}$ as the error variance. In particular, the OLS estimate of the error variance,

$$
s^{2}=\frac{1}{n-2} \sum_{t=1}^{n} e_{t}^{2}, \quad e_{t} \equiv y_{t}-\hat{c}-\hat{\phi} y_{t-1} \text { where }(\hat{c}, \hat{\phi}) \text { are OLS estimates, }
$$

is consistent for $\sigma^{2}$.

\section*{Estimation of AR(p)}
Similarly for $\operatorname{AR}(p)$, the autoregressive coefficients $\left(\phi_{1}, \ldots, \phi_{p}\right)$ can be consistently estimated by regressing $y_{t}$ on the constant and lagged $y$ 's, ( $y_{t-1}, \ldots, y_{t-p}$ ). We assume the sample is $\left(y_{-p+1}, y_{-p+2}, \ldots, y_{0}, y_{1}, \ldots, y_{n}\right)$, inclusive of $p$ lagged values prior to $y_{1}$, so that the $\operatorname{AR}(p)$ equation for $t=1$ can be included and the sample period is from $t=1$ to $n$.

Proposition 6.7 (Estimation of AR coefficients): Let $\left\{y_{t}\right\}$ be the AR( $p$ ) process following (6.2.6) with the stationarity condition (6.1.18). Suppose further that $\left\{\varepsilon_{t}\right\}$ is independent white noise. Then the OLS estimator of ( $c, \phi_{1}, \phi_{2}, \ldots, \phi_{p}$ )\\
obtained by regressing $y_{t}$ on the constant and $p$ lagged values of $y$ for the sample period of $t=1,2, \ldots, n$ is consistent and asymptotically normal. Letting $\boldsymbol{\beta}=\left(c, \phi_{1}, \phi_{2}, \ldots, \phi_{p}\right)^{\prime}$ and $\mathbf{x}_{t}=\left(1, y_{t-1}, \ldots, y_{t-p}\right)^{\prime}$, and $\widehat{\boldsymbol{\beta}}=$ OLS estimate of $\beta$, we have

$$
\operatorname{Avar}(\widehat{\boldsymbol{\beta}})=\sigma^{2} \mathrm{E}\left(\mathbf{x}_{t} \mathbf{x}_{t}^{\prime}\right)^{-1}
$$

which is consistently estimated by

$$
\widehat{\operatorname{Avar}(\widehat{\boldsymbol{\beta}})}=s^{2} \cdot\left(\frac{1}{n} \sum_{t=1}^{n} \mathbf{x}_{t} \mathbf{x}_{t}^{\prime}\right)^{-1}
$$

where $s^{2}$ is the OLS estimate of the error variance given by


\begin{equation*}
s^{2}=\frac{1}{n-p-1} \sum_{t=1}^{n}\left(y_{t}-\hat{c}-\hat{\phi}_{1} y_{t-1}-\cdots-\hat{\phi}_{p} y_{t-p}\right)^{2} \tag{6.4.5}
\end{equation*}


The proof is to verify the assumptions of Proposition 2.5, which can be done in the same way as we did for $\operatorname{AR}(1)$. The only difficult part is to show that $\mathrm{E}\left(\mathbf{x}_{t} \mathbf{x}_{t}^{\prime}\right)$ is nonsingular (Assumption 2.4), which is equivalent to requiring the $p \times p$ autocovariance matrix $\operatorname{Var}\left(y_{t}, \ldots, y_{t-p+1}\right)$ to be nonsingular (see Review Question 1(b) below). It can be shown that the autocovariance matrix of a covariance-stationary process is nonsingular for any $p$ if $\gamma_{0}>0$ and $\gamma_{j} \rightarrow 0$ as $j \rightarrow \infty$. ${ }^{6}$ So Assumption 2.4, too, is met.

We have not covered the maximum likelihood estimation of $\operatorname{AR}(p)$, but for those of you who are curious about the connection to ML, the OLS estimator is numerically the same as the conditional Gaussian ML estimator and asymptotically equivalent to the exact Gaussian ML estimator. See Section 8.7.

\section*{Choice of Lag Length}
All these nice results about estimation of autoregressions assume that the order (the lag length) of autoregression, $p$, is known. How should we proceed if the order $p$ is unknown? First of all, it is clear that Proposition 6.7, which is about $p$-th order autoregressions, is applicable to autoregressions of lower order, say $r<p$. That is, even if $\phi_{r} \neq 0$ but $\phi_{r+1}=\phi_{r+2}=\cdots=\phi_{p}=0$, the proposition remains applicable as long as ( $\phi_{1}, \phi_{2}, \ldots, \phi_{r}$ ) satisfies the stationarity condition. In particular, the OLS estimates of ( $\phi_{r+1}, \phi_{r+2}, \ldots, \phi_{p}$ ) will converge to the true value of zero. To put the same argument differently, suppose that the true order of

\footnotetext{${ }^{6}$ See, e.g., Proposition 5.1.1 of Brockwell and Davis (1991).
}
the autoregression is $p$ (so $\phi_{p} \neq 0$ ) and suppose that all we know about $p$ is that it is less than or equal to some known integer $p_{\max }$. By setting the $r$ in the argument just given to $p$ and the $p$ to $p_{\text {max }}$, we see that Proposition 6.7 is applicable to autoregressions with $p_{\text {max }}$ lags: if the $\left\{y_{t}\right\}$ is stationary $\operatorname{AR}(p)$ and if we regress $y_{t}$ on the constant and $p_{\text {max }}$ lags of $y$, then the OLS coefficient estimates will be consistent and asymptotically normal.

It is natural to think that the true lag length may be estimated from this OLS estimate of $\left(\phi_{1}, \ldots, \phi_{p_{\max }}\right)$. We consider two classes of rules for determining the lag length. The first is the

General-to-specific sequential $t$ rule: Start with an autoregression of $p_{\text {max }}$ lags.\\
If the last lag is significant at some prespecified significance level (say, 10 percent), then end the procedure and set the lag length to $p_{\max }$. Otherwise drop the last lag, reestimate the autoregression with one fewer lag, and repeat the same test. If this process continues until there is only one lag in the autoregression and if that lag is insignificant, then set the lag length to 0 (no lags).

This rule has the following properties, to be illustrated in a moment. Since the $t$-test is consistent, the lag length thus chosen will never be less than $p$ (the true lag length) in large samples. However, the probability of overfitting (that the lag length chosen is greater than the true lag length $p$ ) is not zero, even in large samples. That is, if $\hat{p}$ is the lag length chosen by the sequential $t$ rule,


\begin{equation*}
\lim _{n \rightarrow \infty} \operatorname{Prob}(\hat{p}<p)=0 \quad \text { and } \quad \lim _{n \rightarrow \infty} \operatorname{Prob}(\hat{p}>p)>0 . \tag{6.4.6}
\end{equation*}


To illustrate, suppose $p=2$ and $p_{\text {max }}=3$. The sequential $t$ rule starts out with three lags in the autoregression:

$$
y_{t}=c+\phi_{1} y_{t-1}+\phi_{2} y_{t-2}+\phi_{3} y_{t-3}+\varepsilon_{t}
$$

We test the null hypothesis that $\phi_{3}=0$ at the 10 percent significance level. With a probability of 10 percent, we reject the (true) null and set the lag length to 3 , or accept the null with a probability of 90 percent, in large samples. In the event of accepting the null, we set $\phi_{3}=0$, estimate an autoregression with two lags, and test the hypothesis that $\phi_{2}=0$. Since $\phi_{2} \neq 0$ in truth (i.e., since $p=2$ ), the absolute value of the $t$-value on the OLS estimate of $\phi_{2}$ would be very large in large samples, so we would never accept the (false) null that the second lag is zero. Therefore, in this example, $\operatorname{Prob}(\hat{p}=2)=90 \%$ and $\operatorname{Prob}(\hat{p}=3)=10 \%$ in large samples.

There are two minor variations in the choice of the sample period with data $\left(y_{1}, \ldots, y_{n}\right)$. One is to use the same sample period of $t=p_{\text {max }}+1, p_{\text {max }}+2, \ldots, n$\\
throughout. The other is to let the sample period be "elastic" by expanding it as fewer and fewer lags are included. That is, when the autoregression being estimated has $j$ lags, the sample period is $t=j+1, j+2, \ldots, n$.

The second rule for selecting the lag length uses either the Akaike information criterion (AIC) or the Schwartz information criterion (SIC), also called the Bayesian information criterion (BIC). This procedure sets the lag length to the $j$ that minimizes


\begin{equation*}
\log \left(S S R_{j} / n\right)+(j+1) C(n) / n, \tag{6.4.7}
\end{equation*}


over $j=0,1,2, \ldots, p_{\text {max }}$, where $S S R_{j}$ is the sum of squared residuals for the autoregression with $j$ lags. In this formula, " $j+1$ " enters as the number of parameters (the coefficients of the constant and $j$ lags). For the AIC, $C(n)=2$ while for the BIC, $C(n)=\log (n)$. A few comments about this information-criterion-based rule:

\begin{itemize}
  \item As the number of parameters increases with $j$, the first term of the objective function (6.4.7) declines because the fit of the equation gets better, but the second term increases. The information criterion strikes a balance between a better fit and model parsimony.
  \item Without an upper bound $p_{\text {max }}$, a ridiculously large value of $j$ might be chosen (indeed, the information criterion achieves a minimum of $-\infty$ at $j+1=n$ because $S S R_{j}=0$ for $j=n-1$ ). In the present context where the true DGP is a finite-order autoregression, a reasonable upper bound is the maximum possible lag length.
  \item Again there can be two ways to set the sample period. You can use the fixed sample period of $t=p_{\text {max }}+1, p_{\text {max }}+2, \ldots, n$ throughout, in which case $S S R_{j}$ is the sum of $n-p_{\text {max }}$ squared residuals. Or you can use the "elastic" sample period, in which case $S S R_{j}$ is the sum of $n-j$ squared residuals. In either case, there is yet another variation, which is to replace the $n$ in the information criterion (6.4.7) by the actual sample size. So you replace the $n$ by $n-p_{\text {max }}$ if the fixed sample is to be used. Based on simulations and some theoretical considerations, Ng and Perron (2000) recommend the use of a fixed, rather than an "elastic," sample period, with $n-p_{\max }$ replacing $n$ in the information criterion.
\end{itemize}

Let $\hat{p}_{\mathrm{AIC}}$ and $\hat{p}_{\mathrm{BIC}}$ be the lag lengths chosen by the AIC and BIC, respectively. Like the lag length selected by the general-to-specific $t$ rule, they are functions of data, and hence are random variables. For them it is possible to show the following (see, e.g., Lütkepohl, 1993, Section 4.3):\\
(1) $\hat{p}_{\mathrm{BIC}} \leq \hat{p}_{\mathrm{AIC}}$ unless the sample size is very small. (If the fixed sample of $t=p_{\text {max }}+1, \ldots, n$ is used and " $n-p_{\text {max }}$ " replaces $n$ in (6.4.7), this is true for any finite sample with $n \geq p_{\text {max }}+8$.) This is an algebraic result and holds for any sample on $y$.\\
(2) Suppose that $\left\{y_{t}\right\}$ is stationary $\operatorname{AR}(p)$ and $\left\{\varepsilon_{t}\right\}$ is independent white noise with a finite fourth moment. Then


\begin{equation*}
\lim _{n \rightarrow \infty} \hat{p}_{\mathrm{BIC}}=p, \quad \lim _{n \rightarrow \infty} \operatorname{Prob}\left(\hat{p}_{\mathrm{AIC}}<p\right)=0, \quad \lim _{n \rightarrow \infty} \operatorname{Prob}\left(\hat{p}_{\mathrm{AIC}}>p\right)>0 . \tag{6.4.8}
\end{equation*}


That is, just like the general-to-specific sequential $t$ rule, the AIC rule has a positive probability of overfitting. In contrast, the BIC rule is consistent. Furthermore, Hannan and Deistler (1988, Theorem 5.4.1(c) and its Remark 1) have shown that the consistency of $\hat{p}_{\mathrm{BIC}}$ holds when the upper bound $p_{\text {max }}$ is made to increase at a rate of $\log (n)$ (i.e., when it is set to $[c \log (n)]$ [the integer part of $c \log (n)]$ for any given $c>0$ ). This means that you do not need to know an upper bound in the consistent estimation by BIC of the order of an autoregression.

\section*{Estimation of VARs}
Estimation of VAR coefficients is equally easy. If $\mathbf{y}_{t}$ is $M$-dimensional, the VAR( $p$ ) (6.3.9) is an $M$-equation system and can be written as


\begin{equation*}
y_{t m}=\mathbf{x}_{t}^{\prime} \delta_{m}+\varepsilon_{t m} \quad(m=1,2, \ldots, M) \tag{6.4.9}
\end{equation*}


where $y_{t m}$ is the $m$-th element of $\mathbf{y}_{t}, \varepsilon_{t m}$ is the $m$-th element of $\varepsilon_{t}$, and

$$
\begin{gathered}
\underset{(M p+1) \times \mathbf{I}}{\mathbf{x}_{t}}=\left[\begin{array}{c}
1 \\
\mathbf{y}_{t-1} \\
\mathbf{y}_{t-2} \\
\vdots \\
\mathbf{y}_{t-p}
\end{array}\right], \quad \boldsymbol{\delta}_{m}=\left[\begin{array}{c}
c_{m} \\
\boldsymbol{\phi}_{1 m} \\
\boldsymbol{\phi}_{2 m} \\
\vdots \\
\boldsymbol{\phi}_{p m}
\end{array}\right], \\
c_{m}=m \text {-th element of } \mathbf{c}, \quad \boldsymbol{\phi}_{j m}^{\prime}=m \text {-th row of } \mathbf{\Phi}_{j} .
\end{gathered}
$$

So the $M$ equations share the same set of regressors. It is easy to verify that $\mathbf{x}_{t}$ is orthogonal to the error term in each equation and that the error is conditionally homoskedastic (the proof is very similar to the proof for $\operatorname{AR}(p)$ ). So the $M$-equation system with i.i.d. $\left\{\varepsilon_{t}\right\}$ satisfies all the conditions for the multivariate regression model of Section 4.5. Thus, efficient GMM estimation is very easy:\\
just do equation-by-equation OLS. The expressions for the asymptotic variance and its consistent estimate can be obtained from Proposition 4.6 with $\mathbf{z}_{i m}=\mathbf{x}_{t}$ (so (4.5.13') on page 280 becomes $\left.\frac{1}{n} \sum_{t} \mathbf{x}_{t} \mathbf{x}_{t}^{\prime}\right)$. If $\boldsymbol{\delta}$ is the $(M p+1) M$-dimensional stacked vector created from ( $\boldsymbol{\delta}_{1}, \ldots, \boldsymbol{\delta}_{M}$ ) and if $\hat{\boldsymbol{\delta}}$ is the equation-by-equation OLS estimate of $\boldsymbol{\delta}$, we have from (4.5.17)


\begin{equation*}
\widehat{\operatorname{Avar}(\hat{\boldsymbol{\delta}})}=\widehat{\boldsymbol{\Omega}} \otimes\left(\frac{1}{n} \sum_{t=1}^{n} \mathbf{x}_{t} \mathbf{x}_{t}^{\prime}\right)^{-1} \tag{6.4.10}
\end{equation*}


where


\begin{equation*}
\widehat{\boldsymbol{\Omega}} \equiv \frac{1}{n-M p-1} \sum_{t=1}^{n} \hat{\varepsilon}_{t} \hat{\varepsilon}_{t}^{\prime}, \quad \hat{\varepsilon}_{t}=\mathbf{y}_{t}-\hat{\mathbf{c}}-\widehat{\boldsymbol{\Phi}}_{1} \mathbf{y}_{t-1}-\cdots-\widehat{\boldsymbol{\Phi}}_{p} \mathbf{y}_{t-p} \tag{6.4.11}
\end{equation*}


If we do not know the lag length $p$, we estimate it from data. In the information-criteria-based rule, the lag length we choose is the $k$ that minimizes


\begin{equation*}
\log \left(\left|\frac{1}{n} \sum_{t=1}^{n} \hat{\varepsilon}_{t} \hat{\varepsilon}_{t}^{\prime}\right|\right)+\left(k M^{2}+M\right) \cdot \frac{C(n)}{n}, \tag{6.4.12}
\end{equation*}


over $k=0,1, \ldots, p_{\text {max }}$. Here, as in the univariate case, $C(n)=2$ for the AIC and $\log (n)$ for the BIC. Exactly the same results, listed as (1) and (2) above for the univariate case, hold for the multivariate case (see, e.g, Lütkepohl, 1993, Section 4.3).

\section*{Estimation of ARMA $(p, q)$}
Letting

$$
\begin{gathered}
u_{t} \equiv \varepsilon_{t}+\theta_{1} \varepsilon_{t-1}+\theta_{2} \varepsilon_{t-2}+\cdots+\theta_{q} \varepsilon_{t-q} \\
\mathbf{z}_{t}=\left(1, y_{t-1}, y_{t-2}, \ldots, y_{t-p}\right)^{\prime}, \quad \delta=\left(c, \phi_{1}, \phi_{2}, \ldots, \phi_{p}\right)^{\prime}
\end{gathered}
$$

the univariate $\operatorname{ARMA}(p, q)$ equation (6.2.9) can be written as


\begin{equation*}
y_{t}=\mathbf{z}_{t}^{\prime} \delta+u_{t} . \tag{6.4.13}
\end{equation*}


The moving-average component of ARMA creates two problems. First, obviously, the error term $u_{t}$ is serially correlated (in fact, it is MA( $q$ )). Second, since the lagged dependent variables included in $\mathbf{z}_{t}$ are correlated with lags of $\varepsilon$ included in the error term, the regressors $\mathbf{z}_{t}$ are not orthogonal to the error term $u_{t}$. That is, serial correlation in the presence of lagged dependent variables results in biased estimation. The second problem could be dealt with by using suitably lagged\\
dependent variables ( $y_{t-q-1}, y_{t-q-2}, \ldots$ ) as instruments. Regarding the first problem, the method to be developed later in this chapter provides consistent estimates of the coefficient vector in the presence of serial correlation. Given a consistent estimate of $\boldsymbol{\delta}$, the MA parameters ( $\theta$ 's) can be estimated from the residuals. This procedure, however, is not efficient. ${ }^{7}$ We will not cover the topic of efficient estimation of ARMA( $p, q$ ) processes with Gaussian errors. Interested readers are referred to, e.g., Fuller (1996, Chapter 8) or Hamilton (1994, Chapter 5).

\section*{QUESTIONS FOR REVIEW}
\begin{enumerate}
  \item (Avar of AR coefficients)\\
(a) For AR(1), show that
\end{enumerate}

$$
\operatorname{Avar}(\hat{\phi})=\frac{\sigma^{2}}{\gamma_{0}}=1-\phi^{2},
$$

where $\hat{\phi}$ is the OLS estimator of $\phi$. Hint: It is $\sigma^{2}$ times the $(2,2)$ element of the inverse of (6.4.4).\\
(b) For $\operatorname{AR}(p)$, let $\hat{\boldsymbol{\phi}}$ be the OLS estimate of $\left(\phi_{1}, \phi_{2}, \ldots, \phi_{p}\right)^{\prime}$. Show that

$$
\operatorname{Avar}(\hat{\boldsymbol{\phi}})=\sigma^{2} \mathbf{V}^{-1}
$$

where $\mathbf{V}$ is the $p \times p$ autocovariance matrix

$$
\mathbf{V}=\operatorname{Var}\left(y_{t}, \ldots, y_{t-p+1}\right)=\left[\begin{array}{ccccc}
\gamma_{0} & \gamma_{1} & \gamma_{2} & \cdots & \gamma_{p-1} \\
\gamma_{1} & \gamma_{0} & \gamma_{1} & \cdots & \gamma_{p-2} \\
\vdots & \ddots & \ddots & \ddots & \vdots \\
\gamma_{p-2} & \cdots & \gamma_{1} & \gamma_{0} & \gamma_{1} \\
\gamma_{p-1} & \cdots & \gamma_{2} & \gamma_{1} & \gamma_{0}
\end{array}\right] .
$$

Hint: If $\mathbf{x}_{t}$ is as in Proposition 6.7, then

$$
\mathrm{E}\left(\mathbf{x}_{t} \mathbf{x}_{t}^{\prime}\right)=\left[\begin{array}{cc}
1 & \mu \mathbf{1}^{\prime} \\
\mu \mathbf{1} & \mathbf{V}+\mu^{2} \mathbf{1 \mathbf { 1 } ^ { \prime }}
\end{array}\right],
$$

where $\mathbf{1}$ is a $p \times 1$ vector of ones. Show that the inverse of this matrix is

\footnotetext{${ }^{7}$ The model yields infinitely many orthogonality conditions, $\mathrm{E}\left(u_{t} \cdot y_{s}\right)=0$ for all $s \leq t-q-1$. The GMM procedure can exploit only finitely many orthogonality conditions.
}
given by
$$
\mathbf{A} \equiv\left[\begin{array}{cc}
a & \mathbf{c}^{\prime} \\
\mathbf{c} & \mathbf{V}^{-1}
\end{array}\right] \quad \text { where } a \equiv 1+\mu^{2} \mathbf{1}^{\prime} \mathbf{V}^{-1} \mathbf{1}, \mathbf{c} \equiv-\mu \mathbf{V}^{-1} \mathbf{1} .
$$

This also shows that $\mathrm{E}\left(\mathbf{x}_{t} \mathbf{x}_{t}^{\prime}\right)$ is nonsingular if and only if $\mathbf{V}$ is. (If $\mathbf{V}$ is nonsingular, then $\mathrm{E}\left(\mathbf{x}_{t} \mathbf{x}_{t}^{\prime}\right)$ is nonsingular because it is invertible. If $\mathbf{V}$ is singular, then there exists a $p$-dimensional vector $\mathbf{d} \neq \mathbf{0}$ such that $\mathbf{V d}=\mathbf{0}$ and $\mathrm{E}\left(\mathbf{x}_{t} \mathbf{x}_{t}^{\prime}\right) \mathbf{f}=\mathbf{0}$ where $\mathbf{f}^{\prime}=\left(-\mu \mathbf{1}^{\prime} \mathbf{d}, \mathbf{d}^{\prime}\right)$.)\\
2. (Estimated VAR(1) coefficients) Consider a VAR(1) without the intercept: $\mathbf{y}_{t}=\mathbf{A y}_{t-1}+\boldsymbol{\varepsilon}_{t}$. Verify that the estimate of $\mathbf{A}$ by multivariate regression using the sample $\left(y_{1}, y_{2}, \ldots, y_{n}\right)$ (so the sample period is $t=2,3, \ldots, n$ ) is

$$
\widehat{\mathbf{A}}=\left(\sum_{t=2}^{n} \mathbf{y}_{t} \mathbf{y}_{t-1}^{\prime}\right)\left(\sum_{t=2}^{n} \mathbf{y}_{t-1} \mathbf{y}_{t-1}^{\prime}\right)^{-1} .
$$

Hint: The coefficient vector in the $m$-th equation is the $m$-th row of $\mathbf{A}$.

\subsection*{6.5 Asymptotics for Sample Means of Serially Correlated Processes}
This section studies the asymptotic properties of the sample mean

$$
\bar{y} \equiv \frac{1}{n} \sum_{t=1}^{n} y_{t}
$$

for serially correlated processes. (It should be kept in mind that $\bar{y}$ depends on $n$, although the notation does not make it explicit.) For the consistency of the sample mean, we already have a sufficient condition in the Ergodic Theorem of Chapter 2. The first proposition of this section provides another sufficient condition in the form of restrictions on covariance-stationary processes. We also provide two CLTs for serially correlated processes, one for linear processes and the other for ergodic stationary processes. Chapter 2 includes a CLT for ergodic stationary processes, but it rules out serial correlation because the process is assumed to be a martingale difference sequence. The CLT for ergodic stationary processes of this section generalizes it by allowing for serial correlation.

\section*{LLN for Covariance-Stationary Processes}
Proposition 6.8 (LLN for covariance-stationary processes with vanishing autocovariances): Let $\left\{y_{t}\right\}$ be covariance-stationary with mean $\mu$ and $\left\{\gamma_{j}\right\}$ be the autocovariances of $\left\{y_{t}\right\}$. Then\\
(a) $\bar{y} \rightarrow_{\text {m.s. }} \mu$ as $n \rightarrow \infty$ if $\lim _{j \rightarrow \infty} \gamma_{j}=0$.\\
(b)


\begin{equation*}
\lim _{n \rightarrow \infty} \operatorname{Var}(\sqrt{n} \bar{y})=\sum_{j=-\infty}^{\infty} \gamma_{j}<\infty \quad \text { if }\left\{\gamma_{j}\right\} \text { is summable. }{ }^{8} \tag{6.5.1}
\end{equation*}


Since a mean square convergence implies a convergence in probability, part (a) shows that a very mild condition on covariance-stationary processes suffices for the sample mean to be consistent for $\mu$. To prove part (a), by Chebychev's LLN (see Section 2.1), it suffices to show that $\lim _{n \rightarrow \infty} \operatorname{Var}(\bar{y})=0$, which is fairly easy (Analytical Exercise 9) once the expression for $\operatorname{Var}(\bar{y})$ is derived. Unlike in the case without serial correlation in $\left\{y_{t}\right\}$, calculating the variance of $\bar{y}$ is more cumbersome because covariances between two different terms have to be taken into account:


\begin{align*}
n \cdot \operatorname{Var}(\bar{y})= & \operatorname{Var}(\sqrt{n} \bar{y}) \\
= & \frac{1}{n}\left[\operatorname{Cov}\left(y_{1}, y_{1}+\cdots+y_{n}\right)+\cdots+\operatorname{Cov}\left(y_{n}, y_{1}+\cdots+y_{n}\right)\right] \\
= & \frac{1}{n}\left[\left(\gamma_{0}+\gamma_{1}+\cdots+\gamma_{n-2}+\gamma_{n-1}\right)+\left(\gamma_{1}+\gamma_{0}+\gamma_{1}+\cdots+\gamma_{n-2}\right)\right. \\
& \left.\quad+\cdots+\left(\gamma_{n-1}+\gamma_{n-2}+\cdots+\gamma_{1}+\gamma_{0}\right)\right] \\
= & \frac{1}{n}\left[n \gamma_{0}+2(n-1) \gamma_{1}+\cdots+2(n-j) \gamma_{j}+\cdots+2 \gamma_{n-1}\right] \\
= & \gamma_{0}+2 \sum_{j=1}^{n-1}\left(1-\frac{j}{n}\right) \gamma_{j} \tag{6.5.2}
\end{align*}


(If you have difficulty deriving this, set $n$ to, say, 3 to verify.) For each fixed $j \geq 1$, the coefficient of $\gamma_{j}$ in (6.5.2) goes to 1 as $n \rightarrow \infty$, so that all the autocovariances eventually enter the sum with a coefficient of one. This is not enough to show the desired result (6.5.1). Filling in the rest of the proof of (b) is Analytical Exercise 10.

For a covariance-stationary process $\left\{y_{t}\right\}$, we define the long-run variance to be the limit as $n \rightarrow \infty$ of $\operatorname{Var}(\sqrt{n} \bar{y})$ (if it exists). So part (b) of the proposition says that the long-run variance equals $\sum_{j=-\infty}^{\infty} \gamma_{j}$, which in turn equals the value of

\footnotetext{${ }^{8}$ Anderson (1971, Lemma 8.3.1, p. 460). Absolute summability of $\left\{\gamma_{j}\right\}$ is not needed.
}
the autocovariance-generating function $g_{Y}(z)$ at $z=1$ (see (6.2.15)) or $2 \pi$ times the spectrum at frequency zero (see (6.2.17)).

\section*{Two Central Limit Theorems}
We present two CLTs to cover serial correlation. The first is a generalization of Lindeberg-Levy.

Proposition 6.9 (CLT for MA( $\infty$ )): Let $y_{t}=\mu+\sum_{j=0}^{\infty} \psi_{j} \varepsilon_{t-j}$ where $\left\{\varepsilon_{t}\right\}$ is independent white noise and $\sum_{j=0}^{n}\left|\psi_{j}\right|<\infty$. Then


\begin{equation*}
\sqrt{n}(\bar{y}-\mu) \underset{\mathrm{d}}{\rightarrow} N\left(0, \sum_{j=-\infty}^{\infty} \gamma_{j}\right) . \tag{6.5.3}
\end{equation*}


We will not prove this result (see Anderson, 1971, Theorem 7.7.8 or Brockwell and Davis, 1991, Theorem 7.1.2 for a proof), but we should not be surprised to see that $\operatorname{Avar}(\bar{y})$ (the variance of the limiting distribution of $\sqrt{n}(\bar{y}-\mu)$ ) equals the long-run variance $\sum_{j=-\infty}^{\infty} \gamma_{j}$. We know from Lemma 2.1 that " $x_{n} \rightarrow_{\mathrm{d}} x$," $" \operatorname{Var}\left(x_{n}\right) \rightarrow c<\infty " \Rightarrow " \operatorname{Var}(x)=c . "$ Here, set $x_{n}=\sqrt{n}(\bar{y}-\mu)$. By (6.5.1), $\operatorname{Var}\left(x_{n}\right)$ converges to a finite limit $\sum \gamma_{j}$. So the variance of the limiting distribution of $\sqrt{n}(\bar{y}-\mu)$ is $\sum \gamma_{j}$.

By Proposition 6.1(d), the process in Proposition 6.9 is (stationary and) ergodic. Remember that ergodicity limits serial correlation by requiring that two random variables sufficiently far apart in the sequence be almost independent. But ergodic stationarity alone is not sufficient to ensure the asymptotic normality of $\bar{y}$. A natural question arises: is there a further restriction on ergodicity, short of requiring the process to be linear as in Proposition 6.9, that delivers asymptotic normality? The restriction that has become increasingly popular is what this book calls Gordin's condition for ergodic stationary processes. It has three parts. The first two parts follow:\\
(a) $\mathrm{E}\left(y_{t}^{2}\right)<\infty$. (This is a restriction on [strictly] stationary processes because, strictly speaking, a stationary process might not have finite second moments.)\\
(b) $\mathrm{E}\left(y_{t} \mid y_{t-j}, y_{t-j-1}, \ldots\right) \rightarrow_{\text {m.s. }} 0$ as $j \rightarrow \infty$. (Since $\left\{y_{t}\right\}$ is stationary, this condition is equivalent to $\mathrm{E}\left(y_{0} \mid y_{-j}, y_{-j-1}, \ldots\right) \rightarrow_{\text {m.s. }} 0$ as $j \rightarrow \infty$. Because $\mathrm{E}\left(y_{t} \mid y_{t-j}, y_{t-j-1}, \ldots\right)$ is a random variable, the convergence cannot be the usual convergence for real numbers.)

This condition implies that the unconditional mean is zero, $\mathrm{E}\left(y_{t}\right)=0 .{ }^{9}$ As the forecast (the conditional expectation) is based on less and less information, it should approach the forecast based on no information, i.e., the unconditional expectation. To prepare for the third part, let $I_{t} \equiv\left(y_{t}, y_{t-1}, y_{t-2}, \ldots\right)$ and write $y_{t}$ as

$$
\begin{aligned}
y_{t}= & y_{t}-\left[\mathrm{E}\left(y_{t} \mid I_{t-1}\right)-\mathrm{E}\left(y_{t} \mid I_{t-1}\right)\right]-\left[\mathrm{E}\left(y_{t} \mid I_{t-2}\right)-\mathrm{E}\left(y_{t} \mid I_{t-2}\right)\right]-\cdots \\
& -\left[\mathrm{E}\left(y_{t} \mid I_{t-j}\right)-\mathrm{E}\left(y_{t} \mid I_{t-j}\right)\right] \\
= & {\left[y_{t}-\mathrm{E}\left(y_{t} \mid I_{t-1}\right)\right]+\left[\mathrm{E}\left(y_{t} \mid I_{t-1}\right)-\mathrm{E}\left(y_{t} \mid I_{t-2}\right)\right]+\cdots } \\
& +\left[\mathrm{E}\left(y_{t} \mid I_{t-j+1}\right)-\mathrm{E}\left(y_{t} \mid I_{t-j}\right)\right]+\mathrm{E}\left(y_{t} \mid I_{t-j}\right) \\
= & \left(r_{t 0}+r_{t 1}+\cdots+r_{t, j-1}\right)+\mathrm{E}\left(y_{t} \mid y_{t-j}, y_{t-j-1}, \ldots\right)
\end{aligned}
$$

where

$$
r_{t k} \equiv \mathrm{E}\left(y_{t} \mid I_{t-k}\right)-\mathrm{E}\left(y_{t} \mid I_{t-k-1}\right)
$$

This $r_{t k}$ is the revision of expectations about $y_{t}$ as the information set increases from $I_{t-k-1}$ to $I_{t-k}$. By (b), $y_{t}-\left(r_{t 0}+r_{t 1}+\cdots+r_{t, j-1}\right)$ converges in mean square to zero as $j \rightarrow \infty$ for each $t$. So $y_{i}$ can be written as

$$
y_{t}=\sum_{j=0}^{\infty} r_{t j}
$$

This sum is called the telescoping sum. The third part of Gordin's condition is\\
(c) $\sum_{j=0}^{\infty}\left[\mathrm{E}\left(r_{t j}^{2}\right)\right]^{1 / 2}<\infty$.

The telescoping sum indicates how the "shocks" represented by ( $r_{t 0}, r_{t 1}, \ldots$ ) influence the current value of $y$. Condition (c) says, roughly speaking, that shocks that occurred a long time ago do not have disproportionately large influence. As such, the condition restricts the extent of serial correlation in $\left\{y_{t}\right\}$.

To understand the condition, take for example a zero-mean AR(1) with independent errors:

$$
y_{t}=\phi y_{t-1}+\varepsilon_{t}, \quad|\phi|<1, \quad\left\{\varepsilon_{t}\right\} \text { independent white noise with } \sigma^{2}=\operatorname{Var}\left(\varepsilon_{t}\right) .
$$

Evidently, Gordin's condition (a) is satisfied. Condition (b) too is satisfied because

\footnotetext{${ }^{9}$ See Lemma 5.14 of White (1984) for a proof.
}

\begin{equation*}
\mathrm{E}\left(y_{t} \mid y_{t-j}, y_{t-j-1}, \ldots\right)=\phi^{j} y_{t-j} \tag{6.5.4}
\end{equation*}

and $\mathrm{E}\left[\left(\phi^{j} y_{t-j}\right)^{2}\right] \rightarrow 0$ as $j \rightarrow \infty$. The expectation revision $r_{t j}$ can be written as

\begin{equation*}
r_{t j}=\phi^{j} y_{t-j}-\phi^{j+1} y_{t-j-1}=\phi^{j}\left(y_{t-j}-\phi y_{t-j-1}\right)=\phi^{j} \varepsilon_{t-j} . \tag{6.5.5}
\end{equation*}


So the telescoping sum for $\operatorname{AR}(1)$ is the MA( $\infty$ ) representation. Condition (c) is satisfied because $\left[\mathrm{E}\left(r_{t j}^{2}\right)\right]^{1 / 2}=|\phi|^{j} \sigma$ and

$$
\sum_{j=0}^{\infty}|\phi|^{j} \sigma=\frac{\sigma}{1-|\phi|}<\infty
$$

The CLT we have been seeking is

\section*{Proposition 6.10 (Gordin's CLT for zero-mean ergodic stationary processes): ${ }^{10}$}
Suppose $\left\{y_{t}\right\}$ is stationary and ergodic and suppose Gordin's condition is satisfied. Then $\mathrm{E}\left(y_{t}\right)=0$, the autocovariances $\left\{\gamma_{j}\right\}$ are absolutely summable, and

$$
\sqrt{n} \bar{y} \underset{\mathrm{~d}}{\rightarrow} N\left(0, \sum_{j=-\infty}^{\infty} \gamma_{j}\right) .
$$

Since Gordin's condition is satisfied if $\left\{y_{t}\right\}$ is a martingale difference sequence, which exhibits no serial correlation, Proposition 6.10 is a generalization of the ergodic stationary Martingale Difference Sequence CLT of Chapter 2.

\section*{Multivariate Extension}
Extension of the above results to the multivariate case is straightforward. Let

$$
\overline{\mathbf{y}} \equiv \frac{1}{n} \sum_{t=1}^{n} \mathbf{y}_{t}
$$

be the sample mean of a vector process $\left\{\mathbf{y}_{t}\right\}$.

\begin{itemize}
  \item The multivariate version of Proposition 6.8 is that\\
(a) $\overline{\mathbf{y}} \rightarrow_{\text {m.s. }} \boldsymbol{\mu}$ if each diagonal element of $\boldsymbol{\Gamma}_{j}$ goes to zero as $j \rightarrow \infty$, and\\
(b) $\lim _{n \rightarrow \infty} \operatorname{Var}(\sqrt{n} \overline{\mathbf{y}})$ (which is the long-run covariance matrix of $\left\{\mathbf{y}_{t}\right\}$ ) equals $\sum_{j=-\infty}^{\infty} \boldsymbol{\Gamma}_{j}$ if $\left\{\boldsymbol{\Gamma}_{j}\right\}$ is summable (i.e., if each component of $\boldsymbol{\Gamma}_{j}$ is summable).
\end{itemize}

\footnotetext{${ }^{10}$ This result is due to Gordin (1969), restated as Theorem 5.15 of White (1984). The claim about the absolute summability of autocovariances is noted in footnote 19 of Hansen (1982).
}\begin{itemize}
  \item Therefore, the long-run covariance matrix of $\left\{\mathbf{y}_{t}\right\}$ can be written as
\end{itemize}


\begin{equation*}
\mathbf{G}_{\mathbf{Y}}(1)=2 \pi \mathbf{s}_{Y}(0)=\sum_{j=-\infty}^{\infty} \boldsymbol{\Gamma}_{j}=\boldsymbol{\Gamma}_{0}+\sum_{j=1}^{\infty}\left(\boldsymbol{\Gamma}_{j}+\boldsymbol{\Gamma}_{j}^{\prime}\right), \tag{6.5.6}
\end{equation*}


where $\mathbf{G}_{\mathbf{Y}}(z)$ (the autocovariance-generating function) and $\mathbf{s}_{Y}(\omega)$ (the spectrum) for the vector process $\left\{\mathbf{y}_{t}\right\}$ are defined in Section 6.3.

\begin{itemize}
  \item The multivariate version of Proposition 6.9 is that
\end{itemize}


\begin{equation*}
\sqrt{n}(\overline{\mathbf{y}}-\boldsymbol{\mu}) \underset{\mathrm{d}}{\rightarrow} N\left(\mathbf{0}, \sum_{j=-\infty}^{\infty} \boldsymbol{\Gamma}_{j}\right) \tag{6.5.7}
\end{equation*}


if $\left\{\mathbf{y}_{t}\right\}$ is vector MA( $\infty$ ) with absolutely summable coefficients and $\left\{\boldsymbol{\varepsilon}_{t}\right\}$ is vector independent white noise.

\begin{itemize}
  \item The multivariate extension of Gordin's condition on ergodic stationary processes is obvious: ${ }^{11}$
\end{itemize}

\section*{Gordin's condition on ergodic stationary processes}
(a) $\mathrm{E}\left(\mathbf{y}_{t} \mathbf{y}_{t}^{\prime}\right)$ exists and is finite.\\
(b) $\mathrm{E}\left(\mathbf{y}_{t} \mid \mathbf{y}_{t-j}, \mathbf{y}_{t-j-1}, \ldots\right) \rightarrow_{\text {m.s. }} \mathbf{0}$ as $j \rightarrow \infty$.\\
(c) $\sum_{j=0}^{\infty}\left[\mathrm{E}\left(\mathbf{r}_{t j}^{\prime} \mathbf{r}_{t j}\right)\right]^{1 / 2}$ is finite, where

$$
\mathbf{r}_{t j} \equiv \mathrm{E}\left(\mathbf{y}_{t} \mid \mathbf{y}_{t-j}, \mathbf{y}_{t-j-1}, \ldots\right)-\mathrm{E}\left(\mathbf{y}_{t} \mid \mathbf{y}_{t-j-1}, \mathbf{y}_{t-j-2}, \ldots\right)
$$

\begin{itemize}
  \item The multivariate version of Proposition 6.10 is obvious: Suppose Gordin's condition holds for vector ergodic stationary process $\left\{\mathbf{y}_{t}\right\}$. Then $\mathrm{E}\left(\mathbf{y}_{t}\right)=\mathbf{0},\left\{\boldsymbol{\Gamma}_{j}\right\}$ is absolutely summable, and
\end{itemize}


\begin{equation*}
\sqrt{n} \overline{\mathbf{y}} \underset{\mathrm{~d}}{\rightarrow} N\left(\mathbf{0}, \sum_{j=-\infty}^{\infty} \boldsymbol{\Gamma}_{j}\right) . \tag{6.5.8}
\end{equation*}


\section*{QUESTIONS FOR REVIEW}
\begin{enumerate}
  \item Show that $\operatorname{Var}(\sqrt{n} \bar{y})=\mathbf{1}^{\prime} \operatorname{Var}\left(y_{t}, y_{t-1}, \ldots, y_{t-n+1}\right) \mathbf{1} / n$, where $\operatorname{Var}\left(y_{t}\right.$, $\left.y_{t-1}, \ldots, y_{t-n+1}\right)$ is the $n \times n$ autocovariance matrix of a covariance-stationary process $\left\{y_{t}\right\}$.
\end{enumerate}

\footnotetext{${ }^{11}$ Stated as Assumption 3.5 in Hansen (1982).
}
2. (Avar( $\bar{y})$ for MA $(\infty))$ Verify that $\operatorname{Avar}(\bar{y})$ for the MA( $\infty$ ) process in Proposition 6.9 equals
$$
\sigma^{2} \cdot\left(\sum_{j=0}^{\infty} \psi_{j}\right)^{2} .
$$

Hint: It should equal $g_{Y}(1)$ for MA( $\infty$ ).\\
3. (Asymptotics of $\bar{y}$ for $\operatorname{AR}(1)$ ) Consider the $\operatorname{AR}(1)$ process $y_{t}=c+\phi y_{t-1}+\varepsilon_{t}$ with $|\phi|<1$. Does $\bar{y} \rightarrow_{\mathrm{p}} \mu$ ? Add the assumption that $\left\{\varepsilon_{t}\right\}$ is an independent white noise process. Calculate $\operatorname{Avar}(\bar{y})$. [Answer: [ $(1+\phi) /(1-\phi)] \gamma_{0}$.]\\
4. (Long-run covariance matrix of filtered series) Let $\mathbf{H}(L)=\mathbf{H}_{0}+\mathbf{H}_{1} L$ and $\left\{\mathbf{x}_{t}\right\}$ be covariance-stationary with absolutely summable autocovariances. Verify that the long-run covariance matrix of $\mathbf{y}_{t}=\mathbf{H}(L) \mathbf{x}_{t}$ is given by

$$
\left(\mathbf{H}_{0}+\mathbf{H}_{1}\right) \mathbf{G}_{\mathbf{x}}(1)\left(\mathbf{H}_{0}^{\prime}+\mathbf{H}_{1}^{\prime}\right) .
$$

Hint: Use (6.3.13) and (6.5.6).\\
5. (Long-run variance of "unit-root MA" processes) What is the long-run variance of $y_{t}=\varepsilon_{t}-\varepsilon_{t-1}$, where $\varepsilon_{t}$ is white noise? [Answer: zero.]

\subsection*{6.6 Incorporating Serial Correlation in GMM}
\section*{The Model and Asymptotic Results}
We have now all the tools needed to introduce serial correlation in the singleequation GMM model of Chapter 3. Recall that $\mathbf{g}_{t}$ in Chapter 3 (with the subscript " $i$ " replaced by " $t$ ") is a $K$-dimensional vector defined as $\mathbf{x}_{t} \cdot \varepsilon_{t}$, the product of the $K$-dimensional vector of instruments $\mathbf{x}_{t}$ and the scalar error term $\varepsilon_{t}$. By Assumption 3.3 (orthogonality conditions), the mean of $\mathbf{g}_{t}$ is zero. We assumed in Assumption 3.5 that $\left\{\mathbf{g}_{t}\right\}$ was a martingale difference sequence. Since no serial correlation is allowed under this assumption, the matrix $\mathbf{S}$, defined to be the asymptotic variance of $\overline{\mathbf{g}}\left(\equiv \frac{1}{n} \sum_{t=1}^{n} \mathbf{g}_{t}\right)$, was simply the variance of $\mathbf{g}_{t}$. We can now allow for serial correlation in $\left\{\mathbf{g}_{t}\right\}$ by relaxing Assumption 3.5 as

Assumption 3.5' (Gordin's condition restricting the degree of serial correlation): $\left\{\mathbf{g}_{t}\right\}$ satisfies Gordin's condition. Its long-run covariance matrix is nonsingular.

Then by the multivariate version of Proposition 6.10, $\sqrt{n} \overline{\mathbf{g}}$ converges to a normal distribution. The variance of this limiting distribution (namely, $\operatorname{Avar}(\overline{\mathbf{g}})$ ), denoted S, equals the long-run covariance matrix,


\begin{equation*}
\mathbf{S}=\sum_{j=-\infty}^{\infty} \boldsymbol{\Gamma}_{j}=\boldsymbol{\Gamma}_{0}+\sum_{j=1}^{\infty}\left(\boldsymbol{\Gamma}_{j}+\boldsymbol{\Gamma}_{j}^{\prime}\right) \tag{6.6.1}
\end{equation*}


where $\boldsymbol{\Gamma}_{j}$ is the $j$-th order autocovariance matrix


\begin{equation*}
\boldsymbol{\Gamma}_{j}=\mathrm{E}\left(\mathbf{g}_{t} \mathbf{g}_{t-j}^{\prime}\right) \quad(j=0, \pm 1, \pm 2, \ldots) \tag{6.6.2}
\end{equation*}


(Since $\mathrm{E}\left(\mathbf{g}_{t}\right)=\mathbf{0}$ by the orthogonality conditions, there is no need to subtract the mean from $\mathbf{g}_{t}$.)

That is, $\mathbf{S}(\equiv \operatorname{Avar}(\overline{\mathbf{g}}))$ equals $\boldsymbol{\Gamma}_{0}$ under Assumption 3.5 and $\sum_{j=-\infty}^{\infty} \boldsymbol{\Gamma}_{j}$ under Assumption $3.5^{\prime}$. This is the only difference between the GMM model with and without serial correlation. Accordingly, all the results we developed in Sections 3.5-3.7 carry over, provided that $\mathbf{S}$ is now given by (6.6.1) and $\widehat{\mathbf{S}}$, some consistent estimation of $\mathbf{S}$, is redefined to take into account serial correlation. More specifically:

\begin{itemize}
  \item The GMM estimator $\hat{\delta}(\widehat{\mathbf{W}})$ remains consistent and asymptotically normal. Its asymptotic variance is consistently estimated by (3.5.2), provided that the $\widehat{\mathbf{S}}$ there is consistent for $\mathbf{S}$ defined in (6.6.1). This estimated asymptotic variance sometimes gets the adjective heteroskedasticity and autocorrelation consistent (HAC).
  \item The GMM estimator achieves the minimum variance when $\operatorname{plim}_{n \rightarrow \infty} \widehat{\mathbf{W}}=\mathbf{S}^{-1}$ where $\mathbf{S}$ is given by (6.6.1).
  \item The two-step procedure described in Section 3.5 still provides an efficient GMM estimator, provided that $\widehat{\mathbf{S}}$ calculated in Step 1 is consistent for $\mathbf{S}$.
  \item With $\widehat{\mathbf{S}}$ thus properly defined, the expressions for the $t, W, J, C$, and $L R$ statistics of Sections 3.5-3.7 remain valid; those statistics retain the same asymptotic distributions in the presence of serial correlation.
\end{itemize}

\section*{Estimating S When Autocovariances Vanish after Finite Lags}
All these results assume that you have a consistent estimate, $\widehat{\mathbf{S}}$, of the long-run variance matrix $\mathbf{S}$. Obtaining such an estimate takes some intellectual effort, particularly when autocovariances do not vanish after finite lags.

We start our quest with consistent estimation of individual autocovariances. The natural estimator is


\begin{equation*}
\widehat{\boldsymbol{\Gamma}}_{j}=\frac{1}{n} \sum_{t=j+1}^{n} \hat{\mathbf{g}}_{t} \hat{\mathbf{g}}_{t-j}^{\prime} \quad(j=0,1, \ldots, n-1) \tag{6.6.3}
\end{equation*}


where

$$
\hat{\mathbf{g}}_{t} \equiv \mathbf{x}_{t} \cdot \hat{\varepsilon}_{t}, \quad \hat{\varepsilon}_{t}=y_{t}-\mathbf{z}_{t}^{\prime} \hat{\boldsymbol{\delta}}, \quad \hat{\delta} \text { consistent for } \delta .
$$

If we had the true value $\mathbf{g}_{t}$ in place of the estimated value $\hat{\mathbf{g}}_{t}$ in this formula, then $\widehat{\Gamma}_{j}$ would be consistent for $\Gamma_{j}$ under Assumptions 3.1 and 3.2, the assumptions implying ergodic stationarity for $\left\{\mathbf{g}_{t}\right\}$. But we have to use the estimated series $\left\{\hat{\mathbf{g}}_{t}\right\}$ rather than the true series to calculate autocovariances. For this reason we need some suitable fourth-moment condition (which we will not bother to state) for the estimated autocovariances to be consistent. The proof of consistency is not given here because it is very similar to the proof of Proposition 3.4.

Given the estimated autocovariances, there are two cases to consider for the calculation of $\widehat{\mathbf{S}}$. If we know a priori that $\boldsymbol{\Gamma}_{j}=\mathbf{0}$ for $j>q$ where $q$ is known and finite, then clearly $\mathbf{S}$ can be consistently estimated by


\begin{equation*}
\widehat{\mathbf{S}}=\widehat{\boldsymbol{\Gamma}}_{0}+\sum_{j=1}^{q}\left(\widehat{\boldsymbol{\Gamma}}_{j}+\widehat{\boldsymbol{\Gamma}}_{j}^{\prime}\right)=\sum_{j=-q}^{q} \widehat{\boldsymbol{\Gamma}}_{j} \quad\left(\text { recall: } \widehat{\boldsymbol{\Gamma}}_{-j}=\widehat{\boldsymbol{\Gamma}}_{j}^{\prime}\right) \tag{6.6.4}
\end{equation*}


More difficult is the other case where we do not know $q$ (which may or may not be infinite). How can we estimate the long-run variance matrix, which involves possibly infinitely many parameters? There are two approaches.

\section*{Using Kernels to Estimate S}
Recently, several procedures have been proposed to estimate the long-run variance matrix S. ${ }^{12}$ A class of estimators, called the kernel-based (or "nonparametric") estimators, can be expressed as a weighted average of estimated autocovariances:


\begin{equation*}
\widehat{\mathbf{S}}=\sum_{j=-n+1}^{n-1} k\left(\frac{j}{q(n)}\right) \cdot \widehat{\Gamma}_{j} . \tag{6.6.5}
\end{equation*}


\footnotetext{${ }^{12}$ For the results to be stated in this and the next subsection to hold, we need to impose a set of technical conditions (in addition to Gordin's condition above) that further restrict the nature of serial correlation. For a statement of those conditions and a more detailed discussion of the subject treated in this and the next subsection, see Den Haan and Levin (1996b).
}Here, the function $k(\cdot)$, which gives weights for autocovariances, is called a kernel and $q(n)$ is called the bandwidth. The bandwidth can depend on the sample size $n$. The estimator (6.6.4) above is a special kernel-based estimator with $q(n)=q$ and

\[
k(x)= \begin{cases}1 & \text { for }|x| \leq 1,  \tag{6.6.6}\\ 0 & \text { for }|x|>1 .\end{cases}
\]

This kernel is called the truncated kernel. In the present case of unknown $q$, we could use the truncated kernel with a bandwidth $q(n)$ that increases with the sample size. As $q(n)$ increases to infinity, more and more (and ultimately all) autocovariances can be included in the calculation of $\widehat{\mathbf{S}}$. This truncated kernel-based $\widehat{\mathbf{S}}$, however, is not guaranteed to be positive semidefinite in finite samples. This is a problem because then the estimated asymptotic variance of the GMM estimator may not be positive semidefinite.

Newey and West (1987) noted that the kernel-based estimator can be made nonnegative definite in finite samples if the kernel is the Bartlett kernel:

\[
k(x)= \begin{cases}1-|x| & \text { for }|x| \leq 1,  \tag{6.6.7}\\ 0 & \text { for }|x|>1 .\end{cases}
\]

The Bartlett kernel-based estimator of $\mathbf{S}$ is called (in econometrics) the NeweyWest estimator. For example, for $q(n)=3$, the kernel-based estimator includes autocovariances up to two (not three) lags:


\begin{equation*}
\widehat{\mathbf{S}}=\widehat{\boldsymbol{\Gamma}}_{0}+\left(\frac{2}{3}\right)\left(\widehat{\boldsymbol{\Gamma}}_{1}+\widehat{\boldsymbol{\Gamma}}_{1}^{\prime}\right)+\left(\frac{1}{3}\right)\left(\widehat{\boldsymbol{\Gamma}}_{2}+\widehat{\boldsymbol{\Gamma}}_{2}^{\prime}\right) . \tag{6.6.8}
\end{equation*}


In his definitive treatment of kernel-based estimation of the long-run variance, Andrews (1991) has examined positive semidefinite kernels, namely, the class of kernels (including Bartlett's) that yield a positive semidefinite $\widehat{\mathbf{S}}$. He shows, under some suitable regularity conditions, how the speed at which $\widehat{\mathbf{S}}$ converges (in mean square) to $\mathbf{S}$ depends on the kernel and $q(n)$. The most rapid possible rate is $n^{2 / 5}$ (i.e., $n^{2 / 5}$. $(\widehat{\mathbf{S}}-\mathbf{S})$ remains stochastically bounded), ${ }^{13}$ which occurs when $q(n)$ grows at rate $n^{1 / 5}$ (i.e., $q(n) / n^{1 / 5}$ remains bounded) and the kernel is a member of a certain subset of positive semidefinite kernels. The Bartlett kernel is not in this subset. For Bartlett, the most rapid possible rate is $n^{1 / 3}$ which occurs when $q(n)$ grows at rate $n^{1 / 3}$. Therefore, at least on asymptotic grounds, positive semidefinite

\footnotetext{${ }^{13}$ A sequence of random variables $\left\{x_{n}\right\}$ is said to be stochastically bounded and written $x_{n}=O_{p}$ if for any $\varepsilon$ there exists an $M>0$ such that $\operatorname{Prob}\left(\left|x_{n}\right|>M\right)<\varepsilon$ for all $n$. If a sequence converges in distribution to a random variable, then the sequence is stochastically bounded.
}
kernels in this subset should be preferred to Bartlett. An example of such kernels is the quadratic spectral (QS) kernel:

\begin{equation*}
k(x)=\frac{25}{12 \pi^{2} x^{2}}\left(\frac{\sin (6 \pi x / 5)}{6 \pi x / 5}-\cos (6 \pi x / 5)\right) \tag{6.6.9}
\end{equation*}


Since $k(x) \neq 0$ for $|x|>1$ in the QS kernel, all the estimated autocovariances $\widehat{\Gamma}_{j} (j=0,1, \ldots, n-1)$ enter the calculation of $\widehat{\mathbf{S}}$ even if $q(n)<n-1$.

To be sure, the knowledge of the growth rate is not enough to determine the value of the bandwidth for any particular data of finite length. Andrews (1991) also considered a data-dependent formula for determining the bandwidth that not only grows at the optimal rate but also minimizes some appropriate criterion (the asymptotic mean square error). But even in this formula, some parameter values need to be provided by the user, so the formula does not really provide automatic data-based selection of the bandwidth.

\section*{VARHAC}
The other procedure for estimating $\mathbf{S}$, called the VAR Heteroskedasticity and Autocorrelation Consistent (VARHAC) estimator, was proposed by Den Haan and Levin (1996a). The idea is to fit a finite-order VAR to the $K$-dimensional series $\left\{\hat{\mathbf{g}}_{t}\right\}$ and then construct the long-run covariance matrix implied by the estimated VAR. This is not to say that $\mathbf{g}_{t}$ is assumed to be a finite-order VAR. In fact, in this procedure (and also in the kernel-based procedure), $\mathbf{g}_{t}$ does not even have to be a linear process, let alone a finite-order autoregression. The VARHAC procedure consists of two steps.

Step 1: Lag length selection for each VAR equation. In this step, for each VAR equation, you determine the lag length by the Bayesian information criterion (BIC). Let $\hat{g}_{k i}$ be the $k$-th element of the $K$-dimensional vector $\hat{\mathbf{g}}_{i}$. The $k$-th VAR equation can be written as


\begin{align*}
\hat{g}_{k t}= & \phi_{11}^{(k)} \hat{g}_{1, t-1}+\phi_{12}^{(k)} \hat{g}_{1, t-2}+\cdots+\phi_{1 p}^{(k)} \hat{g}_{1, t-p} \\
& +\phi_{21}^{(k)} \hat{g}_{2, t-1}+\phi_{22}^{(k)} \hat{g}_{2, t-2}+\cdots+\phi_{2 p}^{(k)} \hat{g}_{2, t-p} \\
& \cdots \\
& +\phi_{K 1}^{(k)} \hat{g}_{K, t-1}+\phi_{K 2}^{(k)} \hat{g}_{K, t-2}+\cdots+\phi_{K p}^{(k)} \hat{g}_{K, t-p}+e_{k t} \\
= & \phi_{1}^{(k)} \hat{\mathbf{g}}_{t-1}+\cdots+\phi_{p}^{(k)} \hat{\mathbf{g}}_{t-p}+e_{k t} \tag{6.6.10}
\end{align*}


where

$$
\underset{(1 \times K)}{\boldsymbol{\phi}_{j}^{(k)}}=\left(\phi_{1 j}^{(k)}, \ldots, \phi_{K j}^{(k)}\right) \quad(j=1,2, \ldots, p) .
$$

In picking the lag length by the BIC, the maximum lag length $p_{\text {max }}$ needs to be specified. In Section 6.4, when the BIC was introduced in the context of estimating the true order (lag length) of a finite-order autoregression, $p_{\text {max }}$ was set to some integer known to be greater than or equal to the true lag length. In the present case where $\mathbf{g}_{t}$ is not necessarily a finiteorder VAR, Den Haan and Levin (1996a, Theorem 2(c)) show that the VARHAC estimator is consistent for $\mathbf{S}$ when $p_{\text {max }}$ grows at rate $n^{1 / 3}$ and recommend setting $p_{\text {max }}$ to $\left[n^{1 / 3}\right]$ (the integer part of $n^{1 / 3}$ ). To emphasize its dependence on the sample size, we denote the maximum lag length by $p_{\text {max }}(n)$. Let $S S R_{p}^{(k)}$ be the sum of squared residuals (SSR) from the OLS estimation of (6.6.10) (the $k$-th equation in the VAR) for $t=p_{\text {max }}(n)+ 1, p_{\text {max }}(n)+2, \ldots, n$. ${ }^{14}$ For $p=0, S S R_{p}^{(k)}=\sum_{t=p_{\max }(n)+1}^{n}\left(\hat{g}_{k t}\right)^{2}$, the SSR from the regression of $\hat{g}_{k t}$ on no regressors. Since there are $p K$ coefficients in the $k$-th equation, the BIC criterion is


\begin{equation*}
\log \left(S S R_{p}^{(k)} / n\right)+p K \cdot \log (n) / n \tag{6.6.11}
\end{equation*}


(The $n$ in this expression could be replaced by the actual sample size $n- p_{\text {max }}$ with no asymptotic consequences.) Let $p(k)$ be the $p$ that minimizes the information criterion over $p=0,1, \ldots, p_{\max }(n)$ for the $k$-th VAR equation. Also let $\hat{\boldsymbol{\phi}}_{j}^{(k)}(j=1,2, \ldots, p(k))$ be the OLS estimate of the VAR coefficients $\boldsymbol{\phi}_{j}^{(k)}(j=1,2, \ldots, p(k))$ when (6.6.10) is estimated by OLS for $p=p(k)$.\\
Step 2: Calculating the implied long-run variance. Let

$$
P \equiv \max \{p(1), p(2), \ldots, p(K)\},
$$

the largest lag length in the $K$ equations. By construction, $P \leq p_{\max }(n)$. Step 1 produces an estimated VAR which can be written as


\begin{align*}
\underset{(K \times 1)}{\hat{\mathbf{g}}_{t}}=\underset{(K \times K)(K \times 1)}{\widehat{\mathbf{\Phi}}_{1}} \underset{(K \times K)(K \times 1)}{\hat{\mathbf{g}}_{t-1}}+\cdots & +\underset{(K \times 1)}{\widehat{\mathbf{\Phi}}_{P}}{\underset{(K \times P)}{\hat{\mathbf{g}}_{t-P}}}^{\hat{\mathbf{e}}_{t}} \\
& \left(t=p_{\text {max }}(n)+1, p_{\text {max }}(n)+2, \ldots, n\right) . \tag{6.6.12}
\end{align*}


\footnotetext{${ }^{14}$ So the sample period is fixed throughout the process of picking the lag length.
}Since the lag length differs across equations, some rows of $\widehat{\boldsymbol{\Phi}}_{p}$ ( $p= 1,2, \ldots, P$ ) may be zero. The rows of $\widehat{\boldsymbol{\Phi}}_{p}$ are related to $\hat{\boldsymbol{\phi}}_{j}^{(k)}(j=1,2, \ldots$, $p(k)$ ) obtained in Step 1 as

\[
k \text {-th row of } \underset{(K \times K)}{\widehat{\Phi}_{p}}= \begin{cases}\hat{\boldsymbol{\phi}}_{p}^{(k)} & \text { for } 1 \leq p \leq p(k)  \tag{6.6.13}\\ 0 & \text { for } p(k)<p \leq P\end{cases}
\]

Let $\widehat{\Omega}$ be the VAR residual variance matrix:


\begin{equation*}
\underset{(K \times K)}{\widehat{\Omega}} \equiv \frac{1}{n} \sum_{t=p_{\max }(n)+1}^{n} \hat{\mathbf{e}}_{t} \hat{\mathbf{e}}_{t}^{\prime} \tag{6.6.14}
\end{equation*}


Recall that the long-run variance matrix of a VAR is given by (6.3.16) for $z=1$. The estimator of $\mathbf{S}$ implied by the estimated VAR (6.6.12) is then


\begin{equation*}
\widehat{\mathbf{S}}=\left[\mathbf{I}_{K}-\sum_{p=1}^{P} \widehat{\boldsymbol{\Phi}}_{p}\right]^{-1} \widehat{\mathbf{\Omega}}\left[\mathbf{I}_{K}-\sum_{p=1}^{P} \widehat{\boldsymbol{\Phi}}_{p}^{\prime}\right]^{-1} \tag{6.6.15}
\end{equation*}


This VARHAC estimator is positive semidefinite by construction. Den Haan and Levin (1996a) show that (under suitable regularity conditions) the estimator converges to $\mathbf{S}$ at a faster rate than any positive semidefinite kernel-based estimator for almost all autocovariance structures of $\mathbf{g}_{t}$ (which is not necessarily a finite-order VAR). In particular, if $\mathbf{g}_{t}$ is finite-order VARMA, then the lag order chosen by the BIC grows at a logarithmic rate and the VARHAC estimator converges at a rate arbitrarily close to $n^{1 / 2}$, which is faster than the maximum rate $n^{2 / 5}$ achieved by positive semidefinite kernel-based estimators.

\section*{QUESTIONS FOR REVIEW}
\begin{enumerate}
  \item When $\mathbf{x}_{t}=\mathbf{z}_{t}$, what is the efficient GMM estimator?
  \item (Consequence of ignoring serial correlation) Consider the efficient GMM estimator that presumes $\left\{\mathbf{g}_{t}\right\}$ to be serially uncorrelated. If $\left\{\mathbf{g}_{t}\right\}$ is in fact serially correlated, is the estimator consistent? Efficient?
\end{enumerate}

\subsection*{6.7 Estimation under Conditional Homoskedasticity (Optional)}
We saw in Section 3.8 that the GMM estimator reduced to the 2SLS estimator under conditional homoskedasticity. This section shows how the 2SLS can be generalized to incorporate serial correlation. It will then become clear that the GLS estimator is not a GMM estimator with serial correlation. The latter half of this section examines the relationship between the GLS and GMM.

\section*{Kernel-Based Estimation of $\mathbf{S}$ under Conditional Homoskedasticity}
The relationship between serial correlation in $\mathbf{g}_{t} \equiv \mathbf{x}_{t} \cdot \varepsilon_{t}$ and serial correlation in the error term $\varepsilon_{t}$ becomes clearer under conditional homoskedasticity. Let


\begin{equation*}
\omega_{j} \equiv \mathrm{E}\left(\varepsilon_{t} \varepsilon_{t-j}\right) \quad \text { for all } j \tag{6.7.1}
\end{equation*}


Since $\left\{\varepsilon_{t}\right\}$ is stationary by Assumptions 3.1 and 3.2, $\omega_{j}$ does not depend on $t$. If $\mathrm{E}\left(\varepsilon_{t}\right)=0$ (which is the case if $\mathbf{x}_{t}$ includes the constant), then $\omega_{j}$ equals the $j$-th order autocovariance of $\left\{\varepsilon_{t}\right\}$. Suppose the error term is conditionally homoskedastic:\\
conditional homoskedasticity:


\begin{equation*}
\mathrm{E}\left(\varepsilon_{t} \varepsilon_{t-j} \mid \mathbf{x}_{t}, \mathbf{x}_{t-j}\right)=\omega_{j} \quad(j=0, \pm 1, \pm 2, \ldots) \tag{6.7.2}
\end{equation*}


Under this additional condition, the usual argument utilizing the Law of Total Expectations can be applied to $\boldsymbol{\Gamma}_{j}$ :


\begin{align*}
\boldsymbol{\Gamma}_{j} & =\mathrm{E}\left(\mathbf{g}_{t} \mathbf{g}_{t-j}^{\prime}\right) \quad\left(\text { since } \mathrm{E}\left(\mathbf{g}_{t}\right)=\mathbf{0}\right) \\
& =\mathrm{E}\left(\varepsilon_{t} \varepsilon_{t-j} \mathbf{x}_{t} \mathbf{x}_{t-j}^{\prime}\right) \\
& =\mathrm{E}\left[\mathrm{E}\left(\varepsilon_{t} \varepsilon_{t-j} \mid \mathbf{x}_{t}, \mathbf{x}_{t-j}\right) \mathbf{x}_{t} \mathbf{x}_{t-j}^{\prime}\right] \quad \text { (by the Law of Total Expectations) } \\
& =\omega_{j} \mathrm{E}\left(\mathbf{x}_{t} \mathbf{x}_{t-j}^{\prime}\right) \tag{6.7.3}
\end{align*}


Thus, unless $\mathrm{E}\left(\mathbf{x}_{t} \mathbf{x}_{t-j}^{\prime}\right)$ is zero, $\left\{\mathbf{g}_{t}\right\}$ is serially correlated if and only if $\left\{\varepsilon_{t}\right\}$ is.\\
To estimate $\boldsymbol{\Gamma}_{j}$, we exploit the fact that $\boldsymbol{\Gamma}_{j}$ is a product of two second moments, $\mathrm{E}\left(\mathbf{x}_{t} \mathbf{x}_{t-j}^{\prime}\right)$ and $\mathrm{E}\left(\varepsilon_{t} \varepsilon_{t-j}\right)$. A natural estimator of $\mathrm{E}\left(\mathbf{x}_{t} \mathbf{x}_{t-j}^{\prime}\right)$ is $\frac{1}{n} \sum_{t=j+1}^{n} \mathbf{x}_{t} \mathbf{x}_{t-j}^{\prime}$. It is easy to show that $\omega_{j}$ is consistently estimated by $\frac{1}{n} \sum_{t=j+1}^{n} \hat{\varepsilon}_{t} \hat{\varepsilon}_{t-j}$ where $\hat{\varepsilon}_{t}$ is the residual from some consistent estimator (the proof is almost the same as in the case with no serial correlation, see Proposition 3.2). Thus, a natural estimate of $\boldsymbol{\Gamma}_{j}$\\
under conditional homoskedasticity is


\begin{equation*}
\widehat{\Gamma}_{j}=\left(\frac{1}{n} \sum_{t=j+1}^{n} \hat{\varepsilon}_{t} \hat{\varepsilon}_{t-j}\right)\left(\frac{1}{n} \sum_{t=j+1}^{n} \mathbf{x}_{t} \mathbf{x}_{t-j}^{\prime}\right) . \tag{6.7.4}
\end{equation*}


Given these estimated autocovariances, the kernel-based estimation of $\widehat{\mathbf{S}}$ is exactly the same as in the case without conditional homoskedasticity described in the previous section.

\section*{Data Matrix Representation of Estimated Long-Run Variance}
For the purpose of relating the GMM estimator with serial correlation but with conditional homoskedasticity to conventional estimators, the data matrix representation of $\widehat{\mathbf{S}}$ is useful. By defining an $n \times n$ matrix $\widehat{\boldsymbol{\Omega}}$ suitably, we can write $\widehat{\mathbf{S}}$ as


\begin{equation*}
\widehat{\mathbf{S}}=\frac{1}{n} \mathbf{X}^{\prime} \widehat{\mathbf{\Omega}} \mathbf{X}, \tag{6.7.5}
\end{equation*}


where $\mathbf{X}$ is the $n \times K$ data matrix whose $t$-th row is $\mathbf{x}_{t}^{\prime}$. The matrix $\widehat{\boldsymbol{\Omega}}$ looks much like the autocovariance matrix of $\left\{\varepsilon_{t}\right\}$ :

\[
\underset{(n \times n)}{\widehat{\boldsymbol{\Omega}}}=\left[\begin{array}{ccccc}
\hat{\omega}_{0} & \hat{\omega}_{1} & \hat{\omega}_{2} & \cdots & \hat{\omega}_{n-1}  \tag{6.7.6}\\
\hat{\omega}_{1} & \hat{\omega}_{0} & \hat{\omega}_{1} & \cdots & \hat{\omega}_{n-2} \\
\vdots & \ddots & \ddots & \ddots & \vdots \\
\hat{\omega}_{n-2} & \cdots & \hat{\omega}_{1} & \hat{\omega}_{0} & \hat{\omega}_{1} \\
\hat{\omega}_{n-1} & \cdots & \hat{\omega}_{2} & \hat{\omega}_{1} & \hat{\omega}_{0}
\end{array}\right]
\]

If we know a priori that $\operatorname{Cov}\left(\varepsilon_{t}, \varepsilon_{t-j}\right)=0$ (so that $\boldsymbol{\Gamma}_{j}=\mathbf{0}$ ) for $j>q$, specify the elements of this $\widehat{\boldsymbol{\Omega}}$ as

\[
\hat{\omega}_{j}= \begin{cases}\frac{1}{n} \sum_{t=j+1}^{n} \hat{\varepsilon}_{t} \hat{\varepsilon}_{t-j} & \text { if } 0 \leq j \leq q,  \tag{6.7.7}\\ 0 & \text { if } j>q .\end{cases}
\]

Then the $\widehat{\mathbf{S}}$ in (6.6.4) can be written as (6.7.5). If we do not know that there is such a $q$, then use the kernel-based estimator of $\mathbf{S}$. For the Bartlett kernel, let

\[
\hat{\omega}_{j}= \begin{cases}{\left[1-\frac{j}{q(n)}\right] \frac{1}{n} \sum_{t=j+1}^{n} \hat{\varepsilon}_{t} \hat{\varepsilon}_{t-j}} & \text { if } 0 \leq j \leq q(n)-1,  \tag{6.7.8}\\ 0 & \text { if } j>q(n)-1 .\end{cases}
\]

Then the consistent estimator (6.6.3) with (6.6.4) equals (6.7.5).

It is now easy to show the following extension of the 2SLS estimator to encompass serial correlation:

\begin{itemize}
  \item The efficient GMM estimator under conditional homoskedasticity can be written as
\end{itemize}


\begin{equation*}
\hat{\boldsymbol{\delta}}\left(\widehat{\mathbf{S}}^{-1}\right)=\left[\mathbf{Z}^{\prime} \mathbf{X}\left(\mathbf{X}^{\prime} \widehat{\boldsymbol{\Omega}} \mathbf{X}\right)^{-1} \mathbf{X}^{\prime} \mathbf{Z}\right]^{-1} \mathbf{Z}^{\prime} \mathbf{X}\left(\mathbf{X}^{\prime} \widehat{\mathbf{\Omega}} \mathbf{X}\right)^{-1} \mathbf{X}^{\prime} \mathbf{y}, \tag{6.7.9}
\end{equation*}


where $\mathbf{Z}$ and $\mathbf{y}$ are the data matrices of the regressors and the dependent variable. This generalizes (3.8.3') on page 230.

\begin{itemize}
  \item The consistent estimator of $\operatorname{Avar}\left(\hat{\delta}\left(\widehat{\mathbf{S}}^{-1}\right)\right)$ reduces to
\end{itemize}


\begin{align*}
\widehat{\operatorname{Avar}\left(\hat{\delta}\left(\widehat{\mathbf{S}}^{-1}\right)\right)} & =\left[\frac{1}{n} \mathbf{Z}^{\prime} \mathbf{X}\left(\frac{1}{n} \mathbf{X}^{\prime} \widehat{\mathbf{\Omega}} \mathbf{X}\right)^{-1} \frac{1}{n} \mathbf{X}^{\prime} \mathbf{Z}\right]^{-1} \\
& =n \cdot\left[\mathbf{Z}^{\prime} \mathbf{X}\left(\mathbf{X}^{\prime} \widehat{\mathbf{\Omega}} \mathbf{X}\right)^{-1} \mathbf{X}^{\prime} \mathbf{Z}\right]^{-1} . \tag{6.7.10}
\end{align*}


So the square roots of the diagonal elements of $\left[\mathbf{Z}^{\prime} \mathbf{X}\left(\mathbf{X}^{\prime} \widehat{\mathbf{\Omega}} \mathbf{X}\right)^{-1} \mathbf{X}^{\prime} \mathbf{Z}\right]^{-1}$ are the standard errors. This generalizes (3.8.5') on page 230.

\section*{Relation to GLS}
The procedure for incorporating serial correlation into the GMM estimation we have described is different from the GLS procedure. To see the difference most clearly, consider the case where $\mathbf{x}_{t}=\mathbf{z}_{t}$. If there are $L$ regressors, $\mathbf{X}^{\prime} \mathbf{Z}$ and $\mathbf{X}^{\prime} \widehat{\mathbf{\Omega}} \mathbf{X}$ in (6.7.9) are all $L \times L$ matrices, so the efficient GMM estimator that exploits the orthogonality conditions $\mathbf{E}\left(\mathbf{z}_{t} \cdot \varepsilon_{t}\right)=\mathbf{0}$ is the OLS estimator $\hat{\boldsymbol{\delta}}_{\mathrm{OLS}}=\left(\mathbf{Z}^{\prime} \mathbf{Z}\right)^{-1} \mathbf{Z}^{\prime} \mathbf{y}$. The consistent estimate of its asymptotic variance is obtained by setting $\mathbf{X}=\mathbf{Z}$ in (6.7.10):


\begin{equation*}
\widehat{\operatorname{Avar}\left(\hat{\boldsymbol{\delta}}_{\mathrm{OLS}}\right)}=n \cdot\left(\mathbf{Z}^{\prime} \mathbf{Z}\right)^{-1}\left(\mathbf{Z}^{\prime} \hat{\mathbf{\Omega}} \mathbf{Z}\right)\left(\mathbf{Z}^{\prime} \mathbf{Z}\right)^{-1} . \tag{6.7.11}
\end{equation*}


This may seem strange because we know from Section 1.6 that the GLS estimator is more efficient than OLS when the error is serially correlated. This is a finite sample result, but, given that we have a consistent estimator $\widehat{\boldsymbol{\Omega}}$ of $\boldsymbol{\Omega}$, we could take the GLS formula and estimate $\boldsymbol{\delta}$ as

$$
\hat{\boldsymbol{\delta}}_{\mathrm{GLS}}=\left(\mathbf{Z}^{\prime} \hat{\boldsymbol{\Omega}}^{-1} \mathbf{Z}\right)^{-1} \mathbf{Z}^{\prime} \hat{\boldsymbol{\Omega}}^{-1} \mathbf{y} .
$$

Is not this estimator asymptotically more efficient than OLS in that its asymptotic variance is smaller? The problem with GLS is that consistency is not guaranteed,\\
if the regressors are merely predetermined, as shown below. As we emphasized in Chapter 2, the regressors are not strictly exogenous in most time series models. It follows that GLS should not be used to correct for serial correlation in the error term for models lacking strict exogeneity. In particular, for the $\mathbf{x}_{t}=\mathbf{z}_{t}$ case, the correct procedure to adjust for serial correlation is to leave the point estimate unchanged while incorporating serial correlation in the estimate of the asymptotic variance.

That the GLS estimator is generally inconsistent can be seen as follows. To focus on the problem at hand, assume that the true autocovariance matrix is proportional to some known $n \times n$ matrix $\mathbf{V}$ :

$$
\operatorname{Var}\left(\varepsilon_{1}, \varepsilon_{2}, \ldots, \varepsilon_{n}\right)=\sigma^{2} \cdot \mathbf{V} .
$$

As shown in Section 1.6, the GLS estimator of the coefficient $\boldsymbol{\delta}$ in the regression model $\mathbf{y}=\mathbf{Z} \boldsymbol{\delta}+\boldsymbol{\varepsilon}$ can be calculated as the OLS estimator on the transformed model $\mathbf{C y}=\mathbf{C Z} \boldsymbol{\delta}+\mathbf{C} \boldsymbol{\varepsilon}$, where $\mathbf{C}$ is the "square root" of $\mathbf{V}^{-1}$ such that $\mathbf{C}^{\prime} \mathbf{C}=\mathbf{V}^{-1}$. But look at the $t$-th observation of the transformed model, which can be written as


\begin{equation*}
\tilde{y}_{t}=\tilde{\mathbf{z}}_{t}^{\prime} \boldsymbol{\delta}+\tilde{\varepsilon}_{t}, \tag{6.7.12}
\end{equation*}


where

$$
\begin{aligned}
& \tilde{y}_{t}=c_{t 1} y_{1}+\cdots+c_{t n} y_{n}, \\
& \tilde{\mathbf{z}}_{t}=c_{t 1} \mathbf{z}_{1}+\cdots+c_{t n} \mathbf{z}_{n}, \\
& \tilde{\varepsilon}_{t}=c_{t 1} \varepsilon_{1}+\cdots+c_{t n} \varepsilon_{n},
\end{aligned}
$$

$\left(c_{t 1}, \ldots, c_{t n}\right)$ is the $t$-th row of $\mathbf{C}$.

For the GLS estimator to be consistent, it is necessary that $\mathrm{E}\left(\tilde{\mathbf{z}}_{t} \cdot \tilde{\varepsilon}_{t}\right)=\mathbf{0}$. The left-hand side of this condition can be written as

$$
\begin{gathered}
\mathrm{E}\left(\tilde{\mathbf{z}}_{t} \cdot \tilde{\varepsilon}_{t}\right)=(\mathrm{a})+(\mathrm{b}), \\
(\mathrm{a})=\sum_{s=1}^{n} c_{t s}^{2} \mathrm{E}\left(\mathbf{z}_{s} \cdot \varepsilon_{s}\right), \quad(\mathrm{b})=\sum_{s \neq v} c_{t s} \cdot c_{t v} \mathrm{E}\left(\mathbf{z}_{s} \cdot \varepsilon_{v}\right) .
\end{gathered}
$$

By the orthogonality condition $\mathrm{E}\left(\mathbf{z}_{t} \cdot \varepsilon_{t}\right)=\mathbf{0}$, (a) is zero. For (b), because the regressors are merely predetermined and not assumed to be strictly exogenous, $\mathrm{E}\left(\mathbf{z}_{s} \cdot \varepsilon_{v}\right)$ is not guaranteed to be zero for $s \neq v$. So (b) is not zero, unless the $c$ 's take particular values to offset the nonzero terms.

There is one important special case where GLS is consistent, and that is when the error is a finite-order autoregressive process. To take the simplest case, suppose\\
that $\left\{\varepsilon_{t}\right\}$ is $\operatorname{AR}(1): \varepsilon_{t}=\phi \varepsilon_{t-1}+\eta_{t}$. Its autocovariance matrix is

$$
\operatorname{Var}\left(\varepsilon_{1}, \varepsilon_{2}, \ldots, \varepsilon_{n}\right)=\sigma^{2} . \mathbf{V}, \mathbf{V}=\frac{1}{1-\phi^{2}}\left[\begin{array}{ccccc}
1 & \phi & \phi^{2} & \cdots & \phi^{n-1} \\
\phi & 1 & \phi & \cdots & \phi^{n-2} \\
& & \ddots & & \\
\phi^{n-2} & \cdots & \phi & 1 & \phi \\
\phi^{n-1} & \cdots & \phi^{2} & \phi & 1
\end{array}\right] .
$$

It is easy to verify that the square root, $\mathbf{C}$, of $\mathbf{V}^{-1}$ (satisfying $\mathbf{C}^{\prime} \mathbf{C}=\mathbf{V}$ ) is given by

\[
\mathbf{C}=\left[\begin{array}{cccccc}
\sqrt{1-\phi^{2}} & 0 & 0 & \cdots & 0 & 0  \tag{6.7.13}\\
-\phi & 1 & 0 & \cdots & 0 & 0 \\
0 & -\phi & 1 & \cdots & 0 & 0 \\
\vdots & \vdots & \vdots & \cdots & \vdots & \vdots \\
0 & 0 & 0 & \cdots & -\phi & 1
\end{array}\right] .
\]

So the $t$-th transformed equation is


\begin{equation*}
y_{t}-\phi y_{t-1}=\left(\mathbf{z}_{t}-\phi \mathbf{z}_{t-1}\right)^{\prime} \boldsymbol{\delta}+\eta_{t} . \tag{6.7.14}
\end{equation*}


The GLS estimate of $\boldsymbol{\delta}$, which is the OLS estimate in the regression of $y_{t}-\phi y_{t-1}$ on $\mathbf{z}_{t}-\phi \mathbf{z}_{t-1}$ is consistent if $\eta_{t}$ is orthogonal to both $\mathbf{z}_{t}$ and $\mathbf{z}_{t-1}$. However, those orthogonality conditions are different from the orthogonality conditions for the untransformed model, which are that $\mathrm{E}\left(\mathbf{z}_{t} \cdot \varepsilon_{t}\right)=\mathbf{0}$.

\section*{QUESTIONS FOR REVIEW}
\begin{enumerate}
  \item Verify $(6.7 .5)$ for $n=3, q=1$.
  \item (Sargan's statistic) Derive the generalization of Sargan's statistic (3.8.10') on page 230 to the case of serially correlated errors.
  \item Suppose that $\mathbf{z}_{t}$ is a strict subset of $\mathbf{x}_{t}$. We have seen in Section 3.8 that the 2SLS estimator reduces to the OLS estimator. Is this true here as well? That is, does (6.7.9) reduce to $\left(\mathbf{Z}^{\prime} \mathbf{Z}\right)^{-1} \mathbf{Z}^{\prime} \mathbf{y}$ ? [Answer: No.]
  \item (Data matrix representation of $\widehat{\mathbf{S}}$ without conditional homoskedasticity) Define $\widehat{\Omega}$ suitably so that (6.7.5) equals the Bartlett kernel-based estimator of $\widehat{\mathbf{S}}$ without conditional homoskedasticity, namely, (6.6.5) with (6.6.7). Hint: Write $\widehat{\boldsymbol{\Omega}}$ as (6.7.6). If the lag length $q$ is known, $\hat{\omega}_{j}=\hat{\varepsilon}_{j} \hat{\varepsilon}_{t-j}$ for $0 \leq j \leq q$ and 0 for $j>q$. How would you define $\hat{\omega}_{j}$ for the Bartlett kernel-based estimator?
  \item Consider the simple regression model: $y_{t}=\beta z_{t}+\varepsilon_{t}$, where $\left\{\varepsilon_{t}\right\}$ is $\operatorname{AR}(1), \varepsilon_{t}= \phi \varepsilon_{t-1}+\eta_{t}$ with known $\phi$. Assume $\mathrm{E}\left(z_{t} \varepsilon_{t}\right)=0, \mathrm{E}\left(z_{t} \eta_{t}\right)=0, \mathrm{E}\left(z_{t-1} \eta_{t}\right)=0$, and conditional homoskedasticity. Also, to simplify, assume $\mathrm{E}\left(z_{t}\right)=0$. Show that
\end{enumerate}

$$
\operatorname{Avar}\left(\widehat{\beta}_{\mathrm{OLS}}\right)=\frac{\sigma_{\eta}^{2}}{\sigma_{z}^{2}} \frac{1+2 \sum_{j=1}^{\infty} \phi^{j} \rho_{j}}{1-\phi^{2}}, \quad \operatorname{Avar}\left(\widehat{\beta}_{\mathrm{GLS}}\right)=\frac{\sigma_{\eta}^{2}}{\sigma_{z}^{2}} \frac{1}{\left(1+\phi^{2}\right)-2 \phi \rho_{1}}
$$

where $\sigma_{z}^{2} \equiv \operatorname{Var}\left(z_{t}\right), \sigma_{\eta}^{2} \equiv \operatorname{Var}\left(\eta_{t}\right), \rho_{j} \equiv \operatorname{Corr}\left(z_{t}, z_{t-j}\right)$. Find a configuration of $\left(\phi,\left\{\gamma_{j}\right)\right)$ such that $\operatorname{Avar}\left(\widehat{\beta}_{\mathrm{OLS}}\right)>\operatorname{Avar}\left(\widehat{\beta}_{\mathrm{GLS}}\right)$. Explain why this is consistent with the fact that $\beta_{\mathrm{OLS}}$ is an efficient GMM estimator.

\subsection*{6.8 Application: Forward Exchange Rates as Optimal Predictors}
In foreign exchange markets, the percentage difference between the forward exchange rate and the spot exchange rate is called the forward premium. The relation between the forward premium and the expected rate of currency appreciation/depreciation over the life of the forward contract has been the subject of a large number of empirical studies. The simplest hypothesis about the relationship is that the two are the same. In this section, we test this hypothesis using the same weekly data set used by Bekaert and Hodrick (1993).

The data cover the period of 1975-1989 and include the following variables:\\
$S_{t}=$ spot exchange rate, stated in units of the foreign currency per dollar (e.g., 125 yen per dollar) on the Friday of week $t$,\\
$F_{t}=30$-day forward exchange rate on Friday of week $t$, namely, the price in foreign currency units of the dollar deliverable 30 days from Friday of week $t$,\\
$S 30_{t}=$ the spot rate on the delivery date on a 30-day forward contract made on Friday of week $t$.

The important feature of the data is that the maturity of the contract covers several sampling intervals; the delivery date is after the Friday of week $t+4$ but before the Friday of week $t+5$.

\section*{The Market Efficiency Hypothesis}
To indicate the assumptions underlying the hypothesis, consider three strategies for investing dollar funds for 30 days at date $t$. The first is to invest in a domestic money market instrument (e.g., U.S. Treasury bills) whose maturity is 30 days. Let $i_{t}$ be the rate of return over the 30 -day period. The second investment strategy is to convert dollars into foreign currency units, buy a 30 -day money market instrument denominated in the foreign currency whose interest rate is denoted $i_{t}^{*}$, and convert back to dollars at maturity. The rate of return from this investment strategy is

$$
S_{t} \cdot\left(1+i_{t}^{*}\right) / S 30_{t}
$$

This strategy involves foreign exchange risk because at time $t$ of investing you do not know $S 30_{t}$ which is the spot rate 30 days hence. The third investment strategy differs from the second in that it hedges against exchange risk by buying dollars forward at price $F_{t}$ at the same time the dollar fund is converted into foreign currency units. This hedged foreign investment provides a sure return equal to

$$
S_{t} \cdot\left(1+i_{t}^{*}\right) / F_{t}
$$

Because the first and the third strategies involve no risk, arbitrage requires that the rates of return be the same:

$$
1+i_{t}=S_{t} \cdot\left(1+i_{t}^{*}\right) / F_{t} .
$$

Taking logs of both sides and using the approximation $\log (1+x) \approx x$, we obtain the covered interest parity equation:


\begin{equation*}
i_{t}=s_{t}-f_{t}+i_{t}^{*} \quad \text { or } \quad f_{t}-s_{t}=i_{t}^{*}-i_{t}, \tag{6.8.1}
\end{equation*}


where lower case letters for $S$ and $F$ are natural logs. That is, the forward premium, $f_{t}-s_{t}$, equals the interest rate differential. This relationship holds almost exactly in real data. In fact, people routinely calculate the forward premium as the interest rate differential.

The rate of return from the second investment strategy is uncertain, but if investors are risk-neutral, its expected return is equal to the sure rate of return. Furthermore, if their expectations are rational, the expected return is the conditional expectation based on the information currently available, so


\begin{equation*}
\mathrm{E}\left[S_{t} \cdot\left(1+i_{t}^{*}\right) / S 30_{t} \mid I_{t}\right]=S_{t} \cdot\left(1+i_{t}^{*}\right) / F_{t}, \tag{6.8.2}
\end{equation*}


where $\mathrm{E}\left(\cdot \mid I_{t}\right)$ is the expectations operator and $I_{t}$ is the information set as of date $t$. Since $S_{t}, F_{t}$, and $i_{t}^{*}$ are known at date $t$, this equality reduces to

$$
\mathrm{E}\left(F_{t} / S 30_{t} \mid I_{t}\right)=1
$$

Again, by $\log$ approximation, $F_{t} / S 30_{t} \approx 1+f_{t}-s 30_{t}$. Therefore,


\begin{equation*}
\mathrm{E}\left(s 30_{t} \mid I_{t}\right)=f_{t} \text { or } \mathrm{E}\left(\varepsilon_{t} \mid I_{t}\right)=0 \quad \text { where } \varepsilon_{t} \equiv s 30_{t}-f_{t} . \tag{6.8.3}
\end{equation*}


This is the market efficiency hypothesis. As is clear from the derivation, the hypothesis assumes risk-neutrality and rational expectations. The hypothesis is sometimes described as the forward rate being the optimal forecast of future spot rates, because the conditional expectation is the optimal forecast in that it minimizes the mean square error (see Proposition 2.7).

\section*{Testing Whether the Unconditional Mean Is Zero}
Taking the unconditional expectation of both sides of (6.8.3), we obtain the unconditional relationship:


\begin{equation*}
\mathrm{E}\left(s 30_{t}\right)=\mathrm{E}\left(f_{t}\right) \quad \text { or } \quad \mathrm{E}\left(\varepsilon_{t}\right)=0 \tag{6.8.4}
\end{equation*}


That is, the forecast must be correct on average. Figure 6.1 is a plot of the forecast error $\varepsilon_{t}$ for the yen/dollar exchange rate. It looks like a stationary series, and we proceed under the assumption that it is. ${ }^{15}$ Because the forecast error is observable, we can test the hypothesis that its unconditional mean is zero. This is in contrast to the Fama exercise, where the forecast error (about future inflation) is observable only up to a constant.

To test the hypothesis of market efficiency, we need to derive the asymptotic distribution of the sample mean, $\bar{\varepsilon}$, under the hypothesis. The task is somewhat involved because $\left\{\varepsilon_{t}\right\}$ is serially correlated. To see why the forecast error has serial correlation, note that $I_{t}$, while including $\left\{\varepsilon_{t-5}, \varepsilon_{t-6}, \ldots\right\}$, does not include $\left\{\varepsilon_{t}, \ldots, \varepsilon_{t-4}\right\}$ because $\varepsilon_{t}$, which depends on $s 30_{t}$, cannot be observed until after $t+4$. Therefore,


\begin{equation*}
\mathrm{E}\left(\varepsilon_{t} \mid \varepsilon_{t-5}, \varepsilon_{t-6}, \ldots\right)=0 \tag{6.8.5}
\end{equation*}


It follows that $\operatorname{Cov}\left(\varepsilon_{t}, \varepsilon_{t-j}\right)=0$ for $j \geq 5$ but not necessarily for $j \leq 4$. In

\footnotetext{${ }^{15}$ Using covered interest rate parity, $\varepsilon_{t}=\left(s 30_{t}-s_{t}\right)-\left(i_{t}^{*}-i_{t}\right)$. If you believe that both the 30 -day rate of change $s 30_{t}-s_{t}$ and the interest rate differential $i_{t}^{*}-i_{t}$ are stationary, you have to believe that the forecast error is stationary.
}\begin{figure}[h]
\begin{center}
  \includegraphics[alt={},max width=\textwidth]{14fb0795-6d81-4e76-9788-792298f1ddae-438_803_1189_264_250}
\captionsetup{labelformat=empty}
\caption{Figure 6.1: Forecast Error, Yen/Dollar}
\end{center}
\end{figure}

the Fama exercise, the ex-post real interest rate (which is the inflation forecast error plus a constant) was serially uncorrelated because the maturity for the interest rate and the sampling interval were the same. Here, the forecast error has serial correlation because the maturity of the forward contract ( 30 days) is longer than the sampling interval (a week). The estimated correlogram is displayed along with the two-standard error band in Figure 6.2. ${ }^{16}$ The standard error is calculated under the assumption that the forecast error is i.i.d. So it equals $1 / \sqrt{n}$, as in Table 2.1. The pattern of autocorrelation is broadly consistent with the hypothesis: it stays well above two standard errors for the first four lags and generally lies within the band thereafter.

Under the null of market efficiency, since serial correlation vanishes after a finite lag, it is easy to see that $\left\{\varepsilon_{t}\right\}$ satisfies the essential part of Gordin's condition restricting serial correlation. By the Law of Iterated Expectations, (6.8.5) implies $\mathrm{E}\left(\varepsilon_{t} \mid \varepsilon_{t-j}, \varepsilon_{t-j-1}, \ldots\right)=0$ for $j \geq 5$, which immediately implies part (b) of Gordin's condition. Since the revision of expectations, $r_{t j}$, is zero for $j \geq 5$, part (c) is satisfied, provided that $r_{t j}(0 \leq j \leq 4)$ have finite second moments. It then

\footnotetext{${ }^{16}$ Although theory implies that the population mean is zero, I subtracted the sample mean in the calculation of sample autocorrelations.
}\begin{figure}[h]
\begin{center}
  \includegraphics[alt={},max width=\textwidth]{14fb0795-6d81-4e76-9788-792298f1ddae-439_716_1123_265_313}
\captionsetup{labelformat=empty}
\caption{Figure 6.2: Correlogram of $s 30-f$, Yen/Dollar}
\end{center}
\end{figure}

follows from Proposition 6.10 that


\begin{equation*}
\frac{\sqrt{n} \bar{\varepsilon}}{\sqrt{\gamma_{0}+2 \sum_{j=1}^{4} \gamma_{j}}}=\frac{\bar{\varepsilon}}{\sqrt{\left(\gamma_{0}+2 \sum_{j=1}^{4} \gamma_{j}\right) / n}} \mathrm{~d}^{\mathrm{d}} N(0,1) . \tag{6.8.6}
\end{equation*}


Furthermore, the denominator is consistently estimated by replacing $\gamma_{j}$ by its sample counterpart $\hat{\gamma}_{j}$. So the ratio


\begin{equation*}
\frac{\bar{\varepsilon}}{\sqrt{\left(\hat{\gamma}_{0}+2 \sum_{j=1}^{4} \hat{\gamma}_{j}\right) / n}} \tag{6.8.7}
\end{equation*}


is also asymptotically standard normal. The denominator of this expression will be called the standard error of the sample mean.

Table 6.1 reports results for testing the implication of the hypothesis that the unconditional mean of the forecast error is zero, for three foreign exchange rates against the dollar. The three currencies are the Deutsche mark, the British pound, and the Japanese yen. Columns 1 and 2 of the table report the mean and the standard deviation of the actual rate of change of the spot rate, $s 30_{t}-s_{t}$, and the forward premium, $f_{t}-s_{t}$. (Because the spot rate is per dollar, a positive rate of change means that the dollar strengthened.) The monthly rates of change and the

\begin{table}[h]
\begin{center}
\captionsetup{labelformat=empty}
\caption{Table 6.1: Means of Rates of Change, Forward Premiums, and Differences between the Two: 1975-1989}
\begin{tabular}{|l|l|l|l|}
\hline
\multirow{2}{*}{Exchange rate} & \multicolumn{3}{|c|}{Means and Standard Deviations} \\
\hline
 & $s 30-s$ (actual rate of change) & $f-s$ (expected rate of change) & Difference (unexpected rate of change) \\
\hline
¥/\$ & -4.98 (41.6) & -3.74 (3.65) & -1.25 (42.4) standard error $=3.56$ \\
\hline
DM/\$ & -1.78 (40.6) & -3.91 (2.17) & 2.13 (41.1) standard error $=3.26$ \\
\hline
f/\$ & 3.59 (39.2) & 2.16 (3.49) & 1.43 (39.9) standard error $=3.29$ \\
\hline
\end{tabular}
\end{center}
\end{table}

Note: $s 30-s$ is the rate of change over 30 days and $f-s$ is the forward premium, expressed as annualized percentage changes. Standard deviations are in parentheses. The standard errors are the values of the denominator of (6.8.7). The data are weekly, from 1975 to 1989. The sample size is 778.\\
forward premium have been multiplied by 1,200 to express the rates in annualized percentages. The market efficiency hypothesis is that the forward premium is the expected rate of change of the spot rate. It is evident that the actual rate of change is much more volatile than the forward premium. The third column reports the sample mean of $\varepsilon_{t}$, which equals the difference between the first two columns for the sample mean. The numbers in parentheses below the sample mean are the standard errors calculated as described above. In no case is there significant evidence against the null hypothesis of zero mean.

\section*{Regression Tests}
The unconditional test, however, does not exploit the implication of market efficiency that the forecast error is orthogonal to conditioning information $I_{t}$. Probably the most important variable in the information set $I_{t}$ for future spot rates is the forward rate $f_{t}$. We can test whether the forecast error is uncorrelated with the forward rate by regressing $\varepsilon_{t}$ on the constant and $f_{t}$. This is equivalent to running

\begin{figure}[h]
\begin{center}
  \includegraphics[alt={},max width=\textwidth]{14fb0795-6d81-4e76-9788-792298f1ddae-441_808_1157_278_309}
\captionsetup{labelformat=empty}
\caption{Figure 6.3: Yen/Dollar Spot Rate, Jan. 1975-Dec. 1989}
\end{center}
\end{figure}

the following regression


\begin{equation*}
s 30_{t}=\beta_{0}+\beta_{1} f_{t}+\varepsilon_{t}, \tag{6.8.8}
\end{equation*}


and testing whether $\beta_{0}=0$ and $\beta_{1}=1$. Because the error term is the forecast error orthogonal to anything known at date $t$ including $f_{t}$, the regressors are guaranteed to be orthogonal to the error term. So it might appear that we can estimate (6.8.8) by OLS and adjust for serial correlation in the error term as indicated in Section 6.6 for correct inference.

However, as the graph in Figure 6.3 suggests, the spot exchange rate looks like a process with increasing variance. As we will verify in Chapter 9, we cannot reject the hypothesis that the process has a unit root. The forward rate has the same property. (Its plot is not shown here because it is indistinguishable from the plot of the spot rate.) Thus, Assumption 3.2, which requires the regressor ( $f_{t}$ ) and the dependent variable ( $s 30_{t}$ ) to be stationary, is not satisfied. The problem actually runs deeper. We have observed in Figure 6.1 that the difference between $s 30_{t}$ and $f_{t}$ looked like a stationary series. In the language to be introduced in Chapter 10, $s 30_{t}$ and $f_{t}$ are "cointegrated" with a cointegrating vector of $(1,-1)$, in which case the OLS estimate of $\beta_{1}$ in (6.8.8) converges quickly to 1 . This can be verified in Figure 6.4 where $s 30_{t}$ is plotted against $f_{t}$. If a regression line is fitted, its slope

\begin{figure}[h]
\begin{center}
  \includegraphics[alt={},max width=\textwidth]{14fb0795-6d81-4e76-9788-792298f1ddae-442_756_1128_261_270}
\captionsetup{labelformat=empty}
\caption{Figure 6.4: Plot of $s 30$ against $f$, Yen/Dollar}
\end{center}
\end{figure}

is 0.99 . The problem from the point of view of testing market efficiency is that the OLS estimate converges to 1 regardless of whether or not the hypothesis is true. Therefore, the test has no power.

The problem can be dealt with by estimating a different regression,


\begin{equation*}
s 30_{t}-s_{t}=\beta_{0}+\beta_{1}\left(f_{t}-s_{t}\right)+\varepsilon_{t} . \tag{6.8.9}
\end{equation*}


Judging from the plot of $s 30_{t}-s_{t}$ and $f_{t}-s_{t}$ (not shown), it is reasonable to assume that they are stationary, which is consistent with Assumption 3.2 (ergodic stationarity). That the equation satisfies other GMM assumptions under market efficiency can be verified as follows. For $\beta_{0}=0$ and $\beta_{1}=1$, the error term $\varepsilon_{t}$ equals the forecast error $s 30_{t}-f_{t}$ and so the orthogonality conditions $\mathrm{E}\left(\varepsilon_{t}\right)=0$ and $\mathrm{E}\left[\left(f_{t}-s_{t}\right) \cdot \varepsilon_{t}\right]=0$ (Assumption 3.3 with $\mathbf{z}_{t}=\mathbf{x}_{t}$ ) are satisfied. The identification condition is that the second moment of $\left(1, f_{t}-s_{t}\right)$ be nonsingular, which is true if and only if $\operatorname{Var}\left(f_{t}-s_{t}\right)>0$. This is evidently satisfied; if the population variance were zero, we would not observe any variations in $f_{t}-s_{t}$. This leaves us with Assumption 3.5', which requires that there is not too much serial correlation in

$$
\mathbf{g}_{t} \equiv \mathbf{x}_{t} \cdot \varepsilon_{t}=\left[\begin{array}{c}
\varepsilon_{t} \\
\left(f_{t}-s_{t}\right) \varepsilon_{t}
\end{array}\right] .
$$

\begin{table}[h]
\begin{center}
\captionsetup{labelformat=empty}
\caption{Table 6.2: Regression Tests of Market Efficiency: 1975-1989}
\begin{tabular}{|l|l|l|l|l|}
\hline
\multicolumn{5}{|c|}{$s 30_{t}-s_{t}=\beta_{0}+\beta_{1}\left(f_{t}-s_{t}\right)+\varepsilon_{t}$} \\
\hline
\multirow{2}{*}{Currency} & \multicolumn{2}{|c|}{Regression Coefficients} & \multirow{2}{*}{$R^{2}$} & \multirow{2}{*}{Wald Statistic for $\mathrm{H}_{0}: \beta_{0}=0, \beta_{1}=1$} \\
\hline
 & Constant & Forward premium &  &  \\
\hline
¥/\$ & -12.8 (4.01) & -2.10 (0.738) & 0.034 & \( \begin{gathered} 18.6 \\ (p=0.009 \%) \end{gathered} \) \\
\hline
DM/\$ & -13.6 (5.72) & -3.01 (1.37) & 0.026 & \( \left(\begin{array}{c} 8.7 \\ (p=1.312 \%) \end{array}\right. \) \\
\hline
f/\$ & 7.96 (3.54) & -2.02 (0.85) & 0.033 & 12.9 ( $p=0.156 \%$ ) \\
\hline
\end{tabular}
\end{center}
\end{table}

NOTE: Standard errors are in parentheses. They are heteroskedasticity-robust and allow for serial correlation up to four lags calculated by the formula (6.8.11). The sample size is 778 . Our results differ slightly from those in Bekaert and Hodrick (1993) because they used the Bartlett kernel-based estimator with $q(n)=4$ to estimate S.

As already noted, since $I_{t}$ includes $\left\{\varepsilon_{t-5}, \varepsilon_{t-6}, \ldots\right\}$, the forecast error satisfies (6.8.5). Review Question 3 will ask you to show that $\left\{\mathbf{g}_{t}\right\}$ inherits the same property, namely,


\begin{equation*}
\mathrm{E}\left(\mathbf{g}_{t} \mid \mathbf{g}_{t-5}, \mathbf{g}_{t-6}, \ldots\right)=\mathbf{0} . \tag{6.8.10}
\end{equation*}


So, as in the unconditional test, Gordin's condition is satisfied, provided that the relevant second moments are finite. Since $\mathrm{E}\left(\mathbf{g}_{t} \mathbf{g}_{t-j}^{\prime}\right)=\mathbf{0}$ for $j \geq 5$, the long-run covariance matrix $\mathbf{S}$ can be estimated as in (6.6.4) with $q=4$. (In the empirical exercise, you will be asked to estimate $\widehat{\mathbf{S}}$ by the Bartlett kernel and also by VARHAC.)

To summarize, the efficient GMM estimator of $\boldsymbol{\beta}=\left(\beta_{0}, \beta_{1}\right)^{\prime}$ in (6.8.9) is OLS and the consistent estimator of its asymptotic variance is given by


\begin{equation*}
\widehat{\operatorname{Avar}\left(\widehat{\boldsymbol{\beta}}_{\mathrm{OLS}}\right)}=\mathbf{S}_{\mathbf{x x}}^{-1} \widehat{\mathbf{S}} \mathbf{S}_{\mathbf{x x}}^{-1}, \tag{6.8.11}
\end{equation*}


where $\widehat{\mathbf{S}}$ is given in (6.6.4) with $q=4$. Table 6.2 reports the regression test results for the three currencies. The corresponding plot for the yen/\$ exchange rate is in Figure 6.5. Notice that the estimated slope coefficient is not only significantly different from 1 but also negative, for all three currencies. The null hypothesis

\begin{figure}[h]
\begin{center}
  \includegraphics[alt={},max width=\textwidth]{14fb0795-6d81-4e76-9788-792298f1ddae-444_755_1180_271_246}
\captionsetup{labelformat=empty}
\caption{Figure 6.5: Plot of $s 30-s$ against $f-s$, Yen/Dollar}
\end{center}
\end{figure}

implied by market efficiency that $\beta_{0}=0$ and $\beta_{1}=1$ can be rejected decisively. This is reflected in the too many observations in the north-west quadrant of Figure 6.5 ; too often expected dollar depreciation (negative $f_{t}-s_{t}$ ) is associated with actual dollar appreciation (positive $s 30_{t}-s_{t}$ ).

\section*{QUESTIONS FOR REVIEW}
\begin{enumerate}
  \item (The standard error in the unconditional test) Let $\varepsilon_{t}=s 30_{t}-f_{t}$ and consider a regression of $\varepsilon_{t}$ on the constant. To account for serial correlation, use (6.6.4) with $q=4$ for $\widehat{\mathbf{S}}$. Verify that the $t$-value for the hypothesis that the intercept is zero is given by (6.8.7).
  \item (Lack of identification) For simplicity suppose $F_{t}$ is the two-week forward exchange rate so that the $s 30_{t}$ in (6.8.9) can be replaced by $s_{t+2}$. Suppose that the spot exchange rate is a random walk and $I_{t}=\left\{s_{t}, s_{t-1}, \ldots\right\}$. What should $f_{t}$ be under the hypothesis? Which of the required assumptions are violated? [Answer: The identification condition.] Suppose, instead, that $s_{t}$ is $\operatorname{AR}(1)$ satisfying $s_{t}=c+\phi s_{t-1}+\eta_{t}$ where $\left\{\eta_{t}\right\}$ is i.i.d. and $|\phi|<1$. Verify that $f_{t}=(1+\phi) c+\phi^{2} s_{t}$. Hint: $s_{t+2}=(1+\phi) c+\phi^{2} s_{t}+\left(\eta_{t+2}+\phi \eta_{t+1}\right)$. Verify that $\varepsilon_{t}=\eta_{t+2}+\phi \eta_{t+1}$.
  \item Prove (6.8.10). Hint: $\mathrm{E}\left(\varepsilon_{t} \mid I_{t}\right)=0$. Verify that $I_{t}$ includes $s_{t}-f_{t}$ and $\left\{\mathbf{g}_{t-5}, \mathbf{g}_{t-6}, \ldots\right\}$. Use the Law of Iterated Expectations.
\end{enumerate}

\section*{PROBLEM SET FOR CHAPTER 6}
\section*{ANALYTICAL EXERCISES}
\begin{enumerate}
  \item (Mean square limits and their uniqueness) In Chapter 2, we defined the mean square convergence of a sequence of random variables $\left\{z_{n}\right\}$ to a random variable $z$. In this and other questions, we will deal with sequences of random variables with finite second moments. For those sequences, the definition of mean square convergence is
\end{enumerate}

DEFINITION. Let $\left\{z_{n}\right\}$ be a sequence of random variables with $\mathrm{E}\left(z_{n}^{2}\right)<\infty$. We say that $\left\{z_{n}\right\}$ converges in mean square to a random variable $z$ with $\mathrm{E}\left(z^{2}\right)< \infty$, written $z_{n} \rightarrow_{\text {m.s. }} z$, if $\mathrm{E}\left[\left(z_{n}-z\right)^{2}\right] \rightarrow 0$ as $n \rightarrow \infty$.\\
(a) Show that, if $z_{n} \rightarrow_{\text {m.s. }} z$ and $z_{n} \rightarrow_{\text {m.s. }} z^{\prime}$, then $\mathrm{E}\left[\left(z-z^{\prime}\right)^{2}\right]=0$. Hint: $z-z^{\prime}=\left(z-z_{n}\right)+\left(z_{n}-z^{\prime}\right)$. Define $\|x\|=\sqrt{\mathrm{E}\left(x^{2}\right)}$. Then

$$
\begin{aligned}
& \left\|z-z^{\prime}\right\|^{2}=\left\|z-z_{n}\right\|^{2}+2 \mathrm{E}\left[\left(z-z_{n}\right)\left(z_{n}-z^{\prime}\right)\right]+\left\|z_{n}-z^{\prime}\right\|^{2} \\
& \leq\left\|z-z_{n}\right\|^{2}+2 \sqrt{\mathrm{E}\left[\left(z-z_{n}\right)^{2}\right]} \sqrt{\mathrm{E}\left[\left(z_{n}-z^{\prime}\right)^{2}\right]}+\left\|z_{n}-z^{\prime}\right\|^{2}
\end{aligned}
$$

(by Cauchy-Schwartz inequality that $\mathrm{E}(x y) \leq \sqrt{\mathrm{E}\left(x^{2}\right)} \sqrt{\left.\mathrm{E}\left(y^{2}\right)\right)}$

$$
=\left(\left\|z-z_{n}\right\|+\left\|z_{n}-z^{\prime}\right\|\right)^{2} .
$$

So $\left\|z-z^{\prime}\right\| \leq\left\|z-z_{n}\right\|+\left\|z_{n}-z^{\prime}\right\|$. This is called the triangle inequality.\\
(b) Show that $\operatorname{Prob}\left(z=z^{\prime}\right)=1$. Hint: Use Chebychev's inequality:

$$
\operatorname{Prob}\left(\left|z-z^{\prime}\right|>\varepsilon\right) \leq \frac{1}{\varepsilon^{2}} \mathrm{E}\left[\left\|z-z^{\prime}\right\|^{2}\right] .
$$

\begin{enumerate}
  \setcounter{enumi}{1}
  \item (Proof of Proposition 6.1) A well-known result about mean square convergence (see, e.g., Brockwell and Davis, 1991, Proposition 2.7.1) is\\
(i) $\left\{z_{n}\right\}$ converges in mean square if and only if $\mathrm{E}\left[\left(z_{m}-z_{n}\right)^{2}\right] \rightarrow 0$ as $m, n \rightarrow \infty$\\
(ii) If $x_{n} \rightarrow_{\text {m.s. }} x$ and $z_{n} \rightarrow_{\text {m.s. }} z$, then
\end{enumerate}

$$
\lim _{n \rightarrow \infty} \mathrm{E}\left(x_{n}\right)=\mathrm{E}(x) \quad \text { and } \quad \lim _{n \rightarrow \infty} \mathrm{E}\left(x_{n} z_{n}\right)=\mathrm{E}(x z) .
$$

Take this result for granted in answering.\\
(a) Prove that (6.1.6) converges in mean square as $n \rightarrow \infty$ under the hypothesis of Proposition 6.1. Hint: Let

$$
y_{t, n}=\mu+\sum_{j=0}^{n} \psi_{j} \varepsilon_{t-j}
$$

What needs to be proved is that $\left\{y_{t, n}\right\}$ converges in mean square as $n \rightarrow \infty$ for each $t$. Given (i) above, it is sufficient to prove that (assuming $m>n$ without loss of generality)

$$
\mathrm{E}\left[\left(\sum_{j=n+1}^{m} \psi_{j} \varepsilon_{t-j}\right)^{2}\right]=\sigma^{2} \sum_{j=n+1}^{m} \psi_{j}^{2} \rightarrow 0 \quad \text { as } m, n \rightarrow \infty .
$$

Use the fact that an absolutely summable sequence is square summable, i.e.,

$$
\sum_{j=0}^{\infty}\left|\psi_{j}\right|<\infty \Rightarrow \sum_{j=0}^{\infty} \psi_{j}^{2}<\infty
$$

and the fact that a sequence of real numbers $\left\{\alpha_{n}\right\}$ converges (to a finite limit) if and only if $\alpha_{n}$ is a Cauchy sequence:

$$
\alpha_{n} \rightarrow \alpha \Leftrightarrow\left|\alpha_{m}-\alpha_{n}\right| \rightarrow 0 \quad \text { as } m, n \rightarrow \infty .
$$

Set $\alpha_{n}=\sum_{j=0}^{n} \psi_{j}^{2}$.\\
(b) Prove that $\mathrm{E}\left(y_{t}\right)=\mu$. Hint: Let $y_{t, n}=\mu+\sum_{j=0}^{n} \psi_{j} \varepsilon_{t-j}$ as in (a). It was shown in (a) that $y_{t, n} \rightarrow_{\text {m.s. }} y_{t}$.\\
(c) (Proof of part (b) of Proposition 6.1) Show that

$$
\mathrm{E}\left[\left(y_{t}-\mu\right)\left(y_{t-j}-\mu\right)\right]=\lim _{n \rightarrow \infty} \mathrm{E}\left[\left(y_{t, n}-\mu\right)\left(y_{t-j, n}-\mu\right)\right] .
$$

Have we shown that $\left\{y_{t}\right\}$ is covariance-stationary? [Answer: Yes.]\\
(d) (Optional) Prove part (c) of Proposition 6.1, taking the following facts from calculus for granted.\\
(i) If $\left\{a_{j}\right\}$ is absolutely summable, then $\left\{a_{j}\right\}$ is summable (i.e., $-\infty< \left.\sum_{j=0}^{\infty} a_{j}<\infty\right)$ and

$$
\left|\sum_{j=0}^{\infty} a_{j}\right| \leq \sum_{j=0}^{\infty}\left|a_{j}\right|
$$

(ii) Consider a sequence with two subscripts, $\left\{a_{j k}\right\}(j, k=0,1,2, \ldots)$. Suppose $\sum_{j=0}^{\infty}\left|a_{j k}\right|<\infty$ for each $k$ and let $s_{k} \equiv \sum_{j=0}^{\infty}\left|a_{j k}\right|$. Suppose $\left\{s_{k}\right\}$ is summable. Then

$$
\left|\sum_{j=0}^{\infty}\left(\sum_{k=0}^{\infty} a_{j k}\right)\right|<\infty \quad \text { and } \quad \sum_{j=0}^{\infty}\left(\sum_{k=0}^{\infty} a_{j k}\right)=\sum_{k=0}^{\infty}\left(\sum_{j=0}^{\infty} a_{j k}\right)<\infty
$$

Hint: Derive

$$
\sum_{j=0}^{\infty}\left|\gamma_{j}\right| \leq \sigma^{2} \sum_{j=0}^{\infty}\left(\sum_{k=0}^{\infty}\left|\psi_{j+k}\right| \cdot\left|\psi_{k}\right|\right)<\infty
$$

\begin{enumerate}
  \setcounter{enumi}{2}
  \item (Autocovariances of $\left.h(L) x_{t}\right)$ We wish to derive the autocovariances of $\left\{y_{t}\right\}$ of Proposition 6.2. As in Proposition 6.2, assume throughout this question that $\left\{x_{t}\right\}$ is covariance-stationary. First consider a weighted average of finitely many terms:
\end{enumerate}

$$
y_{t, n}=h_{0} x_{t}+h_{1} x_{t-1}+\cdots+h_{n} x_{t-n}
$$

(a) Show that for this $\left\{y_{t, n}\right\}$,

$$
\gamma_{j}=\sum_{k=0}^{n} \sum_{\ell=0}^{n} h_{k} h_{\ell} \gamma_{j-k+\ell}^{x}
$$

where $\gamma_{j}^{x}$ is the $j$-th order autocovariance of $\left\{x_{t}\right\}$.\\
Verify that this reduces to (6.1.2) if $\left\{x_{t}\right\}$ is white noise. Hint: If you find this difficult, first set $n=1$, and then increase $n$ to establish the pattern.\\
(b) Now consider a weighted average of possibly infinitely many terms:

$$
y_{t}=h_{0} x_{t}+h_{1} x_{t-1}+\cdots
$$

Taking Proposition 6.2(a) for granted, show that

$$
\gamma_{j}=\sum_{k=0}^{\infty} \sum_{\ell=0}^{\infty} h_{k} h_{\ell} \gamma_{j-k+\ell}^{x} \quad \text { if }\left\{h_{j}\right\} \text { is absolutely summable. }
$$

Hint: $y_{t, n} \rightarrow_{\text {m.s. }} y_{t}$ as $n \rightarrow \infty$ by Proposition 6.2. Proceed as in 2(c).\\
4. (Homogeneous linear difference equations) Consider a $p$-th order homogeneous linear difference equation:


\begin{equation*}
y_{j}-\phi_{1} y_{j-1}-\phi_{2} y_{j-2}-\cdots-\phi_{p} y_{j-p}=0 \quad(j=p, p+1, \ldots), \tag{1}
\end{equation*}


where ( $\phi_{1}, \phi_{2}, \ldots, \phi_{p}$ ) are real constants with $\phi_{p} \neq 0$, with the initial condition that $\left(y_{0}, y_{1}, \ldots, y_{p-1}\right)$ are given. The associated polynomial equation is


\begin{equation*}
1-\phi_{1} z-\phi_{2} z^{2}-\cdots-\phi_{p-1} z^{p-1}-\phi_{p} z^{p}=0 . \tag{2}
\end{equation*}


Let $\lambda_{k}(k=1,2, \ldots, K)$ be the distinct roots of this equation (they can be complex numbers) and $r_{k}$ be the multiplicity of $\lambda_{k}$. So $r_{1}+r_{2}+\cdots+r_{K}=p$. The polynomial equation can be written as


\begin{equation*}
\prod_{k=1}^{K}\left(1-\lambda_{k}^{-1} z\right)^{r_{k}}=0 \tag{3}
\end{equation*}


For example, consider the second degree polynomial equation


\begin{equation*}
1-z+\frac{1}{4} z^{2}=0 \tag{4}
\end{equation*}


Its root is 2 with the multiplicity of 2 because

$$
1-z+\frac{1}{4} z^{2}=\left(1-\frac{1}{2} z\right)^{2}
$$

So $K=1$ and $r_{1}=2$ and $\lambda_{1}=2$ in this example.\\
As you verify below, the solution to the $p$-th order homogeneous equation can be written as


\begin{equation*}
y_{j}=\sum_{k=1}^{K} \sum_{n=0}^{r_{k}-1} c_{k n} \cdot(j)^{n} \lambda_{k}^{-j} \tag{5}
\end{equation*}


The coefficients $c_{k n}$ (there are $p$ of those because $r_{1}+\cdots+r_{K}=p$ ) can be determined from the initial conditions which specify the values for $\left(y_{0}, y_{1}, \ldots\right.$, $\left.y_{p-1}\right)$. In the above example of (4),


\begin{equation*}
y_{j}=c_{10} \cdot 2^{-j}+c_{11} \cdot j \cdot 2^{-j} \tag{6}
\end{equation*}


In the important special case where all the roots are distinct, $K=p$ and $r_{k}=1$\\
for all $k$, and (5) becomes


\begin{equation*}
y_{j}=\sum_{k=1}^{p} c_{k 0} \lambda_{k}^{-j} . \tag{7}
\end{equation*}


(a) For the special case of $p=2$ and distinct roots, (7) is


\begin{equation*}
y_{j}=c_{10} \lambda_{1}^{-j}+c_{20} \lambda_{2}^{-j} \tag{8}
\end{equation*}


Verify that this solves the homogeneous difference equation (1) with $p=2$ for any given ( $c_{10}, c_{20}$ ). To determine ( $c_{10}, c_{20}$ ), write down (8) for $j=0$ and $j=1$, and solve for $c$ 's as a function of ( $y_{0}, y_{1}, \lambda_{1}, \lambda_{2}$ ).\\
(b) To learn (remember) how to deal with multiple roots, consider example (4) above. Verify that (6) solves the second-order homogeneous difference equation

$$
y_{j}-y_{j-1}+\frac{1}{4} y_{j-2}=0
$$

(c) (Optional) Prove the following lemma.

Lemma. Let $0 \leq \xi<1$ and let $n$ be a non-negative integer. Then there exist two real numbers, $A$ and $b$, such that $\xi<b<1$ and $(j)^{n} \xi^{j}<A b^{j}$ for $j=0,1,2, \ldots$.

Hint: Pick any $b$ such that $\xi<b<1$. Because $b^{j}$ eventually gets larger than $(j)^{n} \xi^{j}$ as $j$ increases, there exists some $j$, call it $J$, such that $(j)^{n} \xi^{j}<b^{j}$ for all $j \geq J$.\\
(d) Suppose the stability condition holds for the $p$-th order homogeneous linear difference equation. Prove that there exist two numbers, $A$ and $b$, such that $0<b<1$ and

$$
\left|y_{j}\right|<A b^{j} \quad \text { for } j=0,1,2, \ldots .
$$

Hint:

$$
\left|y_{j}\right|=\left|\sum_{k=1}^{K} \sum_{n=0}^{r_{k}-1} c_{k n} \cdot(j)^{n} \lambda_{k}^{-j}\right|
$$

$$
\begin{aligned}
& \leq \sum_{k=1}^{K} \sum_{n=0}^{r_{k}-1}\left|c_{k n}\right| \cdot(j)^{n} \cdot\left|\lambda_{k}^{-1}\right|^{j} \\
& \leq c \sum_{k=1}^{K} \sum_{n=0}^{r_{k}-1}(j)^{n}\left|\lambda_{k}^{-1}\right|^{j}
\end{aligned}
$$

where $c=\max \left\{\left|c_{k n}\right|\right\}$. Since $\left|\lambda_{k}^{-1}\right|<1$ by the stability condition, we can apply the lemma with $\xi=\left|\lambda_{k}^{-1}\right|$ and claim

$$
(j)^{n}\left|\lambda_{k}^{-1}\right|^{j}<A_{k}\left(b_{k}\right)^{j} \quad \text { for all } j
$$

for some $A_{k}>0$ and $\left|\lambda_{k}^{-1}\right|<b_{k}<1$. Set $A^{*}=\max \left\{A_{k}\right\}$ and $b= \max \left\{b_{k}\right\}$, so that $A_{k}\left(b_{k}\right)^{j} \leq A^{*} b^{j}$ for all $k$. So

$$
\left|y_{j}\right|<c \sum_{k=1}^{K} \sum_{n=0}^{r_{k}-1} A^{*} b^{j}=c p A^{*} b^{j}
$$

(Recall that $p=r_{1}+\cdots+r_{K}$.) Then define $A$ to be $c p A^{*}$. You may find this sort of argument using inequalities hard to follow. If so, set $p, K$, and $\left(r_{1}, \ldots, r_{K}\right)$ to some specific numbers (e.g., $p=3, K=2, r_{1}=1$, $r_{2}=2$ ).\\
5. (Yule-Walker equations) In the text, the autocovariances of an $\operatorname{AR}(1)$ process were calculated from the MA( $\infty$ ) representation. Here, we use the YuleWalker equations.\\
(a) From (6.2.1') on page 376, derive

$$
\gamma_{j}-\phi \gamma_{j-1}=\mathrm{E}\left[\varepsilon_{t} \cdot\left(y_{t-j}-\mu\right)\right] \quad(j=0,1,2, \ldots) .
$$

(b) (Easy) From the MA representation, show that

$$
\mathrm{E}\left[\varepsilon_{t} \cdot\left(y_{t-j}-\mu\right)\right]= \begin{cases}\sigma^{2} & (j=0) \\ 0 & (j=1,2, \ldots)\end{cases}
$$

(c) Derive the Yule-Walker equations:


\begin{align*}
\gamma_{0}-\phi \gamma_{1} & =\sigma^{2}  \tag{9}\\
\gamma_{j}-\phi \gamma_{j-1} & =0 \quad(j=1,2, \ldots) \tag{10}
\end{align*}


(d) Calculate $\left(\gamma_{0}, \gamma_{1}, \ldots\right)$ from (9) and (10). Hint: First set $j=1$ in (10) to obtain $\gamma_{1}=\phi \gamma_{0}$. This and (9) can be solved for $\left(\gamma_{0}, \gamma_{1}\right)$.\\
6. (AR(1) without using filters) We wish to prove, without the help of Proposition 6.2, that $\left\{y_{t}\right\}$ following the $\operatorname{AR}(1)$ equation (6.2.1), with the stationarity condition $|\phi|<1$, must have an MA( $\infty$ ) representation if it is covariancestationary.\\
(a) Let $x_{t} \equiv y_{t}-\mu$ so that $\mathrm{E}\left(x_{t}\right)=0$. By successive substitution, derive from (6.2.1') (on page 376)

$$
x_{t}=x_{t, n}+\phi^{n} x_{t-n}
$$

where

$$
x_{t, n}=\varepsilon_{t}+\phi \varepsilon_{t-1}+\phi^{2} \varepsilon_{t-2}+\cdots+\phi^{n-1} \varepsilon_{t-n+1} .
$$

(b) Show that $x_{t, n} \rightarrow_{\text {m.s. }} x_{t}$. Hint: What needs to be shown is that $\mathrm{E}\left[\left(x_{t, n}-\right.\right. \left.\left.x_{t}\right)^{2}\right] \rightarrow 0$ as $n \rightarrow \infty$. $\left(x_{t, n}-x_{t}\right)^{2}=\phi^{2 n} x_{t-n}^{2}$. Because $\left\{y_{t}\right\}$ is assumed to be covariance-stationary, $\mathrm{E}\left(x_{t-n}^{2}\right)<\infty$.\\
(c) (Trivial) Recalling that the infinite sum

$$
\sum_{j=0}^{\infty} \phi^{j} \varepsilon_{t-j}
$$

is defined to be the mean square limit of the partial sum $x_{t, n}$, show that $y_{t}$ can be written as

$$
y_{t}=\mu+\sum_{j=0}^{\infty} \phi^{j} \varepsilon_{t-j} .
$$

\begin{enumerate}
  \setcounter{enumi}{6}
  \item (Companion form) We wish to generalize the argument just developed in question 6 to $\operatorname{AR}(p)$. Without loss of generality, we can assume that $\mu=0$ (or redefine $y_{t}$ to be the deviation from the mean of $y_{t}$ ). Define
\end{enumerate}

$$
\underset{(p \times 1)}{\boldsymbol{\xi}_{t}}=\left[\begin{array}{c}
y_{t} \\
y_{t-1} \\
y_{t-2} \\
\vdots \\
y_{t-p+1}
\end{array}\right], \quad \underset{(p \times 1)}{\mathbf{v}_{t}}=\left[\begin{array}{c}
\varepsilon_{t} \\
0 \\
0 \\
\vdots \\
0
\end{array}\right],
$$

$$
\underset{(p \times p)}{\mathbf{F}}=\left[\begin{array}{cccccc}
\phi_{1} & \phi_{2} & \phi_{3} & \cdots & \phi_{p-1} & \phi_{p} \\
1 & 0 & 0 & \cdots & 0 & 0 \\
0 & 1 & 0 & \cdots & 0 & 0 \\
\vdots & \vdots & \vdots & \cdots & \vdots & \vdots \\
0 & 0 & 0 & \cdots & 1 & 0
\end{array}\right] .
$$

Consider the following first-order vector stochastic difference equation:


\begin{equation*}
\xi_{t}=\mathbf{F} \xi_{t-1}+\mathbf{v}_{t} . \tag{11}
\end{equation*}


In this system of $p$ equations, called the companion form, the first equation is the same as the $\operatorname{AR}(p)$ equation (6.2.6') on page 379 and the rest are just identities about $y_{t-j}(j=1,2, \ldots, p-1)$.\\
(a) By successive substitution, show that


\begin{gather*}
\boldsymbol{\xi}_{t}=\boldsymbol{\xi}_{t, n}+(\mathbf{F})^{n} \boldsymbol{\xi}_{t-n}  \tag{12}\\
\boldsymbol{\xi}_{t, n} \equiv \mathbf{v}_{t}+\mathbf{F} \mathbf{v}_{t-1}+\mathbf{F}^{2} \mathbf{v}_{t-2}+\cdots+(\mathbf{F})^{n-1} \mathbf{v}_{t-n+1}
\end{gather*}


The eigenvalues or characteristic roots of a $p \times p$ matrix $\mathbf{F}$ are the roots of the determinantal equation


\begin{equation*}
\left|\mathbf{F}-\lambda \mathbf{I}_{p}\right|=0 . \tag{13}
\end{equation*}


It can be shown that the left-hand side of this is $\lambda^{p}-\phi_{1} \lambda^{p-1}-\cdots-\phi_{p-1} \lambda-\phi_{p}$. So the determinantal equation becomes the $p$-th order polynomial equation:


\begin{equation*}
\lambda^{p}-\phi_{1} \lambda^{p-1}-\cdots-\phi_{p-1} \lambda-\phi_{p}=0 . \tag{14}
\end{equation*}


(b) Verify this for $p=2$.

In the rest of this question, assume that all the roots of the $p$-th order polynomial equation (which, as just seen, are the eigenvalues of $\mathbf{F}$ ) are distinct. Let $\left(\lambda_{1}, \lambda_{2}, \ldots, \lambda_{p}\right)$ be those distinct eigenvalues. We have from matrix algebra that there exists a nonsingular $p \times p$ matrix $\mathbf{T}$ such that

\[
\mathbf{F}=\mathbf{T} \boldsymbol{\Lambda} \mathbf{T}^{-1}, \underset{(p \times p)}{\boldsymbol{\Lambda}}=\left[\begin{array}{ccccc}
\lambda_{1} & 0 & 0 & \cdots & 0  \tag{15}\\
0 & \lambda_{2} & 0 & \cdots & 0 \\
\vdots & & \ddots & & \vdots \\
0 & \cdots & 0 & \lambda_{p-1} & 0 \\
0 & \cdots & 0 & 0 & \lambda_{p}
\end{array}\right]
\]

(This fact does not depend on the special form of the matrix $\mathbf{F}$; all that is needed for (15) is that the $\lambda$ 's are distinct eigenvalues of $\mathbf{F}$.) From (15) it follows that

$$
(\mathbf{F})^{n}=\mathbf{T}(\mathbf{\Lambda})^{n} \mathbf{T}^{-1}
$$

(c) (Very easy) Prove this for $n=3$.\\
(d) Show that $\boldsymbol{\xi}_{t, n} \rightarrow_{\text {m.s. }} \boldsymbol{\xi}_{t}$ as $n \rightarrow \infty$ if the stationarity condition about the $\operatorname{AR}(p)$ equation is satisfied and if $\left\{y_{t}\right\}$ is covariance-stationary. Hint: The stationarity condition requires that all the roots of (13) be less than 1 in absolute value. By (12) what needs to be shown is that $(\mathbf{F})^{n} \boldsymbol{\xi}_{t-n} \rightarrow_{\text {m.s. }} \mathbf{0}$. Recall that $\mathbf{z}_{n} \rightarrow_{\text {m.s. }} \boldsymbol{\alpha}$ if the mean square convergence occurs component-wise.\\
(e) Show that $y_{t}$ has the MA( $\infty$ ) representation

$$
y_{t}=\psi_{0} \varepsilon_{t}+\psi_{1} \varepsilon_{t-1}+\psi_{2} \varepsilon_{t-2}+\cdots
$$

How is $\psi_{j}$ related to the characteristic roots? Hint: From (d) it immediately follows that

$$
\xi_{t}=\mathbf{v}_{t}+\mathbf{F} \mathbf{v}_{t-1}+(\mathbf{F})^{2} \mathbf{v}_{t-2}+\cdots
$$

The top equation of this is the MA( $\infty$ ) representation for $y_{t}$.\\
(If all the eigenvalues are not distinct, we need what is called the "Jordan decomposition" of $\mathbf{F}$ in place of (15).)\\
8. (AR(1) that started at $t=0$ ) In the text, the MA( $\infty$ ) representation (6.2.2) of an $\operatorname{AR}(1)$ process following $y_{t}=c+\phi y_{t-1}+\varepsilon_{t}$ assumes that the process is defined for $t=-1,-2, \ldots$ as well as for $t=0,1,2, \ldots$. Instead, consider an AR(1) process that started at $t=0$ with the initial value of $y_{0}$. Suppose that $|\phi|<1$ and that $y_{0}$ is uncorrelated with the subsequent white noise process ( $\varepsilon_{1}, \varepsilon_{2}, \ldots$ ).\\
(a) Write the variance of $y_{t}$ as a function of $\sigma^{2}\left(\equiv \operatorname{Var}\left(\varepsilon_{t}\right)\right)$, $\operatorname{Var}\left(y_{0}\right)$, and $\phi$. Hint: By successive substitution,

$$
\begin{aligned}
y_{t}= & \left(1+\phi+\phi^{2}+\cdots+\phi^{t-1}\right) c \\
& +\varepsilon_{t}+\phi \varepsilon_{t-1}+\phi^{2} \varepsilon_{t-2}+\cdots+\phi^{t-1} \varepsilon_{1}+\phi^{t} y_{0}
\end{aligned}
$$

Let $\mu$ and $\left\{\gamma_{j}\right\}$ be the mean and autocovariances of the covariancestationary $\operatorname{AR}(1)$ process, so $\mu=c /(1-\phi)$ and $\left\{\gamma_{j}\right\}$ are as in (6.2.5).

Show that

$$
\lim _{t \rightarrow \infty} \mathrm{E}\left(y_{t}\right)=\mu, \lim _{t \rightarrow \infty} \operatorname{Var}\left(y_{t}\right)=\gamma_{0}, \quad \text { and } \quad \lim _{t \rightarrow \infty} \operatorname{Cov}\left(y_{t}, y_{t-j}\right)=\gamma_{j} .
$$

(b) Show that $\left\{y_{t}\right\}$ is covariance-stationary if $\mathrm{E}\left(y_{0}\right)=\mu$ and $\operatorname{Var}\left(y_{0}\right)=\gamma_{0}$.\\
9. (Proof of Proposition 6.8(a)) We wish to prove Proposition 6.8(a). By Chebychev's inequality, it suffices to show that $\lim _{n \rightarrow \infty} \operatorname{Var}(\bar{y})=0$. From (6.5.2),

$$
\begin{aligned}
\operatorname{Var}(\bar{y}) & =\frac{\gamma_{0}}{n}+\frac{2}{n} \sum_{j=1}^{n-1}\left(1-\frac{j}{n}\right) \gamma_{j} \\
& =\frac{\gamma_{0}}{n}+\frac{2}{n} \sum_{j=1}^{n}\left(1-\frac{j}{n}\right) \gamma_{j} \quad\left(\text { since } 1-\frac{j}{n}=0 \text { for } j=n\right) \\
& \leq \frac{\gamma_{0}}{n}+\frac{2}{n} \sum_{j=1}^{n}\left(1-\frac{j}{n}\right)\left|\gamma_{j}\right| \\
& \leq \frac{\gamma_{0}}{n}+\frac{2}{n} \sum_{j=1}^{n}\left|\gamma_{j}\right| .
\end{aligned}
$$

(a) (Optional) Prove the following result:

$$
\lim _{j \rightarrow \infty} a_{j}=0 \Rightarrow \lim _{n \rightarrow \infty} \frac{1}{n} \sum_{j=1}^{n} a_{j}=0
$$

Hint: $\left\{a_{j}\right\}$ converges to 0 . So (i) $\left|a_{j}\right|<M$ for all $j$, and (ii) for any given $\varepsilon>$ 0 there exists a positive integer $N$ such that $\left|a_{j}\right|<\varepsilon / 2$ for all $j>N$. So

$$
\begin{aligned}
& \left|\frac{1}{n} \sum_{j=1}^{n} a_{j}\right| \leq \frac{1}{n} \sum_{j=1}^{n}\left|a_{j}\right| \\
& \quad=\frac{1}{n} \sum_{j=1}^{N}\left|a_{j}\right|+\frac{1}{n} \sum_{j=N+1}^{n}\left|a_{j}\right|<\frac{1}{n} \sum_{j=1}^{N} M+\frac{1}{n} \sum_{j=N+1}^{n} \frac{\varepsilon}{2} \\
& \quad=\frac{N M}{n}+\frac{(n-N)}{n} \frac{\varepsilon}{2}<\frac{N M}{n}+\frac{\varepsilon}{2} .
\end{aligned}
$$

By taking $n$ large enough, $N M / n$ can be made less than $\varepsilon / 2$.\\
(b) Taking the result in (a) as given, prove that $\operatorname{Var}(\bar{y}) \rightarrow 0$ as $n \rightarrow \infty$.\\
10. (Proof of Proposition 6.8(b))\\
(a) (optional) Prove the following result:

$$
\sum_{j=1}^{\infty} a_{j}<\infty \Rightarrow \lim _{n \rightarrow \infty} \frac{1}{n} \sum_{j=1}^{n} j a_{j}=0 .
$$

Hint:

$$
\begin{aligned}
\sum_{j=1}^{n} j a_{j}= & a_{1}+2 a_{2}+3 a_{3}+\cdots+n a_{n} \\
= & \left(a_{1}+a_{2}+\cdots+a_{n}\right)+\left(a_{2}+a_{3}+\cdots+a_{n}\right)+\cdots \\
& +\left(a_{n-1}+a_{n}\right)+a_{n} \\
= & \sum_{j=1}^{n} \sum_{k=j}^{n} a_{k} \leq \sum_{j=1}^{n}\left|\sum_{k=j}^{n} a_{k}\right| \\
= & \sum_{j=1}^{N}\left|\sum_{k=j}^{n} a_{k}\right|+\sum_{j=N+1}^{n}\left|\sum_{k=j}^{n} a_{k}\right| .
\end{aligned}
$$

Since $\left\{a_{j}\right\}$ is summable, (i) $\left|\sum_{k=j}^{n} a_{k}\right|<M$ for all $j, n$, and (ii) there exists a positive integer $N$ such that $\left|\sum_{k=j}^{n} a_{k}\right|<\varepsilon / 2$ for any given $\varepsilon>0$ for all $j, n>N$ (this follows because $\left\{\sum_{j=1}^{n} a_{j}\right\}$ is Cauchy).\\
(b) Taking this result as given, prove Proposition 6.8(b).

\section*{EMPIRICAL EXERCISES}
\begin{enumerate}
  \item Read at least pp. 115-119 (except the last two paragraphs of p. 119) of Bekaert and Hodrick (1993) before answering. Data files DM.ASC (for the Deutsche Mark), POUND.ASC (British Pound), and YEN.ASC (Japanese Yen) contain weekly data on the following items:
\end{enumerate}

Column 1: the date of the observation (e.g., "19850104" is January 4, 1985)\\
Column 2: the ask price of the dollar in units of the foreign currency in the spot market on Friday of the current week ( $S_{t}$ )\\
Column 3: the ask price of the dollar in units of the foreign currency in the 30-day forward market on Friday of the current week ( $F_{t}$ )\\
Column 4: the bid price of the dollar in units of the foreign currency in the spot market on the delivery date on a current forward contract $\left(S 30_{t}\right)$.

The sample period is the first week of 1975 through the last week of 1989. The sample size is 778. As in the text, define $s_{t} \equiv \log \left(S_{t}\right), f_{t} \equiv \log \left(F_{t}\right)$, $s 30_{t} \equiv \log \left(S 30_{t}\right)$. If $I_{t}$ is the information available on the Friday of week $t$, it includes $\left\{s_{t}, s_{t-1}, \ldots, f_{t}, f_{t-1}, \ldots, s 30_{t-5}, s 30_{t-6}, \ldots\right\}$. Note that $s 30_{t}$ is not observed until after the Friday of week $t+4$. Define $\varepsilon_{t} \equiv s 30_{t}-f_{t}$.

Pick your favorite currency to answer the following questions.\\
(a) (Library/internet work) For the foreign currency of your choice, identify the week when the absolute value of the forward premium is largest. For that week, find some measure of the domestic one-month interest rate (e.g., the one-month CD rate) for the United States and the currency's home country, to verify that the interest rate differential is as large as is indicated in the forward premium.\\
(b) (Correlogram of $\left\{\varepsilon_{t}\right\}$ ) Draw the sample correlogram of $\varepsilon_{t}$ with 40 lags. Does the autocorrelation appear to vanish after 4 lags? (It is up to you to decide whether to subtract the sample mean in the calculation of sample correlations. Theory says the population mean is zero, which you might want to impose in the calculation. In the calculation of the correlogram for the yen/\$ exchange rate shown in Figure 6.2, the sample mean was subtracted.)\\
(c) (Is the log spot rate a random walk?) Draw the sample correlogram of $s_{t+1}-s_{t}$ with 40 lags. For those 40 autocorrelations, use the Box-Ljung statistic to test for serial correlation. Can you reject the hypothesis that $\left\{s_{t}\right\}$ is a random walk with drift?\\
(d) (Unconditional test) Carry out the unconditional test. Can you replicate the results of Table 6.1 for the currency?\\
(e) (Optional, regression test with truncated kernel) Carry out the regression test. Can you replicate the results of Table 6.2 for the currency?

RATS Tip: Use LINREG with the ROBUSTERRORS option. To include autocorrelations up to the fourth, set LAGS $=4$ in LINREG. The DAMP $=$ 0 option instructs LINREG to use the truncated kernel.

TSP Tip: TSP's GMM procedure has the KERNEL option, but the choice is limited to Bartlett and the kernel called Parzen. You would have to use TSP's matrix commands to do the truncated kernel estimation.\\
(f) (Bartlett kernel) Use the Bartlett kernel-based estimator of $\mathbf{S}$ to do the regression test. Newey and West (1994) provide a data-dependent auto-

matic bandwidth selection procedure. Take for granted that the autocovariance lag length determined by this procedure is 12 for yen/ $\$$ (so autocovariances up to the twelfth are included in the calculation of $\widehat{\mathbf{S}}$ ), 8 for DM $/ \$$, and 16 for Pound $/ \$$.

RATS Tip: The ROBUSTERRORS, DAMP $=1$ option instructs LINREG to use Bartlett. For yen/\$, for example, set LAGS $=12$.

TSP Tip: Set het, kernel = bartlett in the GMM procedure. Set NMA = 12 for yen/\$.

The standard error for the $f-s$ coefficient for yen $/ \$$ should be 0.6815 .

(g) (Optional, VARHAC) Use the VARHAC estimator for $\widehat{\mathbf{S}}$ in the regression test. Set $p_{\text {max }}$ to the integer part of $n^{1 / 3}$. In Step 1 of the VARHAC, you would estimate a bivariate VAR, and for each of the two VAR equations, you would pick the lag length $p$ by BIC on the fixed sample of $t=p_{\text {max }}+ 1, \ldots, T$. The information criterion could be (6.6.11) or, alternatively,


\begin{equation*}
\log \left(S S R_{p}^{(k)} /\left(n-p_{\max }\right)\right)+p K \cdot \log \left(n-p_{\max }\right) /\left(n-p_{\max }\right) \tag{16}
\end{equation*}


Just to be clear about how you picked $p$ in your answer, use this latter objective function in selecting $p$ by BIC. For the case of yen/ $\$$, the lag length selected should be 4 for the first VAR equation and 6 for the second equation. For yen $/ \$$, the standard error for the $f-s$ coefficient should be 0.8027.

\begin{enumerate}
  \setcounter{enumi}{1}
  \item (Continuation of the Fama exercise of Chapter 2) Now we know how to use the three-month T-bill rate to test the efficiency of the U.S. Treasury bills market using monthly data. Let
\end{enumerate}

$$
\begin{aligned}
\pi_{t}= & \text { inflation rate (in percent, annual rate) over the three-month } \\
& \text { period ending in month } t-1, \\
R_{t}= & \text { three-month } \mathrm{T} \text {-bill rate (in percent, annual rate) at the begin- } \\
& \text { ning of month } t .
\end{aligned}
$$

Consider the Fama regression for testing market efficiency:

$$
\pi_{t+3}=\beta_{0}+\beta_{1} R_{t}+\varepsilon_{t}
$$

(a) (Trivial) Why is the dependent variable $\pi_{t+3}$ rather than, say, $\pi_{t}$ ?\\
(b) Show that, under the efficient market hypothesis, $\beta_{1}=1$ and the autocovariance of $\left\{\varepsilon_{t}\right\}$ vanishes after the second lag (i.e., $\gamma_{j}=0$ for $j>2$ ).\\
(c) Let $r_{t+3} \equiv R_{t}-\pi_{t+3}$ be the ex-post real rate on the three-month T-bill. Use PAI3 in MISHKIN.ASC as the measure of the three-month inflation rate to calculate the ex-post real rate. (The timing of the variables in MISHKIN.ASC is such that a January interest rate observation uses end-ofDecember bill rate data and a January observation for a three-month inflation rate is from the December to March CPI data. So PAI3 and TB3 (threemonth T-bill rate) for the same month can be paired in the regression.) Draw the sample correlogram for the sample period of $1 / 53-7 / 71$. What is the prediction of the efficient market hypothesis about the correlogram?\\
(d) Test market efficiency by estimating the Fama regression for the sample period of $1 / 53-7 / 71$. (The sample size should be 223; do not throw away two-thirds of the observations, as Fama did in his Tables 6-8.) Use (6.6.4) with $q=2$ for $\widehat{\mathbf{S}}$.

\section*{ANSWERS TO SELECTED QUESTIONS}
$\_\_\_\_$

\section*{ANALYTICAL EXERCISES}
2d. Since $\left\{\psi_{j}\right\}$ is absolutely summable, $\psi_{j} \rightarrow 0$ as $j \rightarrow \infty$. So for any $j$, there exists an $A>0$ such that $\left|\psi_{j+k}\right| \leq A$ for all $j, k$. So $\left|\psi_{j+k} \cdot \psi_{k}\right| \leq A\left|\psi_{k}\right|$. Since $\left\{\psi_{k}\right\}$ (and hence $\left\{A \psi_{k}\right\}$ ) is absolutely summable, so is $\left\{\psi_{j+k} \cdot \psi_{k}\right\}$ ( $k= 0,1,2, \ldots$ ) for any given $j$. Thus by (i),

$$
\left|\gamma_{j}\right|=\sigma^{2}\left|\sum_{k=0}^{\infty} \psi_{j+k} \psi_{k}\right| \leq \sigma^{2} \sum_{k=0}^{\infty}\left|\psi_{j+k} \psi_{k}\right|=\sigma^{2} \sum_{k=0}^{\infty}\left|\psi_{j+k}\right|\left|\psi_{k}\right|<\infty .
$$

Now set $a_{j k}$ in (ii) to $\left|\psi_{j+k}\right| \cdot\left|\psi_{k}\right|$. Then

$$
\sum_{j=0}^{\infty}\left|a_{j k}\right|=\sum_{j=0}^{\infty}\left|\psi_{k}\right|\left|\psi_{j+k}\right| \leq\left|\psi_{k}\right| \sum_{j=0}^{\infty}\left|\psi_{j}\right|<\infty .
$$

Let

$$
M \equiv \sum_{j=0}^{\infty}\left|\psi_{j}\right| \quad \text { and } \quad s_{k} \equiv\left|\psi_{k}\right| \sum_{j=0}^{\infty}\left|\psi_{j+k}\right| .
$$

Then $\left\{s_{k}\right\}$ is summable because $\left|s_{k}\right| \leq\left|\psi_{k}\right| \cdot M$ and $\left\{\psi_{k}\right\}$ is absolutely summable. Therefore, by (ii),

$$
\sum_{j=0}^{\infty}\left(\sum_{k=0}^{\infty}\left|\psi_{j+k}\right| \cdot\left|\psi_{k}\right|\right)<\infty .
$$

This times $\sigma^{2}$ is $\sum_{j=0}^{\infty}\left|\gamma_{j}\right|$. So $\left\{\gamma_{j}\right\}$ is absolutely summable.\\
7e. $\psi_{j}$ is the $(1,1)$ element of $\mathbf{T}(\mathbf{\Lambda})^{j} \mathbf{T}^{-1}$.\\
8a.

$$
\begin{aligned}
& \mathrm{E}\left(y_{t}\right)=\frac{1-\phi^{t}}{1-\phi} c+\phi^{t} \mathrm{E}\left(y_{0}\right) \\
& =\left(1-\phi^{t}\right) \mu+\phi^{t} \mathrm{E}\left(y_{0}\right) \quad(\text { since } \mu=c /(1-\phi)) \\
& =\mu+\phi^{t} \cdot\left[\mathrm{E}\left(y_{0}\right)-\mu\right], \\
& \begin{aligned}
\operatorname{Var}\left(y_{t}\right)= & \left(1+\phi^{2}+\phi^{4}+\cdots+\phi^{2 t-2}\right) \sigma^{2} \\
& +\phi^{2 t} \operatorname{Var}\left(y_{0}\right) \quad\left(\text { since } \operatorname{Cov}\left(\varepsilon_{t}, y_{0}\right)=0\right) \\
= & \frac{1-\phi^{2 t}}{1-\phi^{2}} \sigma^{2}+\phi^{2 t} \operatorname{Var}\left(y_{0}\right) \\
= & \left(1-\phi^{2 t}\right) \gamma_{0}+\phi^{2 t} \operatorname{Var}\left(y_{0}\right) \quad\left(\text { since } \gamma_{0} \equiv \sigma^{2} /\left(1-\phi^{2}\right)\right), \\
= & \gamma_{0}+ \\
& \phi^{2 t} \cdot\left[\operatorname{Var}\left(y_{0}\right)-\gamma_{0}\right], \\
\operatorname{Cov}\left(y_{t}, y_{t-j}\right)= & \left(\phi^{j}+\phi^{j+2}+\cdots+\phi^{2 t-j-2}\right) \sigma^{2}+\phi^{2 t-j} \operatorname{Var}\left(y_{0}\right) \\
= & \phi^{j} \cdot\left(1+\phi^{2}+\cdots+\phi^{2 t-2 j-2}\right) \sigma^{2}+\phi^{2 t-j} \operatorname{Var}\left(y_{0}\right) \\
= & \phi^{j} \cdot \frac{1-\phi^{2(t-j)}}{1-\phi^{2}} \sigma^{2}+\phi^{2 t-j} \operatorname{Var}\left(y_{0}\right) \\
= & \phi^{j} \cdot\left(1-\phi^{2(t-j)}\right) \gamma_{0}+\phi^{2 t-j} \operatorname{Var}\left(y_{0}\right) \\
= & \gamma_{0} \phi^{j}+\phi^{2 t-j} \cdot\left[\operatorname{Var}\left(y_{0}\right)-\gamma_{0}\right] \\
= & \gamma_{j}+\phi^{2 t-j} \cdot\left[\operatorname{Var}\left(y_{0}\right)-\gamma_{0}\right] .
\end{aligned}
\end{aligned}
$$

\section*{EMPIRICAL EXERCISES}
1c. The Ljung-Box statistic with 40 lags is 84.0 with a $p$-value of $0.0058 \%$. So the hypothesis that the spot rate is a random walk with drift can be rejected. See Figure 6.6 for the correlogram for yen/\$.

1f. For the yen/\$ exchange rate, the standard errors are 3.72 for the constant and 0.68 for the $f-s$ coefficient. The Wald statistic for the hypothesis that $\beta_{0}=0$ and $\beta_{1}=1$ is $21.85(p=0.002 \%)$.

\begin{figure}[h]
\begin{center}
  \includegraphics[alt={},max width=\textwidth]{14fb0795-6d81-4e76-9788-792298f1ddae-460_714_1133_267_284}
\captionsetup{labelformat=empty}
\caption{Figure 6.6: Correlogram of $s_{t+1}-s_{t}$, Yen/Dollar}
\end{center}
\end{figure}

\begin{figure}[h]
\begin{center}
  \includegraphics[alt={},max width=\textwidth]{14fb0795-6d81-4e76-9788-792298f1ddae-460_720_1139_1162_282}
\captionsetup{labelformat=empty}
\caption{Figure 6.7: Correlogram of Three-Month Ex-Post Real Rate}
\end{center}
\end{figure}

2c. The correlogram for the three-month ex-post real rate is shown above (Figure 6.7 ). Theory predicts that serial correlation vanishes after the second lag. The correlogram is mostly consistent with this prediction.

\section*{References}
Anderson, T. W., 1971, The Statistical Analysis of Time Series, New York: Wiley.\\
Andrews, D., 1991, "Heteroskedasticity and Autocorrelation Consistent Covariance Matrix Estimation," Econometrica, 59, 817-858.\\
Bekaert, G., and R. Hodrick, 1993, "On Biases in the Measurement of Foreign Exchange Risk Premiums," Journal of International Money and Finance, 12, 115-138.\\
Brockwell, P., and R. Davis, 1991, Time Series: Theory and Methods (2d ed.), Berlin: SpringerVerlag.\\
Den Haan, W., and A. Levin, 1996a, "Inferences from Parametric and Non-Parametric Covariance Matrix Estimation Procedures," Technical Working Paper Series 195, NBER.\\
$\_\_\_\_$ , 1996b, "A Practitioner's Guide to Robust Covariance Matrix Estimation," Technical Working Paper Series 197, NBER.\\
Fuller, W., 1996, Introduction to Statistical Time Series (2d ed.), New York: Wiley.\\
Gordin, M. I., 1969, "The Central Limit Theorem for Stationary Processes," Soviet Math. Dokl., 10, 1174-1176.\\
Gourieroux, C., and A. Monfort, 1997, Time Series and Dynamic Models, Cambridge: Cambridge University Press.\\
Hamilton, J., 1994, Time Series Analysis, Princeton: Princeton University Press.\\
Hannan, E. J., 1970, Multiple Time Series, New York: Wiley.\\
Hannan, E. J., and M. Deistler, 1988, The Statistical Theory of Linear Systems, New York: Wiley.\\
Hansen, L., 1982, "Large Sample Properties of Generalized Method of Moments Estimators," Econometrica, 50, 1029-1054.\\
Lütkepohl, H., 1993, Introduction to Multiple Time Series Analysis (2d ed.), New York: SpringerVerlag.\\
Newey, W., and K. West, 1987, "Hypothesis Testing with Efficient Method of Moments Estimation," International Economic Review, 28, 777-787.\\
$\_\_\_\_$ , 1994, "Automatic Lag Selection in Covariance Matrix Estimation," Review of Economic Studies, 61, 631-653.\\
Ng, S., and P. Perron, 2000, "A Note on Model Selection in Time Series Analysis," mimeo., Boston: Boston College.\\
White, H., 1984, Asymptotic Theory for Econometricians, New York: Academic.


\end{document}
